% ---------- 01-intro.tex
\chapter{Introduction}\label{sec:intro}

\section{Background}\label{sec:intro-background}

Missing-data means, average treatment effects, and expected conditional
covariances under random design are averages over the covariate
distribution of functions of regressions given the covariates. Their
estimation rates depend on what is known about both the regressions and
the covariate distribution. In missing-response problems, Robins and
Ritov~\cite{robinsritov} showed that attainable rates depend on what is
known about the nuisance functions. We study \emph{rough
design}, in which the covariate density is unknown and subject only to
fixed upper and lower bounds, while the regressions are assumed smooth.
We show that below the root-$n$ threshold, rough design reduces the
minimax polynomial exponent from $4s/(4s+d)$ for a known or
sufficiently smooth density to $2s/d$, where $s$ is the average
smoothness of the two nuisance functions and $d$ is the covariate
dimension. The threshold $s=d/4$ itself is unchanged.

For a missing-data mean, let $\alpha$ and $\beta$ be the H\"older
smoothness indices of the inverse observation probability and of the
outcome regression. For the average product of two conditional means,
which underlies the expected conditional covariance, let them be the
indices of the two means. Let $s=(\alpha+\beta)/2$. Estimators of
these parameters impose assumptions on a density weight $g$, which is
the covariate density or its product with the observation probability,
depending on the parameter. When $g$ is known or belongs to a H\"older
class of smoothness $\gamma>0$ such that
$\gamma/(2\gamma+d)>\{\max(\alpha,\beta)/d\}(d-4s)/(d+4s)$,
higher-order influence function (HOIF) estimators attain
root-mean-square error of order
\begin{equation}\label{eq:known-weight-rate}
 n^{-1/2}\vee n^{-4s/(4s+d)},
\end{equation}
where $x\vee y=\max(x,y)$. This condition and the upper bound are
from~\cite[Section~3.1]{robins}, the corrected version
of~\cite{robins2017}, which treats the missing-data mean there and
indicates the extension to other doubly robust functionals
in~\cite[Section~11.3]{robins}. Below the threshold, the rate
\eqref{eq:known-weight-rate} cannot be improved even when $g$ is known,
as shown in~\cite{robins2009}; see~\cite[Section~3.1]{robins}. Rates of
this form are classical for quadratic functionals of a
density~\cite{bickelritov,birgemassart}. For the average product of two
conditional means, McGrath and
Mukherjee~\cite[Corollaries~1 and~16]{mcgrath} obtain this rate for a
uniform or sufficiently smooth known density, and below $s=d/4$ their
results also cover an unknown density that is sufficiently smooth.

The rough-design question arises in the HOIF
literature~\cite{hoif2008,robins2009,li2011,robins}. Immediately after
their equation~(1.3), Robins et al.~\cite[p.~338]{hoif2008} conjecture
that the rate given there is minimax, up to logarithmic factors, when
the nuisance functions are equally smooth and the density
smoothness $\gamma$ is too small for their condition~(1.2), which for
equal indices is the condition stated before
\eqref{eq:known-weight-rate}. They note that, as $\gamma\to0$, this rate
tends to $n^{-2s/d}\log n$, which agrees up to the logarithmic factor
with the minimax rate $n^{-2s/d}$ for a constant conditional variance
under an equally spaced fixed design~\cite{wangvariance,caivariance}.
This limit suggests that the exponent under rough design is $2s/d$.

Robins et al.~\cite[Section~3.1]{robins} also state that heuristic
arguments indicate a rate slower than \eqref{eq:known-weight-rate} for
smaller $\gamma$. Kennedy, Balakrishnan, and
Wasserman~\cite[Section~4]{kennedy2020discussion} describe the minimax
problem with an unknown density as open. They suggest that the rate
with an unknown covariate density might interpolate between the
known-density rate and the fixed-design rate $n^{-2s/d}$, which would
suggest agreement between rough random design and fixed design at the
nonsmooth extreme. For a counterfactual mean, Liu, Mukherjee, and
Robins~\cite[Section~4]{liumukherjeerobins2020} describe an empirical
HOIF estimator, based on~\cite{liu}, that attains $n^{-2s/d}$ up to
logarithmic factors without complexity assumptions on the covariate
density. McClean et al.~\cite[Section~4.2 and
Theorem~1]{mcclean} construct an estimator of the expected conditional
covariance with a convergence rate in probability of $n^{-2s/d}$, up to
a polylogarithmic factor, when $s<d/4$, and they leave open whether
this rate is minimax under rough design.

\section{Main results}\label{sec:intro-results}

Our main result, \cref{thm:generic}, gives minimax bounds in a
nonparametric model in which the weight $g$ is unknown and satisfies
only fixed bounds $0<g_-\le g\le g_+<\infty$ with $g_-<g_+$. It covers
a class of doubly robust functionals~\cite[Section~3.1]{hoif2008} that
includes missing-data means and expected conditional covariances
(\cref{sec:functionals}). Under interior and nondegeneracy
conditions stated in \cref{sec:exactclass}, the minimax
root-mean-square error $\risk$ satisfies
\begin{equation}\label{eq:intro-refinement}
 \log\risk=
 \begin{cases}
 -\dfrac{2s}{d}\log n
 -\dfrac{2}{d}\sqrt{2s(d-4s)\log n\,
   \log\dfrac{\sqrt{g_+}+\sqrt{g_-}}{\sqrt{g_+}-\sqrt{g_-}}}
       +O(\log\log n), & s<d/4,\\
 -\tfrac12\log n+O(1), & s\ge d/4.
 \end{cases}
\end{equation}
Below $s=d/4$, \eqref{eq:intro-refinement} identifies the
leading correction to the polynomial rate, a factor
$e^{-\kappa\sqrt{\log n}}$ with the explicit constant $\kappa$ of
\eqref{eq:subcritical-scale}. In this regime the
upper and lower bounds differ by a polylogarithmic factor, which is
$(\log n)^{1/4}$ when both smoothness indices are at most one
(\cref{rem:gap}). The lower bounds hold even on the subclass in
which the functions $w,a,b$ of \eqref{eq:generic-moments} stay within
an arbitrary fixed distance of constant baseline values
(\cref{sec:exactclass}). In particular, they apply directly to
models with additional overlap or range restrictions. They also persist if the
density must be $C^\infty$, as long as the model imposes no fixed bound
on its derivatives (\cref{rem:smooth-design}).
\Cref{sec:applications} applies the theorem to the targets in
\cref{tab:applications}.

By \cref{thm:generic}, the exponent $2s/d$ is minimax below the
threshold, and our estimator attains it in root-mean-square error,
uniformly over fixed H\"older classes, for every functional in the class
of \cref{thm:generic}, whether or not the two smoothness indices are
equal. For equal indices, $2s/d$ is the exponent suggested
by the limit $\gamma\to0$ of the rate that Robins et al.\ conjecture to
be minimax. By \eqref{eq:intro-refinement}, however, the minimax risk is
smaller than $n^{-2s/d}$ by a factor $e^{-\kappa\sqrt{\log n}+O(\log\log n)}$,
which tends to zero faster than every power of $1/\log n$. The limit
$n^{-2s/d}\log n$ of their rate exceeds the minimax risk
under rough design by more than any polylogarithmic factor. For the
expected conditional covariance, Park~\cite{park} has obtained closely
related bounds independently and concurrently, with the same
threshold, polynomial exponent, and constant $\kappa$; we compare the
two works in \cref{sec:related}.

\section{Ideas of the proofs}\label{sec:ideas}

We pair moment-matching priors with polynomial approximation,
as in minimax estimation of nonsmooth
functionals~\cite{lepski,cailow,jiaovenkathanweissman,wuyangentropy,wuyang}.
For an interval $I=[x_-,x_+]$ with
$0<x_-<x_+$, let $c_I=(x_++x_-)/2$ and $d_I=(x_+-x_-)/2$. Then, for
real $t$,
\begin{equation}\label{eq:intro-cosine}
 \frac1{c_I+d_I\cos t}=\frac1{\sqrt{x_-x_+}}
 \Bigl\{1+2\sum_{k\ge1}(-\rho_I)^k\cos(kt)\Bigr\},
 \qquad
 \rho_I=\frac{\sqrt{x_+}-\sqrt{x_-}}{\sqrt{x_+}+\sqrt{x_-}}
\end{equation}
(\cref{lem:reciprocal}). Since $\cos(kt)$ is a polynomial of
degree $k$ in $c_I+d_I\cos t$, the partial sums of this series are
polynomials that approximate the reciprocal on $I$, with error
$O(\rho_I^m)$ at order $m$. For $I=[g_-,g_+]$, $\rho_I$ is the $\rho$
of \eqref{eq:sharp-parameters}.

\emph{Upper bound.} Recall that $\psi(P)=\E_PW+\lambda\int abg$, with
$a,b,g$ as in \eqref{eq:generic-moments}. We project $a$ and $b$ in
$L^2(g)$ onto piecewise
polynomials on nested dyadic partitions. With $K$ cells, the projection
error of $\int abg$ is the inner product of the two projection errors,
as for the truncation bias of HOIF
estimators~\cite[Theorem~3.8]{hoif2008}, and it is of order
$K^{-\theta}$, where $\theta=2s/d$. By reciprocal approximation, we replace each projection increment by a
polynomial in observable cell moments, with error
$O(K^{-\theta}(m+1)^\nu\rho^m)$ at truncation order $m$, where $\nu$
is defined in \eqref{eq:poly-dimension}
(\cref{lem:projection-polynomials}). We estimate these polynomials
from the $n$ observations without bias by their lift, which replaces
each monomial by a U-statistic~\cite{halmos,hoeffding};
\cref{lem:lift} expresses the covariance of the lifts exactly
through centered derivatives. We then choose the degrees across
resolutions by the variance bounds of \cref{lem:variance}, without
an estimate of $g$. Polynomial approximations of bilinear forms in the
inverse of a Gram matrix, with U-statistic estimates of their terms,
were used before in~\cite{kongvaliant2018,chenliumukherjee2025}
(\cref{sec:related}).

\emph{Lower bound.} We compare two priors by the method of two fuzzy
hypotheses, and we follow the strategy of the lower bound of Aronow and
Lopatto~\cite{aronowlopatto} for a constant conditional variance, with
a random covariate density, a Poissonized sample, and a matching order
$M$ of order $\sqrt{\log n}$ (\cref{sec:related}; see also \cref{sec:lower-ideas}).
Here the matching comes from a random phase. On the bump supports of the two blocks of a
pair, the density equals $r_0\pm h_n\cos\Theta$, with a uniform phase
$\Theta$. The conditional law of the response is perturbed by $u,v$
and is affine in $(u,v)$; the perturbations are bumps of heights
$A_u,A_v$ divided by the density, with signs correlated through
$\pm\cos(M\Theta)$ under the two priors, for an even integer $M$. In
the generic model, they shift $a$ and $b$ by $u/w$ and $v/w$, or both
by $(u+v)/w$ when the response functions $U$ and $V$ of
\cref{sec:functionals} coincide (\cref{sec:generic-embedding}). Each likelihood
coefficient that differs between the priors contains an odd number of
factors of each perturbation. The signs contribute $\cos(M\Theta)$,
and apart from their signs the products $pu$ and $pv$ do not depend on
$\Theta$. Conditional on the other phase parameters, $j$ additional
density factors form a trigonometric polynomial of degree at most $j$
in $\Theta$. We conclude that every such coefficient vanishes for
$j<M$, including those with higher powers of the perturbations
(\cref{lem:poisson-comparison}). The target means separate by a
multiple of $A_uA_v$ times the $M$th Fourier coefficient of the
reciprocal density, which is of order $\rho^M$ by
\eqref{eq:intro-cosine}. With block-constant phases, the case $J=0$ of
\cref{lem:lattice-priors}, we obtain $\risk\ge c\,a_n$
(\cref{prop:reduction}), where $a_n$ is defined in
\eqref{eq:subcritical-scale}. In the lattice refinement, we let
the phase vary across the $R$ fine cells of each block. For
$M\le j<M+\sqrt M$, a coefficient then survives only where the binary
digits of most of its points nearly agree, which gives the factor
$R^{1-M}$; a cruder bound suffices for larger $j$. We use a random gate to keep
the H\"older norms bounded while the bump heights gain a factor
$m\asymp M$ (\cref{lem:lattice-priors}), and the lower bound gains
the factor $m^2\asymp\log n$.

\emph{The common balance.} Let $0<\theta<1/2$, $L=\log n$, and
$\Delta=1-2\theta$, and let $\tau$, $\kappa$ and $\nu$ be as in
\eqref{eq:sharp-parameters}--\eqref{eq:subcritical-scale}. In both bounds, let $x$ be the logarithm of the
number of cells per observation, up to a shift of order $\log L$. The
degree in the upper bound and the frequency $M$ in the lower bound are
then about $\Delta L/x$. This keeps the variance at $n^{-2\theta}$, up
to factors $e^{O(\sqrt L)}$, in the upper bound, and it makes the
coefficients with $j\ge M$ negligible in the lower bound. Relative to
$n^{-\theta}$, the cells cost $e^{-\theta x}$ and the reciprocal
approximation costs $\rho^{\Delta L/x}$. With
$x_*=\sqrt{\tau\Delta L/\theta}$, the leading exponential cost
satisfies
\[
 \theta x+\frac{\tau\Delta L}x
 =\kappa\sqrt L+\frac{\theta(x-x_*)^2}x .
\]
In the lower bound, $M$ is the nearest even integer to
$M_*=\Delta L/x_*$. Since
$\theta\Delta L/M+\tau M-\kappa\sqrt L=\tau(M-M_*)^2/M$, the
rounding costs a factor $e^{O(1/M)}$. A balance between a loss
exponential in the matching order and a loss exponential in $\log n$
divided by that order also underlies the lower bound
of~\cite{aronowlopatto}. In the upper bound, the positive shifted values of $x$ on the capped
resolutions are $O(\sqrt L)$ and pairwise separated by at least $\log2$,
so only $O(L^{1/4})$ resolutions contribute comparable errors
(\cref{lem:window-sum}). In both bounds, the shift of $x$
contributes $L^{\theta/2}$. The lattice gain contributes a factor of
order $L$ to the lower bound, and the factor $(m+1)^\nu$, with truncation order
$m=O(\sqrt L)$, contributes $L^{\nu/2}$ to the upper bound. We obtain
$a_nL$ and $a_nL^{\nu/2+1/4}$, which leaves the factor $L^{1/4}$ when
$\nu=2$.

\section{Related work}\label{sec:related}

\emph{Higher-order influence functions and rough design.} The HOIF
estimators of Robins et al.~\cite{hoif2008} are built from projection
kernels, which they write through inverse Gram
matrices~\cite[p.~357]{hoif2008}, and the truncation bias of the
projected functional is the integral of the product of two projection
residuals~\cite[Theorems~3.7--3.8]{hoif2008}. When $g$ is unknown,
these estimators start from an initial estimate of $g$ and add U-statistic terms of increasing
order, which form a Neumann-type expansion of the inverse of a Gram
matrix about its value under that
estimate~\cite[Theorems~3.14--3.15]{hoif2008}. Robins et
al.~\cite[p.~381]{hoif2008} conjecture that, for $s<d/4$, such
estimators attain the rate $n^{-2s/d}$ in probability, up to a power of
$\log n$. The conjecture assumes an initial estimate of $g$ that is
consistent in supremum norm at a rate $1/r(n)$, an order at least a
multiple of $\log n/\log r(n)$, and a number of basis functions chosen
according to the order. In their Remark~4.1, which concerns the mean of
the square of a regression function, they observe that the empirical
Gram matrix of more than $n$ basis functions is singular. They give this
as the reason why the minimax rate is slower than $n^{-4s/(4s+d)}$
without restrictions on the covariate
density~\cite[Remark~4.1]{hoif2008}. In their Remark~4.2, they refer
to $n^{-2s/d}$, for $s<d/4$, as their conjectured minimax rate for the
variance-weighted average treatment effect when the covariate density
is not assumed smooth and the treatment effect is not assumed
constant~\cite[Remark~4.2]{hoif2008}. This effect is the
overlap-weighted effect of \cref{tab:applications}, and we compare our
result for it with their conjecture in \cref{sec:overlap-effect}.

Our estimator has a form similar to the estimators in the conjecture on
p.~381 of~\cite{hoif2008}: it is a sum of polynomials in cell moments,
estimated without bias by U-statistics, with orders up to a multiple of
$\log n$ at coarse resolutions. It differs from them in two respects. It
approximates the inverses of the local Gram matrices of $g$ by their
Chebyshev series about the known interval $[g_-,g_+]$, and not about an
initial estimate (\cref{sec:proof-tools}). Further, its order decreases
as the resolution increases. Keeping high orders only at coarse
resolutions resembles the truncation of the higher-order kernels
in~\cite[Section~9, (9.6)]{robins} (\cref{sec:tuning}). Since our
estimator has root-mean-square error $o(n^{-2s/d})$ uniformly over the
class, it attains the rate of that conjecture, although it is not one
of the estimators to which the conjecture refers. The empirical HOIF
estimators of~\cite{liu} avoid estimating $g$ by using the inverse of a
sample Gram matrix of $k=o(n)$ basis functions.

Polynomial approximation of a bilinear form in the inverse of a Gram
matrix, with each term estimated by a U-statistic, was used by Kong and
Valiant~\cite[Section~2.2]{kongvaliant2018}. In high-dimensional linear
regression with unknown covariance $\Sigma$, they estimate
$\E[YX^\top]\Sigma^{-1}\E[XY]$ through unbiased estimates of
$\E[YX^\top]\Sigma^j\E[XY]$. Chen, Liu, and
Mukherjee~\cite[Section~3.2]{chenliumukherjee2025} use Chebyshev
polynomials in $\Sigma$ for this purpose when the dimension exceeds the
sample size. We apply this strategy to the local Gram matrices of $g$ in
a nonparametric model, approximate inverses and reciprocals of
determinants jointly by a series about $[g_-,g_+]$, and choose the
degrees across resolutions to obtain the rates of \cref{thm:generic}.

\emph{Estimators that do not use density smoothness.} In dimension one,
the difference-based estimator of Robins et
al.~\cite[p.~380]{hoif2008}, which follows Wang et
al.~\cite{wangvariance}, has bias $O(n^{-2s})$ and variance $O(n^{-1})$
for the expected conditional covariance. Their bias bound combines the
H\"older continuity of the two regressions with the theory of spacings
and uses no smoothness of the density; it applies when both
smoothness indices are at most one. With these bounds, the estimator
attains the root-mean-square rate $n^{-1/2}\vee n^{-2s}$, that is, the
exponent $\min\{2s,1/2\}$ of \cref{thm:generic} when $d=1$ and both
indices are at most one, with root-$n$ risk at $s=1/4$. As mentioned
above, McClean et al.~\cite[Section~4.2 and Theorem~1]{mcclean}
construct an estimator of the expected conditional covariance whose rate
in probability is $n^{-2s/d}$, up to a polylogarithmic factor, when
$s<d/4$. Their estimator applies in every dimension and requires only
bounds on the covariate density and no knowledge of the nuisance
smoothness. The lower bound of
\cref{cor:applications} on $\cP_{\rm bil}(\alpha,\beta)$, which also
holds in probability, shows that the polynomial exponent of this rate
is minimax (\cref{sec:product-class}). Unlike the estimator of McClean
et al., ours uses $\alpha,\beta$, the bounds $g_-,g_+$, and the
remaining class parameters; it attains the exponent in
root-mean-square error and the correction in
\eqref{eq:intro-refinement} up to a polylogarithmic factor.

Above the threshold, root-$n$ estimation without density smoothness was
already known for some of these functionals. Double
cross-fitting~\cite{neweyrobins} and empirical HOIF
estimators~\cite{liu} attain semiparametric efficiency under their
respective approximation and nuisance conditions, and the bounds
of~\cite[p.~380]{hoif2008} recalled above give root-$n$ risk at
$s=d/4$ in dimension one. McClean et al.~\cite[Section~4.2 and
Theorem~1]{mcclean} obtain root-$n$ consistency of their estimator of
the expected conditional covariance when $s>d/4$, with only bounds on
the covariate density. Our estimator has root-$n$ risk for
$s\ge d/4$ in every dimension, for all smoothness indices, and for
every functional in the class of \cref{thm:generic}. In his setting,
Park~\cite[Theorem~1]{park} also obtains root-$n$ risk for $s\ge d/4$,
including $s=d/4$. Stabilized
higher-order estimators of bilinear functionals without sample
splitting were recently developed in~\cite{stabilizedhoif}, and
higher-order estimators of treatment parameters defined implicitly by
estimating equations in~\cite{generalhoif}.

\emph{Structure-agnostic models.} In the structure-agnostic models of
Balakrishnan et~al.~\cite{structureagnostic}, in which only convergence
rates are imposed on black-box nuisance estimators, first-order
estimators are minimax optimal for the expected conditional covariance
and, when the squared pilot-error radius is at least of order $n^{-1}$,
for the two quadratic functionals studied there. Jin and
Syrgkanis~\cite{jinsyrgkanis2025} prove the corresponding optimality for
the average treatment effect, weighted average treatment effects, and
the average treatment effect on the treated. Their subsequent
generalization treats linear functionals of generalized regressions and
includes a structure-agnostic lower bound for the expected conditional
covariance even under a partially linear outcome
restriction~\cite[Theorems~6.1--6.2 and~7.5]{jinsyrgkanis2026general}. Bonvini
et~al.~\cite{bonvini2024smoothsa} study classes that add smoothness
constraints to structure-agnostic ones. The H\"older smoothness that we
assume is what allows higher-order estimators to improve on first-order
ones in the minimax rate. Minimax-rate optimality does not imply
admissibility: in the structure-agnostic model, Liu, Mukherjee, and
Robins~\cite[Theorems~0--2]{liumukherjeerobins2026} show, under their
other stated conditions and when the squared pilot-error radius is at
least a constant multiple of $n^{-1/2}$, that second-order estimators
asymptotically dominate first-order ones for two quadratic functionals.
For the expected conditional covariance, neither estimator
asymptotically dominates the other.

\emph{Moment matching and polynomial approximation.} Recall that, as in
minimax estimation of nonsmooth
functionals~\cite{lepski,cailow,jiaovenkathanweissman,wuyangentropy,wuyang},
we pair moment-matching priors with polynomial approximation, with the
reciprocal in place of the nonsmooth function. To estimate the
support size of a discrete distribution, Wu and Yang~\cite{wuyang} use
best polynomial approximation of the reciprocal on an interval bounded
away from zero. Our priors use the framework of the lower bounds of Robins et
al.~\cite[Theorem~2.1 and Sections~3--4]{robins2009} for these
functionals: independent, randomly signed smooth bumps on the cells of
a partition, with cell masses that do not depend on the prior
parameters, and finite response spaces. Our lower bound follows the
strategy of the lower bound of Aronow and Lopatto~\cite{aronowlopatto} for a constant conditional
variance under rough design. There, the covariate density is random
under both priors, approximate moment matching makes the local laws of
the data under the two priors close simultaneously for every number
$k\le M$ of observations in a local region, the sample is Poissonized, and $M\asymp\sqrt{\log n}$ gives a factor
$e^{-C\sqrt{\log n}}$ in the lower bound, for some constant $C$.
Relative to~\cite{aronowlopatto}, what is new here is the random phase
of the density, under which the two priors differ only at frequency $M$,
and the separation through the $M$th Fourier coefficient of the
reciprocal density. This identifies the constant $\kappa$ in
\eqref{eq:intro-refinement}, which the lower bound of~\cite{park} also
identifies for the expected conditional covariance, and the lattice phases of
\cref{sec:lattice} gain a further factor $\log n$. We compare the two
constructions in more detail at the start of \cref{sec:lower-ideas}.

\emph{Variance estimation.} In the semiparametric model
$\Var(Y\mid X)=\sigma^2$, the variance can be estimated faster under a
random design without density smoothness than under an equally spaced
fixed design~\cite{wangvariance,caivariance} when $0<s\le1$ and
$s<d/4$. This was sketched in~\cite[Section~4]{hoif2008} for
$s<\min\{1,d/4\}$ and proved in~\cite{shen,aronowlopatto}. For $s>1$ and $d>4s$, Aronow and
Lopatto~\cite{aronowlopatto} identify the minimax exponent
$2(s+1)/(d+4)$, and Dobriban et al.~\cite{dobriban2026twoscale} give a
two-scale estimator that attains it and whose analysis uses only upper
and lower bounds on the design density. We discuss these results in \cref{sec:product-class};
they do not apply to the nonparametric models of this paper.

\emph{Discrete covariates.} In a discrete analogue, Zeng
et al.~\cite{zeng} give nearly matching minimax bounds for the average
treatment effect with high-dimensional discrete covariates. Their lower
bound~\cite[Section~3.2]{zeng} uses priors on the covariate
probabilities, propensities, and regressions of the categories that
match moments up to order $\log n$, through the duality between moment
matching and best polynomial approximation. They also show that the
expected conditional covariance and a variance-weighted average
treatment effect can be estimated faster when the covariate
distribution is known, but that this knowledge may not help for the
average treatment effect~\cite[Section~5]{zeng}.

\emph{Concurrent work.} Park~\cite{park} studies the expected
conditional covariance $\E[\Cov(A,Y\mid X)]$ for responses bounded by
one in absolute value, covariates in $[0,1]^d$ with a density between
known bounds $0<g_-<1<g_+$, and conditional means $\E[A\mid X]$ and
$\E[Y\mid X]$ in H\"older classes of orders $\alpha$ and $\beta$. With
$U=A$ and $V=Y$, these are the response and density conditions of
$\cP_{\rm bil}(\alpha,\beta)$ in \cref{sec:product-class}, but the
H\"older norms differ: Park requires extensions to a fixed open
neighborhood of $[0,1]^d$ and sums the supremum norms of the
derivatives, whereas we take their maximum on the closed cube. The
present paper, in an unpublished version dated 2 October 2026,
and~\cite{park}, submitted to arXiv on 4 October 2026, were developed
independently and concurrently. Park states his bounds for the mean
squared error. His exponent $(\alpha+\beta)/d$ and his constant
$\log\{(\sqrt{g_+}+\sqrt{g_-})/(\sqrt{g_+}-\sqrt{g_-})\}$ are our
$\theta$ and $\tau$, and his dimension of the polynomials of degree at
most $\lceil t\rceil-1$ in $d$ variables is our $q_t$. In this
notation, for $0<\theta<1/2$, Park~\cite[Theorem~1]{park} bounds the
minimax mean squared error below and above by constant multiples of
\[
 n^{-2\theta}(\log n)^{\theta}e^{-4\sqrt{\tau\theta(1-2\theta)\log n}}
 \quad\text{and}\quad
 n^{-2\theta}(\log n)^{\theta+q_\alpha+q_\beta+1/2}e^{-4\sqrt{\tau\theta(1-2\theta)\log n}},
\]
and he obtains the order $n^{-1}$ when $\theta\ge1/2$, including
$\theta=1/2$. Park presents these bounds as a resolution of the
conjecture in~\cite[Remark~4.2]{hoif2008}. That remark concerns the
variance-weighted average treatment effect, a ratio whose numerator is
the expected conditional covariance~\cite[Example~1c]{hoif2008}; the
conjecture after their equation~(1.3), which we recall in
\cref{sec:intro-background}, is a different passage. Since
$2\sqrt{\tau\theta(1-2\theta)}$ is the
constant $\kappa$ of \eqref{eq:subcritical-scale}, the threshold, the
polynomial exponent, and the constant in the stretched-exponential
factor agree with those of \cref{thm:generic}. In root-mean-square
error, Park's upper bound is $C\,a_n(\log n)^{\nu/2+1/4}$, which is our
upper bound, and his lower bound is $c\,a_n$, which is the bound that
our priors with block-constant phases give (\cref{rem:gap}). Our lower
bound $c\,a_n\log n$ is larger by the factor $\log n$ that the lattice
refinement of \cref{sec:lattice} gains, and we deduce that it also
holds on Park's class. Since his class is contained in
$\cP_{\rm bil}(\alpha,\beta)$ for a suitable radius, our upper bound
holds on it as well. The hard laws that we use for this target
(\cref{sec:products-proof}) have Rademacher responses with conditional
means $u$ and $v$, and $C^\infty$ densities with values in
$[g_-,g_+]$. The functions $u,v$ are $C^\infty$ and vanish outside a
fixed cube in the interior of $[0,1]^d$. Their extensions by zero
satisfy Park's H\"older conditions once the constant $\epsilon$ of
\cref{prop:local-lower} is small enough, so these laws lie in his
class, and \cref{prop:local-lower} gives the lower bound
$c\,a_n\log n$ there. The unresolved factor in root-mean-square error
is $(\log n)^{\nu/2-3/4}$ here and $(\log n)^{\nu/2+1/4}$
in~\cite{park}, and neither work determines the power of $\log n$ in
the minimax rate.

The two proofs share their main ingredients. Both use polynomial
approximation of the reciprocal on $[g_-,g_+]$, for the bias of the
estimator and for the separation in the lower bound. Both upper bounds
write a projection error as a telescoping sum of increments over
nested partitions, remove the local Taylor polynomials of the two
regressions from each increment, approximate inverses of local Gram
matrices by matrix polynomials, estimate the resulting polynomials
without bias by products over distinct observations, and lower the
degree as the partition is refined. Park refines cubes into $2^d$
children and subtracts his corrections from a double cross-fit local
polynomial estimator; the result is clipped to $[-1,1]$. We bisect one
coordinate at a time and approximate $\int abg$ directly by
polynomials in cell moments, without regression fits. Both lower
bounds randomize the density while keeping the masses of fixed regions
constant, which are cubes in~\cite{park} and pairs of blocks here, and
both choose the priors so that regions holding few observations do not
distinguish them. The main text of~\cite{park} does not describe the
lower-bound construction in more detail.

Park also shows that a double cross-fit local polynomial estimator
whose degree depends only on $d$ attains the exponent without knowledge
of $\alpha$ and $\beta$~\cite[Corollary~1]{park}. Further, his paper
states the lower bound also for conditional covariances that are
smooth and bounded away from zero and, when $\alpha=\beta$, for the
average conditional variance of a binary response. It treats the
normalized covariance $\E[\Cov(A,Y\mid X)]/\E[\Var(Y\mid X)]$, and it
extends the bounds to responses that satisfy moment conditions in place
of boundedness. By a data-driven selection, Park attains the
known-density rate at each fixed density that is piecewise constant on
a dyadic partition. He also reports an estimator that attains his
upper bound up to a larger power of $\log n$ when $(\alpha,\beta)$ is
unknown but lies in a known compact subset of $\{\alpha+\beta<d/2\}$.
The paper also contains simulations and an analysis of diamond price
data. The present paper treats the general class of
\cref{thm:generic} and the applications in \cref{tab:applications}, of
which the expected conditional variance and the overlap-weighted effect
have counterparts in~\cite{park},
refines the lower bound by the factor $\log n$, and gives lower bounds
on the local classes \eqref{eq:local-class} and for smooth densities
(\cref{rem:smooth-design}). The proofs of~\cite{park} are in a
supplement that is not part of the arXiv posting, and we could not
locate the supplement online as of 7 October 2026. Our comparison is
based on the main text of arXiv:2610.05006v1.

\section{Organization}\label{sec:intro-organization}

\Cref{sec:models,sec:applications} state the model,
\cref{thm:generic}, and its applications. \Cref{sec:upper-ideas}
proves the upper bounds. \Cref{sec:lower-ideas} gives the two-point
lemma for the root-$n$ lower bound and reduces the bound below the
threshold to two estimates
on priors constructed in \cref{sec:lattice} and embedded in the
model in \cref{sec:generic-embedding}. \Cref{sec:hard-laws}
completes the proof. \Cref{sec:conclusion} discusses open problems,
and the appendices prove the applications.
% ---------- 02-setting.tex
\chapter{Setting and main result}\label{sec:models}

\section{Functionals and examples}\label{sec:functionals}

We observe $n$ independent copies of $O=(X,Z)$ with unknown law $P$.
Here $X\in[0,1]^d$ is a covariate with unknown Lebesgue density $p$
and $Z$ is a response in a fixed measurable space.
We let $D=D(Z)$, $U=U(Z)$, $V=V(Z)$, and $W=W(Z)$ be known bounded
functions such that $0\le D\le M_0$ and $|U|,|V|\le M_0$.
We fix a known $\lambda\ne0$ and assume that $\E_P[D\mid X]$ is bounded
away from zero. The parameter of interest is
\begin{equation}\label{eq:generic-functional}
 \psi(P)=\E_PW+\lambda\E_P\!\left[
 \frac{\E_P[U\mid X]\E_P[V\mid X]}{\E_P[D\mid X]}\right].
\end{equation}
\cref{ex:mar,ex:conditional} give the choices of
$(D,U,V,W,\lambda)$ for missing-data means, average treatment effects,
and conditional covariances and variances. We write
\begin{equation}\label{eq:generic-moments}
 w=\E_P[D\mid X],\qquad
 a=\frac{\E_P[U\mid X]}w,\qquad
 b=\frac{\E_P[V\mid X]}w,\qquad g=wp.
\end{equation}
Then $\psi(P)=\E_PW+\lambda\int_{[0,1]^d}abg$. We impose H\"older
restrictions on $a$ and $b$ and upper and lower bounds on $g$.

The functional $\psi$ belongs to the doubly robust class of
Robins et al.~\cite[Section~3.1]{hoif2008} and to the mixed-bias class
of Rotnitzky, Smucler, and Robins~\cite{mixedbias}. For bounded
candidate functions $\widetilde a,\widetilde b$, the moment
$\E_P[W+\lambda\{U\widetilde b(X)+V\widetilde a(X)-D\widetilde a(X)\widetilde b(X)\}]$
differs from $\psi(P)$ by $-\lambda\int(\widetilde a-a)(\widetilde b-b)g$.
This bias vanishes when either candidate function is correct and otherwise
depends on the product of their errors. For missing-data means, the
augmented inverse-probability-weighted estimating function is a special
case of the construction of Robins, Rotnitzky, and
Zhao~\cite[equation~(9), p.~849, and Section~8]{rrz}.
The score of the average treatment effect in double/debiased machine
learning has the same product-bias property~\cite[Section~9.4,
equations~(9.4.4)--(9.4.5)]{causalmlbook}.
The upper bound uses the same product structure for projection errors.

\begin{example}[Missing-data means and average treatment effects]\label{ex:mar}
We consider a full-data distribution $\bar P$ on $(X,R,Y)$ such that
$R$ is binary and $Y\in[0,1]$. Let $P$ be the law of the coarsened
observation $O=(X,(R,RY))$. If we take $(D,U,V)=(R,1,RY)$, $W=0$, and
$\lambda=1$, then we have
\[
 w=P(R=1\mid X),\quad
 a=1/w,\quad b=\E_P[Y\mid R=1,X],\quad g=wp,\quad \psi(P)=\E_P[b(X)].
\]
Under the missing-at-random (MAR) condition
$R\perp Y\mid X$~\cite{rubin}, this equals the full-data mean
$\E_{\bar P}Y$.
For treatment observations $(X,A,Y)$ such that $A$ is binary,
we write $b_j=\E_P[Y\mid A=j,X]$ for $j=0,1$.
If we take $R=\mathbf1_{\{A=j\}}$, then we have $\psi(P)=\E_P[b_j(X)]$.
Under consistency, conditional exchangeability, and
positivity~\cite{hernanrobins}, the difference
$\E_P[b_1(X)-b_0(X)]$ is the average treatment effect (ATE).
\end{example}

\begin{example}[Conditional covariance and variance]\label{ex:conditional}
We take $D=1$ with $|U|,|V|\le1$, and we write
$m_U=\E_P[U\mid X]$ and $m_V=\E_P[V\mid X]$. Then $w=1$ and $g=p$.
If we take $W=UV$ and $\lambda=-1$, then we have
\[
 \psi(P)=\E_P[UV]-\E_P[m_U(X)m_V(X)]
        =\E_P[\Cov_P(U,V\mid X)].
\]
If we take $U=V=Y$ and $f=\E_P[Y\mid X]$, then we have
\[
 \psi(P)=\E_P[Y^2]-\E_P[f(X)^2]=\E_P[\Var_P(Y\mid X)].
\]
This is the smallest mean squared error achievable by predicting $Y$
from $X$.
Taking $W=0$ and $\lambda=1$ instead gives $\psi(P)=\E_P[f(X)^2]$.
Subtracting $(\E_PY)^2$, we obtain the explained variance
$\Var_P(f(X))$. When $\Var_P(Y)>0$, the ratio
$\Var_P(f(X))/\Var_P(Y)$ is the population $R^2$.
In a randomized trial with $P(A=1\mid X)=1/2$ and outcome
$Y_{\mathrm{obs}}\in[0,1]$, we take
$Y=(2A-1)(2Y_{\mathrm{obs}}-1)$. Then $f$ is
the conditional average treatment effect (CATE), so that $\Var_P(f(X))$
measures the heterogeneity of treatment effects~\cite{levy}.
\end{example}

The response in the trial of \cref{ex:conditional} is obtained by
centering the outcome at $1/2$ in the modified-outcome construction
discussed by Tian et al.~\cite[Section~2.1]{tian2014}, who attribute it
to Signorovitch's thesis. Its conditional mean under balanced
randomization is the CATE.

Functionals of the form \eqref{eq:generic-functional} also include some
instrumental-variable parameters, and we obtain further treatment-effect
parameters from them through ratios and differences
(\cref{cor:applications}).

\section{Model and nondegeneracy}\label{sec:exactclass}

For $t>0$, we write $k=\lceil t\rceil-1$ and $\vartheta=t-k\in(0,1]$.
For any $f:[0,1]^d\to\mathbb R$ with continuous partial derivatives
of order at most $k$, we define the H\"older norm
\[
 \|f\|_{C^t([0,1]^d)}
 =\max_{|\mathbf j|\le k}\|D^{\mathbf j} f\|_\infty
 +\max_{|\mathbf j|=k}\sup_{x\ne y\in[0,1]^d}
 \frac{|D^{\mathbf j} f(x)-D^{\mathbf j} f(y)|}{\|x-y\|_2^\vartheta}.
\]
Here $D^{\mathbf j}$ is the partial derivative with multi-index
$\mathbf j$, of total order $|\mathbf j|$.
We set $\|f\|_{C^t([0,1]^d)}=+\infty$ if these derivatives do not
exist or are not continuous, and we define the H\"older ball by
$\cH^t(H)=\{f:\|f\|_{C^t([0,1]^d)}\le H\}$.
In particular, for integer $t$ the functions in this ball have Lipschitz
partial derivatives of order $t-1$.

We fix $d\ge1$, $\alpha,\beta,H,\delta>0$, and $0<g_-<g_+<\infty$.
Let $\cP(\alpha,\beta)$ be the class of all laws on the stated
observation space such that
\begin{equation}\label{eq:generic-class}
 w\ge\delta,\qquad
 a\in\cH^\alpha(H),\qquad b\in\cH^\beta(H),\qquad
 g_-\le g\le g_+\quad\text{almost everywhere}.
\end{equation}
These are the only restrictions on the conditional response law
and the covariate density; we impose no smoothness bound on $g$ or
semiparametric shape restrictions. Since $\delta\le w\le M_0$, the covariate density
satisfies $g_-/M_0\le p\le g_+/\delta$. The functions that define the
target, the smoothness indices, the radii, and the bounds are all fixed
as $n$ grows.

For a real-valued parameter $T$ on a class $\cP$ of observation laws,
we define the minimax root-mean-square error (RMSE) by
\[
 \risk(T,\cP)=\inf_{\widehat T}\sup_{P\in\cP}
 \{\E_P[(\widehat T-T(P))^2]\}^{1/2}.
\]
The infimum is over all measurable estimators based on the $n$
observations, including randomized ones.
For the parameter \eqref{eq:generic-functional} on
$\cP(\alpha,\beta)$, we write $\risk=\risk(\psi,\cP(\alpha,\beta))$.
The constants in $\lesssim$ and $\asymp$ never depend on $n$. In
statements about $\cP(\alpha,\beta)$, they may depend on
$d,\alpha,\beta,H,\delta,g_-,g_+,M_0$, on the response functions
$D,U,V,W$, and on $\lambda$. In lower bounds, they may also depend on
the baseline law $\pi_0$ and the radius $r$ that we introduce below; in
a lemma, they may also depend on the lemma's fixed inputs.

For the lower bound, we need a law at which every constraint in
\eqref{eq:generic-class} holds strictly and perturbations of the
conditional law of $Z$ within the model move $a$ and $b$. For example,
if $U=D$, then $\psi(P)=\E_P[W+\lambda V]$ is an observable mean, which
can be estimated at rate $n^{-1/2}$. We let $\pi_0$ be a probability
law for $Z$ with finite support, and we write
\begin{equation}\label{eq:generic-baseline-moments}
 w_0=\E_{\pi_0}D,\qquad
 a_0=\frac{\E_{\pi_0}U}{w_0},\qquad
 b_0=\frac{\E_{\pi_0}V}{w_0}.
\end{equation}
If we take $X$ uniform on $[0,1]^d$ and independent of $Z\sim\pi_0$, we
obtain a joint law $P_0$ under which $w=g=w_0$, $a=a_0$, and $b=b_0$
are constant. For the lower bound, we assume the interior and
nondegeneracy condition
\begin{equation}\label{eq:generic-nondegeneracy}
 \begin{aligned}
 &\max\{\delta,g_-\}<w_0<g_+,\qquad |a_0|<H,\quad |b_0|<H,\\
 &\text{and either }\quad
 \Cov_{\pi_0}\begin{pmatrix}U-a_0D\\V-b_0D\end{pmatrix}
 \text{ is positive definite}\\
 &\phantom{\text{and }}\text{or }\quad
 U=V,\quad \alpha=\beta,\quad \Var_{\pi_0}(U-a_0D)>0.
 \end{aligned}
\end{equation}
The first line states that $P_0$ satisfies \eqref{eq:generic-class}
strictly, since a constant has H\"older norm equal to its absolute value.
The residuals $U-a_0D$ and $V-b_0D$ have mean zero under $\pi_0$;
a positive definite covariance allows us to shift their means, and hence
$a$ and $b$, independently. The second alternative covers the square
of one conditional mean, for which this covariance is singular. For
$U=D$, the first residual vanishes and neither alternative holds.

For $r>0$, we define the neighborhood
\begin{equation}\label{eq:local-class}
 \begin{aligned}
 \cP_r(\pi_0)=\bigl\{P\in\cP(\alpha,\beta):{}&
 |w-w_0|\le r,\ |a-a_0|\le r,\ |b-b_0|\le r\\
 &\text{almost everywhere on }[0,1]^d\bigr\}.
 \end{aligned}
\end{equation}
Under \eqref{eq:generic-nondegeneracy}, this class contains $P_0$, and a
lower bound on it holds on every larger class. For example, suppose
that \eqref{eq:generic-class} is supplemented by
bounds $w_-\le w\le w_+$, $a_-\le a\le a_+$, and $b_-\le b\le b_+$ with
$w_-<w_0<w_+$, $a_-<a_0<a_+$, and $b_-<b_0<b_+$, such as an overlap
condition for missing-data means. If $r$ is smaller than the six
differences $w_0-w_-,\ w_+-w_0,\ a_0-a_-,\ a_+-a_0,\ b_0-b_-,\
b_+-b_0$, then the resulting class contains $\cP_r(\pi_0)$ and is
contained in $\cP(\alpha,\beta)$.

\section{Main theorem}\label{sec:main-theorem}

The rates depend on the smoothness and on the bounds $g_-,g_+$ through
\begin{equation}\label{eq:sharp-parameters}
 s=\frac{\alpha+\beta}{2},\qquad \theta=\frac{2s}{d},\qquad
 \rho=\frac{\sqrt{g_+}-\sqrt{g_-}}{\sqrt{g_+}+\sqrt{g_-}},
 \qquad\tau=-\log\rho.
\end{equation}
The number $\rho$ is the geometric rate of the reciprocal series of
\cref{lem:reciprocal} on $[g_-,g_+]$, and $\tau$ determines the
correction to the polynomial rate. Since
$\rho=(\sqrt{g_+/g_-}-1)/(\sqrt{g_+/g_-}+1)$, it is also the rate in the
classical convergence bound for the conjugate gradient method applied to
a linear system with a symmetric positive definite matrix of condition
number $g_+/g_-$~\cite[Lecture~38]{trefethenbau}. The upper bound also depends on
\begin{equation}\label{eq:poly-dimension}
 q_t=\binom{d+\lceil t\rceil-1}{d}\quad(t>0),\qquad
 \nu=q_\alpha+q_\beta.
\end{equation}
The dimension $q_t$ counts the coefficients of a polynomial of degree
$\lceil t\rceil-1$ in $d$ variables. We have $q_t=1$ if and only if
$0<t\le1$, so that $\nu=2$ if and only if $0<\alpha,\beta\le1$. For
$0<\theta<1/2$ and $\tau>0$, we define
\begin{equation}\label{eq:subcritical-scale}
 \kappa=2\sqrt{\theta(1-2\theta)\tau},\qquad
 a_n=a_n(\theta,\tau)=n^{-\theta}(\log n)^{\theta/2}e^{-\kappa\sqrt{\log n}}.
\end{equation}
We call $\theta=1/2$, that is, $s=d/4$, the threshold.

\begin{theorem}[Minimax bounds]
\label{thm:generic}
Consider the parameter \eqref{eq:generic-functional} on the class
\eqref{eq:generic-class}.
\begin{enumerate}[label=(\alph*)]
\item There are $C>0$ and an integer $n_0\ge3$, depending only on
$d,\alpha,\beta,H,\delta,g_-,g_+,M_0,|\lambda|$, and $\|W\|_\infty$,
such that, if $\cP(\alpha,\beta)$ is nonempty, then for every integer
$n\ge n_0$
\begin{equation}\label{eq:upper-rates}
 \risk\le Cn^{-1/2}\quad\text{if }\theta\ge1/2,\qquad
 \risk\le C\,a_n(\log n)^{\nu/2+1/4}\quad\text{if }0<\theta<1/2.
\end{equation}
\item Suppose that a law $\pi_0$ for $Z$ with finite support satisfies
\eqref{eq:generic-nondegeneracy}, and fix $r>0$. There are $c>0$ and an
integer $n_0\ge3$, depending only on
$d,\alpha,\beta,H,\delta,g_-,g_+,M_0$, the response functions
$D,U,V,W$, $\lambda$, $\pi_0$, and $r$, such that the following holds
for every class $\cP'$ with
$\cP_r(\pi_0)\subseteq\cP'\subseteq\cP(\alpha,\beta)$ and every integer
$n\ge n_0$. If $\theta\ge1/2$, then $\risk(\psi,\cP')\ge cn^{-1/2}$. If
$0<\theta<1/2$, then $\risk(\psi,\cP')\ge c\,a_n\log n$ and
\begin{equation}\label{eq:sharp-probability}
 \inf_{\widehat\psi}\sup_{P\in\cP'}
 \Pp_P\!\left\{
 |\widehat\psi-\psi(P)|\ge
 2c\,a_n\log n
 \right\}\ge\frac14.
\end{equation}
\end{enumerate}
In particular, if such a $\pi_0$ exists, then for all $n\ge n_0$, with
$n_0$ the larger of the two thresholds,
\begin{equation}\label{eq:critical-rate}
 c n^{-1/2}\le\risk\le C n^{-1/2}\qquad(\theta\ge1/2),
\end{equation}
\begin{equation}\label{eq:sharp-bracket}
 c\,a_n\log n\le\risk\le C\,a_n(\log n)^{\nu/2+1/4}
 \qquad(0<\theta<1/2),
\end{equation}
and the same bounds hold for $\risk(\psi,\cP')$ for every class $\cP'$
as in~(b).
\end{theorem}

\begin{remark}[Logarithmic form and the remaining gap]\label{rem:gap}
Taking logarithms in \eqref{eq:sharp-bracket}, we obtain
\begin{equation}\label{eq:log-risk}
 \log\risk=-\theta\log n
 -2\sqrt{\theta(1-2\theta)\tau\log n}+O(\log\log n),
\end{equation}
which is \eqref{eq:intro-refinement}. The priors of
\cref{lem:lattice-priors} with $J=0$, whose phases are constant on
each block, give $\risk\ge c\,a_n$ by the proof of
\cref{prop:local-lower}, and with the upper bound this
implies \eqref{eq:log-risk}. The lattice phases, with
$J\asymp\log\log n$, gain the factor $\log n$, so that
$\risk/a_n\to\infty$, and $a_n$ is not the minimax rate up to fixed
constants. The bounds leave the factor $(\log n)^{\nu/2-3/4}$
unresolved, which is $(\log n)^{1/4}$ when $0<\alpha,\beta\le1$. The
constants may grow as $\theta\uparrow1/2$ or as the residual covariance
in \eqref{eq:generic-nondegeneracy} approaches singularity.
\end{remark}

\begin{remark}[Smooth and nearly uniform densities]
\label{rem:smooth-design}\label{rem:smooth-density}
Every density in the lower-bound constructions of this paper is
$C^\infty$, but its derivatives are not bounded uniformly in $n$. All
lower bounds therefore persist if we require that $p$ be $C^\infty$,
provided that no uniform bound on its derivatives is imposed. When
$D=1$, as in \cref{ex:conditional}, we have $w=1$ and $g=p$.
Applying \cref{thm:generic} with $g_\pm=1\pm\omega$ and
$\delta<1$, we find that, under the remaining conditions of
\eqref{eq:generic-nondegeneracy}, even the bounds
$1-\omega\le p\le1+\omega$, for any fixed $0<\omega<1$, give the
exponent $2s/d$ below the threshold. The same holds, with
$p_\pm=1\pm\omega$, for the rows of \cref{tab:applications} with
the interval $[p_-,p_+]$. The width of the interval changes only the
correction in \eqref{eq:log-risk}.
\end{remark}
% ---------- 03-applications.tex
\chapter{Applications}\label{sec:applications}

\Cref{cor:applications} and \cref{tab:applications} collect
the applications of \cref{thm:generic}, with proofs in
\crefrange{sec:application-reductions}{sec:additional-effects}.

For an interval $I=[i_-,i_+]$ with $0<i_-<i_+<\infty$, we write
$\rho_I$ and $\tau_I$ for the values of $\rho$ and $\tau$ in
\eqref{eq:sharp-parameters} computed from $[g_-,g_+]=I$. For
$\theta,\tau>0$ and $\nu\ge2$, we define
\begin{equation}\label{eq:app-rates}
 \underline a_n(\theta,\tau)=a_n(\theta,\tau)\log n,\qquad
 \overline a_n(\theta,\tau,\nu)=a_n(\theta,\tau)(\log n)^{\nu/2+1/4}
\end{equation}
if $\theta<1/2$, with $a_n(\theta,\tau)$ as in
\eqref{eq:subcritical-scale}, and
$\underline a_n(\theta,\tau)=\overline a_n(\theta,\tau,\nu)=n^{-1/2}$ if
$\theta\ge1/2$.

\begin{definition}\label{def:bracket}
Let $T$ be a real parameter on a class $\cP$ of observation laws, let
$\theta>0$, $\nu\ge2$, and let $I$ be an interval as above. We say that
$(T,\cP)$ satisfies the \emph{lower bracket} $(\theta,I)$ if there are
$c>0$ and an integer $n_0\ge3$ such that, for every integer $n\ge n_0$,
we have $\risk(T,\cP)\ge c\,\underline a_n(\theta,\tau_I)$ and, if
$\theta<1/2$,
\begin{equation}\label{eq:app-probability}
 \inf_{\widehat T}\sup_{P\in\cP}
 \Pp_P\bigl\{|\widehat T-T(P)|\ge2c\,\underline a_n(\theta,\tau_I)\bigr\}
 \ge\frac14 .
\end{equation}
We say that $(T,\cP)$ satisfies the \emph{upper bracket}
$(\theta,I,\nu)$ if there are $C>0$ and $n_0\ge3$ such that
$\risk(T,\cP)\le C\,\overline a_n(\theta,\tau_I,\nu)$ for every integer
$n\ge n_0$. It satisfies the \emph{bracket} $(\theta,I,\nu)$ if it
satisfies both.
\end{definition}

The targets and classes in \cref{tab:applications} are defined below.
Most of these classes bound the covariate density by
\begin{equation}\label{eq:design}
 0<p_-<1<p_+<\infty,\qquad p_-\le p\le p_+ .
\end{equation}
We write $\tau_p=\tau_{[p_-,p_+]}$, we set
$\kappa_p(\theta)=2\sqrt{\theta(1-2\theta)\tau_p}$ for
$0<\theta<1/2$, and we write $I_\zeta=[(1-\zeta)p_-,(1+\zeta)p_+]$ for
$\zeta\in(0,1)$.

\begin{corollary}[Applications]\label{cor:applications}%
\label{cor:mar-ate}\label{cor:covariances}\label{cor:wald}%
\label{cor:overlap-effect}
Fix $d\in\mathbb N$ and the parameters of the classes in
\cref{tab:applications}. Each pair $(T,\cP)$ of a target and a
class in the table satisfies the bracket $(2s/d,I,\nu)$ of its row,
with the following qualification and addition.
\begin{enumerate}[label=(\alph*)]
\item On $\mathcal M^{\rm sep}$ and $\mathcal T^{\rm sep}$, the lower
bracket holds with $I=[p_-,p_+]$, and the upper bracket holds with
$I_\zeta$ in place of $I$, for every fixed $\zeta\in(0,1)$. In
particular, if $s<d/4$, then
\begin{equation}\label{eq:sep-log-risk}
 \log\risk(T,\cP)=-\frac{2s}d\log n-\kappa_p(2s/d)\sqrt{\log n}
 +o(\sqrt{\log n}).
\end{equation}
\item Let $f_0$ be a fixed known measurable predictor with
$B_0=\|f_0\|_\infty<\infty$. On $\cP_{\rm quad}(s)$, the excess loss
$\E_P[(f(X)-f_0(X))^2]$ satisfies the bracket
$(2s/d,[p_-,p_+],2q_s)$ of $\bar\sigma^2$.
\end{enumerate}
The constants $c,C,n_0$ depend only on the fixed quantities that define
$T$ and $\cP$, and also on $\zeta$ in~(a); in~(b), they depend on
$f_0$ only through $B_0$.
\end{corollary}

\begin{table}[ht]
\caption{The applications in \cref{cor:applications}. Here
$\beta_*$ is $\min(\beta_0,\beta_1)$ for the ATE, $\beta_0$ for the
ATT, and $\beta_1$ for the ATU. The last column names the source of the
lower bound: \cref{thm:generic}(b), or the families of laws and the
Markov kernels constructed in
\cref{sec:causal-extension,sec:products-proof,sec:additional-effects}.}
\label{tab:applications}
\centering
\small
\setlength{\tabcolsep}{3pt}
\renewcommand{\arraystretch}{1.2}
\newcolumntype{L}[1]{>{\raggedright\hangindent=1em\hangafter=1\arraybackslash}p{#1}}
\begin{tabular*}{\linewidth}{@{\extracolsep{\fill}}L{0.215\linewidth}lllll@{}}
\toprule
Target & Class & $s$ & $I$ & $\nu$ & Lower bound from\\
\midrule
MAR mean & $\mathcal M_\eta$ & $\frac{\alpha+\beta}2$ & $[\eta,\eta^{-1}]$
 & $q_\alpha+q_\beta$ & \cref{thm:generic}(b)\\
ATE, ATT, ATU & $\mathcal T_\eta$ & $\frac{\alpha+\beta_*}2$ & $[\eta,\eta^{-1}]$
 & $q_\alpha+q_{\beta_*}$ & MAR family, kernel\\
MAR mean & $\mathcal M^{\rm sep}$ & $\frac{\alpha+\beta}2$ & $[p_-,p_+]$
 & $q_\alpha+q_\beta$ & MAR family\\
ATE, ATT, ATU & $\mathcal T^{\rm sep}$ & $\frac{\alpha+\beta_*}2$ & $[p_-,p_+]$
 & $q_\alpha+q_{\beta_*}$ & MAR family, kernel\\
\addlinespace
$\E[\Cov(U,V\mid X)]$, $\E[m_Um_V]$ & $\cP_{\rm bil}(\alpha,\beta)$
 & $\frac{\alpha+\beta}2$ & $[p_-,p_+]$ & $q_\alpha+q_\beta$
 & \cref{thm:generic}(b)\\
$\E[\Var(Y\mid X)]$, $\E[f^2]$ & $\cP_{\rm quad}(s)$ & $s$ & $[p_-,p_+]$
 & $2q_s$ & \cref{thm:generic}(b)\\
$\Var(f(X))$ & $\cP_{\rm quad}(s)$
 & $s$ & $[p_-,p_+]$ & $2q_s$ & Rademacher family\\
$R^2$ & $\cP_{\rm quad}(s;v_{\min})$
 & $s$ & $[p_-,p_+]$ & $2q_s$ & Rademacher family\\
$\E[f^2]$, $\Var(f(X))$ in a trial & $\cP_{\rm trial}(s)$ & $s$
 & $[p_-,p_+]$ & $2q_s$ & kernel from $\cP_{\rm quad}(s)$\\
\addlinespace
Wald ratio & $\mathcal I_{\eta,c_W}$ & $\frac{\alpha+\beta}2$
 & $[\eta,\eta^{-1}]$ & $q_\alpha+q_\beta$ & ATE, kernel\\
Average conditional Wald ratio & $\mathcal J_\eta$ & $\frac{\alpha+\beta}2$
 & $[\eta,\eta^{-1}]$ & $q_\alpha+q_\beta$ & MAR mean, kernel\\
Overlap-weighted effect & $\cP_{\rm OW}(\alpha,\beta)$
 & $\min\{\frac{\alpha+\beta}2,\alpha\}$ & $[p_-,p_+]$
 & $q_\alpha+q_{\min(\alpha,\beta)}$ & two families\\
\bottomrule
\end{tabular*}
\end{table}

In every row, $\risk(T,\cP)=n^{-\min\{1/2,2s/d\}+o(1)}$, and $\nu=2$
when all smoothness indices relevant to the row are at most one. When
$s<d/4$, the lower bound rules out, even up to factors $n^{o(1)}$, the
exponent $4s/(4s+d)$ that higher-order influence function estimators
attain when the density weight $g$ is known or sufficiently
smooth~\cite[Section~3.1 and Corollary~9.1]{robins}. Every density in the
lower-bound constructions is $C^\infty$, so that, as in
\cref{rem:smooth-design}, all these lower bounds persist if $p$
must be $C^\infty$ with no uniform bound on its derivatives.

\section{Missing-data means and treatment effects}\label{sec:mar-implications}

We use the missing-data observations $O=(X,(R,RY))$ of
\cref{ex:mar}, with $(D,U,V,W,\lambda)=(R,1,RY,0,1)$, so that
$w=P(R=1\mid X)$, $a=1/w$, $b=\E_P[Y\mid R=1,X]$, and the target is the
MAR mean $\E_P[b(X)]$. We fix $0<\eta<1/4$ and $H>2$, and we let
$\mathcal M_\eta(\alpha,\beta,H)$ consist of all laws of these
observations such that $Y\in[0,1]$ and
\[
 1/w\in\cH^\alpha(H),\qquad b\in\cH^\beta(H),\qquad
 \eta\le w\le1-\eta,\qquad \eta\le wp\le\eta^{-1}.
\]
This is the class $\cP(\alpha,\beta)$ with $M_0=1$, $\delta=\eta$, and
$[g_-,g_+]=[\eta,\eta^{-1}]$, so that $\rho=(1-\eta)/(1+\eta)$,
together with the overlap condition $w\le1-\eta$.

For treatment effects, we observe $O=(X,A,Y)$ with $A\in\{0,1\}$ and
$Y\in[0,1]$, and we write
\[
 \begin{gathered}
 w(x)=P(A=1\mid X=x),\qquad w_1=w,\qquad w_0=1-w,\\
 a_j=1/w_j,\qquad g_j=w_jp,\qquad b_j(x)=\E_P[Y\mid A=j,X=x]
 \qquad(j=0,1).
 \end{gathered}
\]
Here the subscripts index the treatment arms; $w_0$ and $a_0$ are not
the baseline constants of \cref{sec:models}. We let
$\mathcal T_\eta(\alpha,\beta_0,\beta_1,H)$ consist of all laws of $O$
such that
\begin{equation}\label{eq:two-arm-class}
 \eta\le w\le1-\eta,\qquad
 a_0,a_1\in\cH^\alpha(H),\qquad
 \eta\le g_0,g_1\le\eta^{-1},\qquad
 b_j\in\cH^{\beta_j}(H)\quad(j=0,1).
\end{equation}
These restrictions imply $2\eta\le p\le2\eta^{-1}$. We write
$\mu_j=\E_P[b_j(X)]$ and consider
\begin{equation}\label{eq:treatment-functionals}
 \begin{gathered}
 T_{\rm ATE}=\mu_1-\mu_0,\\
 T_{\rm ATT}=\E_P[b_1(X)-b_0(X)\mid A=1],\qquad
 T_{\rm ATU}=\E_P[b_1(X)-b_0(X)\mid A=0].
 \end{gathered}
\end{equation}
Under consistency, conditional exchangeability, and
positivity~\cite{hernanrobins}, these are the average effects in the
population, among treated units, and among untreated units. Their effective smoothness indices are
\begin{equation}\label{eq:treatment-smoothness}
 s_{\rm ATE}=\frac{\alpha+\min(\beta_0,\beta_1)}2,\qquad
 s_{\rm ATT}=\frac{\alpha+\beta_0}2,\qquad
 s_{\rm ATU}=\frac{\alpha+\beta_1}2.
\end{equation}
To bound the difference from below, we vary one arm regression and keep
the other constant (\cref{sec:causal-extension}).

If we bound the covariate density and the propensity separately, the
correction is determined by the density interval rather than by the
range of $wp$. We fix $H_0>1$ and $0<\epsilon<1/2$. We let
$\mathcal M^{\rm sep}$ consist of the laws of the missing-data
observations such that $Y\in[0,1]$, $p$ satisfies \eqref{eq:design},
$w\in\cH^\alpha(H_0)$, $\epsilon\le w\le1-\epsilon$, and
$b\in\cH^\beta(H_0)$. We let $\mathcal T^{\rm sep}$ consist of the laws
of $(X,A,Y)$ such that $Y\in[0,1]$, $p$ satisfies \eqref{eq:design},
$w=P(A=1\mid X)$ satisfies the same conditions, and
$b_j\in\cH^{\beta_j}(H_0)$ for $j=0,1$. Here $\alpha$ is the smoothness
of $w$ itself, equivalently of $1/w$ up to the H\"older radius, since
$\epsilon\le w\le1-\epsilon$. In all four
classes, the conditional law of the response is otherwise
unrestricted, and no smoothness of $p$ is assumed.

For the upper bounds on $\mathcal M^{\rm sep}$ and
$\mathcal T^{\rm sep}$, we divide $R$ and $RY$ by a smooth pilot
estimate $\widehat w$ of $w$ computed from half of the sample. This
preserves the target and replaces the weight $wp$ by $wp/\widehat w$,
which lies within a factor $1\pm\zeta$ of $p$ except on an event of
exponentially small probability (\cref{sec:broad-inclusions}).
For $\mu_j$, we apply this construction with $R=\mathbf 1_{\{A=j\}}$ and
observation probability $w_j$ as in \cref{ex:mar}; the arm-$0$
pilot estimates $1-w$. The correction computed from $I_\zeta$ tends to
$\kappa_p(2s/d)$ as $\zeta\to0$, which gives \eqref{eq:sep-log-risk}.

\section{Conditional covariances and variances}\label{sec:product-class}

We take $D=1$ as in \cref{ex:conditional}, so that $g=p$, and we
fix $H_0>1$ and the density bounds \eqref{eq:design}. Let
$\cP_{\rm bil}(\alpha,\beta)$ consist of all laws of $(X,U,V)$ that
satisfy \eqref{eq:design}, $|U|,|V|\le1$, $m_U\in\cH^\alpha(H_0)$, and
$m_V\in\cH^\beta(H_0)$, where $m_U=\E_P[U\mid X]$ and
$m_V=\E_P[V\mid X]$. For $U=V=Y$, let $\cP_{\rm quad}(s)$ consist of all
laws of $(X,Y)$ that satisfy \eqref{eq:design}, $|Y|\le1$, and
$f\in\cH^s(H_0)$, where $f=\E_P[Y\mid X]$. We write
\[
 Q_2=\E_P[f(X)^2],\qquad \bar\sigma^2=\E_P[\Var_P(Y\mid X)],\qquad
 E_*=\Var_P(f(X))=Q_2-(\E_PY)^2 .
\]
For fixed $v_{\min}\in(0,1)$, we let $\cP_{\rm quad}(s;v_{\min})$ be the
subclass of $\cP_{\rm quad}(s)$ on which $\Var_P(Y)\ge v_{\min}$, and on
it we define the population coefficient of determination
$R_*^2=E_*/\Var_P(Y)$. In a randomized trial we observe
$(X,A,Y_{\rm obs})$, and we let $\cP_{\rm trial}(s)$ consist of the laws
such that $A\in\{0,1\}$, $Y_{\rm obs}\in[0,1]$, $P(A=1\mid X)=1/2$, $p$
satisfies \eqref{eq:design}, and the conditional average treatment
effect (CATE)
$f=\E_P[Y_{\rm obs}\mid A=1,X]-\E_P[Y_{\rm obs}\mid A=0,X]$ lies in
$\cH^s(H_0)$. All these classes leave the conditional law of the
response otherwise unrestricted.

For $E_*$ and $R_*^2$, the upper bounds combine the estimator of $Q_2$
with sample means of $Y$ and $Y^2$, and the lower bounds use hard laws
on which $\Var_P(Y)\to1$. In Gaussian linear regression, the same
relations through observable moments imply that the signal strength,
the noise variance, and the proportion of explained variance are equally
hard to estimate, up to a parametric
loss~\cite[Section~1.1]{verzelengassiat2018}. In the trial, the ATE
$\E_P[f(X)]$ is an observable mean, whereas $\Var_P(f(X))$ is as hard to
estimate as $E_*$; only $f$ needs to be smooth. Confidence intervals for
this variance in experiments, which remain valid when it is zero, are
constructed in~\cite{sanchezbecerra2023}. In \cref{cor:applications}(b),
the loss $\E_P[(Y-f_0(X))^2]$ is an observable mean, but the excess loss
is as hard to estimate as $\bar\sigma^2$.

For the expected conditional covariance, the probability bound
\eqref{eq:app-probability} on $\cP_{\rm bil}(\alpha,\beta)$ shows that
the polynomial exponent in McClean et al.'s~\cite{mcclean} upper bound
in probability is minimax optimal. Evans and Jones~\cite[Section~1(a)
and Section~2]{evansjones} studied nearest-neighbor estimators of
residual moments and covariance under Lipschitz regression, regularity
conditions on the design density, and additive errors independent of
the covariates; Williamson et
al.~\cite[Section~2.3, Example~1, and Section~3, Theorems~1--2]{williamson}
studied $R^2$ inference under convergence conditions on estimated
regressions. Variance estimation has also been studied in the
semiparametric model $\Var(Y\mid X)=\sigma^2$~\cite{hoif2008,shen}.
There, when $0<s\le1$ and $s<d/4$, a rough random design permits faster
estimation than an equally spaced fixed
design~\cite{wangvariance,caivariance}. This was sketched
in~\cite[Section~4, pp.~382--383]{hoif2008} for $s<\min\{1,d/4\}$,
proved in~\cite{shen} for errors independent of $X$, and proved
in~\cite[Theorem~1.3]{aronowlopatto} for a model in which the law of the
error may depend on $X$. For $s>1$, Aronow and
Lopatto~\cite{aronowlopatto} resolved the uniform minimax formulation of
an open question of Robins, recorded by Richardson and
Rotnitzky~\cite[Section~6]{richardsonrotnitzky}. For $s>1$ and $d>4s$,
the two-scale estimator of Dobriban et
al.~\cite[Section~3.3, Proposition~3.4]{dobriban2026twoscale} has
root-mean-square error $O(n^{-2(s+1)/(d+4)})$, the rate of the upper bound
in~\cite[Theorem~1.1]{aronowlopatto}. The proof for this estimator uses
only upper and lower bounds on the covariate
density~\cite[Appendix~A.5]{dobriban2026twoscale}.
\Cref{cor:applications} concerns the fully nonparametric
model, to which these faster rates do not apply.

\section{Wald ratios and the overlap-weighted effect}\label{sec:further-applications}

We observe $(X,Z,A,Y)$ with a binary instrument $Z$, a binary treatment
$A$, and $Y\in[0,1]$.
We write $w_z(x)=P(Z=z\mid X=x)$,
$m_{Y,z}(x)=\E_P[Y\mid Z=z,X=x]$, $m_{A,z}(x)=\E_P[A\mid Z=z,X=x]$,
$\Delta_Y=m_{Y,1}-m_{Y,0}$, and $\Delta_A=m_{A,1}-m_{A,0}$. When
$\E_P[\Delta_A(X)]>0$, the \emph{Wald ratio} is
$T_{\rm W}=\E_P[\Delta_Y(X)]/\E_P[\Delta_A(X)]$. If $P(Z=1\mid X)=1/2$,
$A\le Z$, and $q=P(A=1\mid Z=1,X)$ is bounded away from zero, then
$\Delta_A=q$, and the \emph{average conditional Wald ratio} is
$T_{\rm CW}=\E_P[b(X)]$ with $b=\Delta_Y/q$. We fix $\alpha,\beta>0$,
$H>2$, $0<\eta<1/4$, and $0<c_W\le1$. We let
$\mathcal I_{\eta,c_W}(\alpha,\beta,H)$ consist of the laws such that,
for $z=0,1$,
\[
 \begin{gathered}
 \eta\le w_z\le1-\eta,\quad 1/w_z\in\cH^\alpha(H),\quad
 \eta\le w_zp\le\eta^{-1},\\
 m_{Y,z},m_{A,z}\in\cH^\beta(H),
 \quad \E_P[\Delta_A(X)]\ge c_W,
 \end{gathered}
\]
and we let $\mathcal J_\eta(\alpha,\beta,H)$ consist of the laws such
that $P(Z=1\mid X)=1/2$, $A\le Z$, and
\[
 1/q\in\cH^\alpha(H),\qquad b\in\cH^\beta(H),\qquad
 \eta\le q\le1-\eta,\qquad \eta\le qp\le\eta^{-1}.
\]

For observations $(X,A,Y)$ with a binary treatment, $Y\in[0,1]$,
propensity $w=P(A=1\mid X)$, and regression $m=\E_P[Y\mid X]$, the
\emph{overlap-weighted effect} is
\begin{equation}\label{eq:overlap-effect}
 \gamma_{\rm OW}=\frac{\E_P[\Cov_P(A,Y\mid X)]}{\E_P[\Var_P(A\mid X)]}
 =\frac{\E_P[w(1-w)(b_1-b_0)(X)]}{\E_P[w(1-w)(X)]},
\end{equation}
with $b_a=\E_P[Y\mid A=a,X]$. It is the coefficient of $A$ in the
least-squares projection of $Y$ onto functions $\gamma A+h(X)$. The first
expression in \eqref{eq:overlap-effect} is the population version of the
partialing-out construction of Robinson~\cite{robinson} for the
partially linear model, which we do not assume to hold. We let
$\cP_{\rm OW}(\alpha,\beta)$ consist of the laws that satisfy
\eqref{eq:design}, $w\in\cH^\alpha(H_0)$, $m\in\cH^\beta(H_0)$, and
$\epsilon\le w\le1-\epsilon$, with $H_0>1$ and $0<\epsilon<1/2$. In all
three classes, the conditional law of the response is otherwise
unrestricted. \Cref{sec:additional-effects} recalls the causal
interpretations of the three targets. The numerator of $\gamma_{\rm OW}$
has effective smoothness $(\alpha+\beta)/2$; its denominator depends only
on the propensity and has effective smoothness $\alpha$. The ratio has
$s_{\rm OW}=\min\{(\alpha+\beta)/2,\alpha\}$, as in \cref{tab:applications}.
% ---------- 04-upper.tex
\chapter{Upper bounds}\label{sec:upper-ideas}

We prove \cref{thm:generic}(a) by approximating $\int abg$ by
polynomials in expectations of known functions of one observation.
These have unbiased estimators, so we never estimate $g$. Our
estimator is a sum of U-statistics of growing order, like the HOIF
estimators of~\cite{hoif2008,robins}, which correct an initial estimate
by U-statistics of increasing order. Our estimator uses no initial
estimate of $g$ (\cref{sec:related}). For linear models, Kong
and Valiant~\cite{kongvaliant2018} developed a method that approximates
a quadratic form in the inverse of an unknown covariance matrix by
polynomials, whose terms are estimated without bias by averages over
distinct observations. Chen, Liu, and
Mukherjee~\cite[Section~3.2]{chenliumukherjee2025} use Chebyshev
polynomials of the reciprocal in this way. We apply this idea cell by
cell, to the local Gram matrices of $g$. We use
reciprocal series to approximate projection increments over nested
partitions, bound the variance of their unbiased estimators, and choose
the resolution and polynomial degrees.

We first reduce to $W=0$ and $\lambda=1$, so that $\psi=\int abg$. For an
estimator $\widehat T$ of $T(P)=\int abg$, by the triangle inequality, we
have
\[
 \left\|\lambda\widehat T+\frac1n\sum_{i=1}^nW(O_i)-\psi(P)\right\|_{L^2(P^n)}
 \le |\lambda|\,\|\widehat T-T(P)\|_{L^2(P^n)}
       +\|W\|_\infty n^{-1/2},
\]
and the last term is at most a fixed multiple of both upper rates. We may
also assume that $\alpha\le\beta$, since exchanging $(a,U,\alpha)$ with
$(b,V,\beta)$ changes neither $\int abg$ nor the class
$\cP(\alpha,\beta)$, and the rates are symmetric in $(\alpha,\beta)$.
We set $\ell=\lceil\beta\rceil-1$, the cellwise polynomial degree of the
estimator, and $s_0=\min(\alpha,\beta)$. For $h\le1$, $h^{s_0}$ bounds both
$h^\alpha$ and $h^\beta$. Recall from \eqref{eq:sharp-parameters} that
$\theta=(\alpha+\beta)/d$ and $\tau=-\log\rho$, and that $\nu$ is given by
\eqref{eq:poly-dimension}.

In this section, constants depend only on
$d,\alpha,\beta,H,\delta,g_-,g_+,M_0$, and, for the general target, on
$|\lambda|$ and $\|W\|_\infty$. The estimator uses these known parameters
and the values of $D,U,V,W$ at the observations. The arguments use only
the bounds $0\le D\le M_0$, $|U|,|V|\le M_0$, and $\|W\|_\infty$, with no
property of the response space or baseline law $\pi_0$.

\section{The reciprocal series}\label{sec:proof-tools}

Reciprocals enter both bounds. In the upper bound, we approximate
products of reciprocals of cell moments of $g$ by polynomials in those
moments. In the lower bound, the separation between the target means of
two priors is a Fourier coefficient of the reciprocal of an oscillating
density. In \cref{lem:reciprocal}, we record the Chebyshev series
of $1/x$ on a positive interval $I$, with geometric rate $\rho_I$; the parameter
$\rho_I^{-1}$ is that of the Bernstein ellipse through the pole at $x=0$,
as in the classical theory of Chebyshev approximation of analytic
functions~\cite[Chapter~8]{trefethen}. The Chebyshev coefficients of
$1/(z-x)$ are recorded explicitly in~\cite[Lemma~1, equation~(10), p.~598]{mathar2006}. For
$I=[g_-,g_+]$, or any positive multiple of it, $\rho_I$ equals the $\rho$
of \eqref{eq:sharp-parameters}.

\begin{lemma}[Reciprocal expansion]\label{lem:reciprocal}
Fix an interval $I=[x_-,x_+]$ with $0<x_-<x_+$, and set
\[
 c_I=\frac{x_++x_-}{2},\qquad d_I=\frac{x_+-x_-}{2},\qquad
 \bar g_I=\sqrt{x_-x_+},\qquad
 \rho_I=\frac{\sqrt{x_+}-\sqrt{x_-}}{\sqrt{x_+}+\sqrt{x_-}}.
\]
Let $T_k$ be the Chebyshev polynomial determined by
$T_k(\cos t)=\cos(kt)$. For $x\in I$ and $|z|<\rho_I^{-1}$, we define
\[
 \mathcal K_I(z;x)=1+2\sum_{k\ge1}(-\rho_I)^kT_k((x-c_I)/d_I)z^k.
\]
Then
\begin{equation}\label{eq:common-generating}
 \mathcal K_I(z;x)=\frac{1-\rho_I^2z^2}
          {1+2\rho_Iz(x-c_I)/d_I+\rho_I^2z^2}.
\end{equation}
In particular, $\mathcal K_I(1;x)=\bar g_I/x$; equivalently, writing
$x=c_I+d_I\cos t$,
\begin{equation}\label{eq:cosine-reciprocal}
 \frac1{c_I+d_I\cos t}
 =\frac1{\bar g_I}
    \left\{1+2\sum_{k\ge1}(-\rho_I)^k\cos(kt)\right\}.
\end{equation}
\end{lemma}
\begin{proof}
We write $x=c_I+d_I\cos t$ with $t$ real and $q=-\rho_Iz$, so that
$T_k((x-c_I)/d_I)=\cos(kt)$ and $|q|<1$. Summing two absolutely
convergent geometric series, we obtain the Poisson-kernel
identity~\cite[Chapter~2]{steinshakarchi}
\[
 1+\sum_{k\ge1}q^k(e^{ikt}+e^{-ikt})
 =\frac1{1-qe^{it}}+\frac1{1-qe^{-it}}-1
 =\frac{1-q^2}{1-2q\cos t+q^2},
\]
which is \eqref{eq:common-generating}. At $z=1$, the identities
$1+\rho_I^2=2\rho_Ic_I/d_I$ and $(1-\rho_I^2)d_I/(2\rho_I)=\bar g_I$ give
\eqref{eq:cosine-reciprocal}.
\end{proof}

For $I=[g_-,g_+]$, the rate $\rho_I=\rho$ determines the constant
$\kappa$ of the upper bound through $\tau=-\log\rho$. The Neumann series
$1/x=c_I^{-1}\sum_{k\ge0}(1-x/c_I)^k$, the most direct expansion of a
reciprocal, converges on $I$ only at the rate
$d_I/c_I=2\rho_I/(1+\rho_I^2)$, which exceeds $\rho_I$. If we replaced
\eqref{eq:cosine-reciprocal} by the Neumann series in the argument of
this section, we would obtain the correction with
$-\log\{2\rho/(1+\rho^2)\}<\tau$ in place of $\tau$, and the upper bound
would no longer match the lower bound. No other choice of polynomials
improves on $\rho_I$: by a classical result of
Chebyshev, the error of the best uniform approximation of $1/x$ on $I$ by
polynomials of degree $m$ is a constant multiple of
$\rho_I^m$~\cite[Theorem~2.1 and Corollary~2.2]{kraus2012}. The HOIF
estimators of Robins et al.\ with an estimated density expand
the inverse of a Gram matrix about its value at an initial
estimate~\cite[Theorems~3.14--3.15]{hoif2008}. At resolutions with order $n$
cells, each cell contains only order-one observations on average, so the
normalized cell moments of $g$ need not be uniformly consistently
estimable under the stated assumptions, and we expand about the known interval $[g_-,g_+]$
instead.

In the upper bound we use these series only on the interval $[g_-,g_+]$,
but with matrix arguments. We set $c_g=(g_-+g_+)/2$, $d_g=(g_+-g_-)/2$,
and $\bar g=\sqrt{g_-g_+}$, the values of $c_I,d_I,\bar g_I$ for this
interval, for which $\rho_I=\rho$. From now on $I$ denotes an identity
matrix, and we regard scalars as $1\times1$ matrices. For a complex square
matrix $M$, we define $A_M=(M-c_gI)/d_g$ and
\begin{equation}\label{eq:matrix-chebyshev-series}
 \begin{gathered}
 Q(z;M)=I+2\rho zA_M+\rho^2z^2I,\\
 \mathcal K(z;M)=(1-\rho^2z^2)Q(z;M)^{-1}
 =I+2\sum_{k\ge1}(-\rho)^kT_k(A_M)z^k.
 \end{gathered}
\end{equation}
The last identity holds as a power series at $z=0$: by the recurrence
$T_{k+1}(A)=2AT_k(A)-T_{k-1}(A)$, with
$T_0(A)=I$ and $T_1(A)=A$, the product of $Q(z;M)$ and the series on the
right is $(1-\rho^2z^2)I$, and $Q(0;M)=I$. The coefficient of $z^k$ is
a polynomial of degree at most $k$ in the entries of $M$. For real
symmetric $M$ with spectrum in $[g_-,g_+]$, $A_M$ has spectrum in
$[-1,1]$, so that $\|T_k(A_M)\|_{\rm op}\le1$. By orthogonal
diagonalization and \cref{lem:reciprocal}, we obtain absolute
convergence of \eqref{eq:matrix-chebyshev-series} for
$|z|<\rho^{-1}$, with $\mathcal K(1;M)=\bar gM^{-1}$.

We expand all reciprocal factors in the same variable $z$ and truncate
their combined index to keep the degree at $m$. The next lemma bounds
the approximation error and, for \cref{lem:variance}, the derivatives
uniformly in the degree. For a function $f$ on
$\mathbb C^p$, a point $\mathsf t_0$, and directions
$v_1,\ldots,v_r\in\mathbb R^p$, we write
$D^rf(\mathsf t_0)[v_1,\ldots,v_r]$ for the derivative
$\partial_{s_1}\cdots\partial_{s_r}f(\mathsf t_0+s_1v_1+\cdots+s_rv_r)$ at
$s=0$.

\begin{lemma}[Reciprocal truncations and derivative bounds]
\label{lem:matrix-approximation}\label{lem:reciprocal-bounds}
\begin{enumerate}[label=(\alph*)]
\item Fix integers $N\ge0$, $p_0\ge1$, and $p_1,\ldots,p_N\ge0$, and set
 $\nu'=p_1+\cdots+p_N$. For real symmetric matrices $M_0,\ldots,M_N$ of
 dimensions $p_0,\ldots,p_N$ and every integer $m\ge0$, we define $G_k$
 and $S_m$ by the expansion in powers of $z$
 \begin{equation}\label{eq:matrix-polynomial}
  \mathcal K(z;M_0)\prod_{i=1}^N\det\mathcal K(z;M_i)
  =\sum_{k\ge0}G_kz^k,\qquad
  S_m(M_0,\ldots,M_N)=\bar g^{-(\nu'+1)}\sum_{k=0}^mG_k,
 \end{equation}
 where a determinant of dimension zero equals one. Then $S_m$ is a matrix
 polynomial of total degree at most $m$ in the entries of
 $M_0,\ldots,M_N$. If the spectra of $M_0,\ldots,M_N$ lie in
 $[g_-,g_+]$, then
 \begin{equation}\label{eq:matrix-approximation-error}
  \begin{gathered}
  \left\|\frac{M_0^{-1}}{\det M_1\cdots\det M_N}
        -S_m(M_0,\ldots,M_N)\right\|_{\rm op}
  \le C(m+1)^{\nu'}\rho^m,\\
  C=2^{\nu'+1}\bar g^{-(\nu'+1)}\sum_{k\ge0}(k+1)^{\nu'}\rho^k.
  \end{gathered}
 \end{equation}
\item Let $\mathcal V\subset\mathbb R^p$ be bounded, and let
 $M_1(\mathsf t),\ldots,M_N(\mathsf t)$ be fixed affine maps from
 $\mathbb C^p$ into complex square matrices that are real symmetric with
 spectra in $[g_-,g_+]$ for $\mathsf t\in\mathcal V$. Let
 $\mathcal S(z;\mathsf t)=\sum_{k\ge0}\mathcal S_k(\mathsf t)z^k$ be formed
 by fixed finite sums and products of entries of the kernels
 $\mathcal K(z;M_i(\mathsf t))$, $1\le i\le N$, and of fixed polynomials
 in $\mathsf t$. Then there are
 $\varepsilon\in(0,1]$ and $C$ such that
 $|\sum_{k=0}^m\mathcal S_k(\mathsf t)|\le C$ for every $m\ge0$ and every
 $\mathsf t\in\mathbb C^p$ with $\|\mathsf t-\mathsf t_0\|_\infty<\varepsilon$
 for some $\mathsf t_0\in\mathcal V$. The constants depend only on $g_\pm$,
 the maps $M_i$, the expression defining $\mathcal S$, and the bound on
 $\mathcal V$.
\item Let $f$ be a polynomial on $\mathbb C^p$, let
 $\mathsf t_0\in\mathbb R^p$ and $\varepsilon>0$, and suppose that $|f(\mathsf t)|\le A$
 whenever $\|\mathsf t-\mathsf t_0\|_\infty<\varepsilon$. Then, for all
 $r\ge1$ and $v_1,\ldots,v_r\in\mathbb R^p$,
 \[
  |D^rf(\mathsf t_0)[v_1,\ldots,v_r]|\le A\,r!\,(e/\varepsilon)^r
     \prod_{i=1}^r\|v_i\|_\infty .
 \]
\end{enumerate}
\end{lemma}
\begin{proof}
(a) The coefficient of $z^k$ in $\mathcal K(z;M)$ has degree at most $k$
in the entries of $M$. Since a determinant is a sum of products of
entries, the coefficient of $z^k$ in $\det\mathcal K(z;M_i)$ also has
total degree at most $k$. The same holds for $G_k$, and $S_m$ has degree
at most $m$. For spectra in
$[g_-,g_+]$, diagonalizing $M_i$, we have
$\det\mathcal K(z;M_i)=\prod_{l=1}^{p_i}\mathcal K(z;\lambda_{il})$ for
$|z|<\rho^{-1}$, where the $\lambda_{il}$ are the eigenvalues of $M_i$.
The left side of \eqref{eq:matrix-polynomial} is therefore a product of
$\nu'+1$ absolutely convergent series, and in each of them the coefficient
of $z^k$ has norm at most $2\rho^k$. There are
$\binom{k+\nu'}{\nu'}\le(k+1)^{\nu'}$ ways to distribute an index $k$
among $\nu'+1$ factors, so that
$\|G_k\|_{\rm op}\le2^{\nu'+1}(k+1)^{\nu'}\rho^k$. In particular,
$\sum_kG_k$ converges absolutely, and since $\mathcal K(1;M)=\bar gM^{-1}$
and $\det(\bar gM_i^{-1})=\bar g^{p_i}/\det M_i$, its sum is
$\bar g^{\nu'+1}M_0^{-1}/(\det M_1\cdots\det M_N)$. Writing $k=m+j$ with
$j\ge1$ and using $m+j+1\le(m+1)(j+1)$, we bound the left side of
\eqref{eq:matrix-approximation-error} by
\[
 \bar g^{-(\nu'+1)}\sum_{k>m}\|G_k\|_{\rm op}
 \le2^{\nu'+1}\bar g^{-(\nu'+1)}(m+1)^{\nu'}\rho^m
 \sum_{j\ge1}(j+1)^{\nu'}\rho^j,
\]
which is at most the right side of \eqref{eq:matrix-approximation-error}.

(b) Let $r_*=\rho^{-1/2}$ and $b=(1-\sqrt\rho)^2$. If $M$ is real
symmetric with spectrum in $[g_-,g_+]$, then $Q(z;M)$ is diagonal in an
orthonormal eigenbasis of $M$, with diagonal entries
$(1+\rho ze^{i\phi})(1+\rho ze^{-i\phi})$ for real $\phi$, and these have
absolute value at least $(1-\rho r_*)^2=b$ for $|z|\le r_*$. Thus
$\|Q(z;M)^{-1}\|_{\rm op}\le b^{-1}$ for $|z|\le r_*$. Let $\Lambda\ge1$ be
such that
$\|M_i(\mathsf t)-M_i(\mathsf t_0)\|_{\rm op}\le\Lambda\|\mathsf t-\mathsf t_0\|_\infty$
for all $i$ and all complex $\mathsf t,\mathsf t_0$, and let
$\varepsilon=\min\{1,d_gb/(4\Lambda)\}$. If $\mathsf t_0\in\mathcal V$ and
$\|\mathsf t-\mathsf t_0\|_\infty<\varepsilon$, then
\[
 \|Q(z;M_i(\mathsf t))-Q(z;M_i(\mathsf t_0))\|_{\rm op}
 \le2\rho r_*\Lambda\varepsilon/d_g\le b/2,\qquad |z|\le r_*,
\]
and the Neumann series shows that $\|Q(z;M_i(\mathsf t))^{-1}\|_{\rm op}\le2/b$
and $\|\mathcal K(z;M_i(\mathsf t))\|_{\rm op}\le2(1+\rho)/b$ there. Since the
inequality $\|\mathsf t-\mathsf t_0\|_\infty<\varepsilon$ is strict, the
kernels are holomorphic on a neighborhood of $|z|\le r_*$. The fixed
polynomials are bounded on the bounded set of such $\mathsf t$, and we
conclude that $\mathcal S(\cdot;\mathsf t)$ is holomorphic and bounded by a uniform $A$
on $|z|\le r_*$. Its Taylor coefficients are $\mathcal S_k(\mathsf t)$ by
\eqref{eq:matrix-chebyshev-series}; by Cauchy's inequality, we have
$|\mathcal S_k(\mathsf t)|\le Ar_*^{-k}$. Summing over
$k\le m$, we obtain the bound $A/(1-\sqrt\rho)$.

(c) By multilinearity we may assume that $\|v_i\|_\infty=1$ for each $i$.
The polynomial $\varphi(w_1,\ldots,w_r)=f(\mathsf t_0+\sum_iw_iv_i)$ on
$\mathbb C^r$ is bounded by $A$ when $\max_i|w_i|<\varepsilon/r$, and
\[
 D^rf(\mathsf t_0)[v_1,\ldots,v_r]=\partial_{w_1}\cdots\partial_{w_r}\varphi(0).
\]
By Cauchy's formula on the torus $|w_1|=\cdots=|w_r|=\varepsilon'$ with
$\varepsilon'<\varepsilon/r$, this derivative is at most
$A\varepsilon'^{-r}$ in absolute value. Letting
$\varepsilon'\uparrow\varepsilon/r$ and using $r^r\le e^rr!$, we obtain
the claim.
\end{proof}

\section{Projection increments and their polynomial approximations}
\label{sec:increments}

We bisect $[0,1]^d$ successively in coordinates $1,\ldots,d$, and we
repeat this cycle. At level $j$, the partition $\mathcal C_j$ has $2^j$
equal-volume rectangles, each of aspect ratio at most two, and each cell
has two children at the next level. Let $\mathbb V_j$ be the space of
functions whose restriction to each cell of $\mathcal C_j$ is a polynomial
of total degree at most $\ell$, and let $\Pi_j$ be the orthogonal
projection onto $\mathbb V_j$ in $L^2(g)$. These spaces are nested. In
dimension one, they are the spaces of piecewise polynomials on dyadic
intervals that underlie the multiwavelet bases of
Alpert~\cite[Section~1.1]{alpert1993}. Haar projections of this kind
were combined with U-statistics to estimate integral functionals of a
density in~\cite{kerkyacharianpicard1996}.
We
define $\psi_j=\langle\Pi_j a,\Pi_j b\rangle_g$, where
$\langle f_1,f_2\rangle_g=\int f_1f_2g$. By orthogonality,
\begin{equation}\label{eq:main-projection}
 \psi-\psi_j=\langle a-\Pi_j a,b-\Pi_j b\rangle_g,
 \qquad |\psi-\psi_j|\le C2^{-j\theta}.
\end{equation}
Indeed, the cellwise Taylor polynomials of degrees $\lceil\alpha\rceil-1$
and $\lceil\beta\rceil-1$ belong to $\mathbb V_j$ and have supremum errors
$C2^{-j\alpha/d}$ and $C2^{-j\beta/d}$. The orthogonal projection
minimizes the $L^2(g)$ error, and the Cauchy--Schwarz inequality bounds
the inner product by $g_+$ times the product of the two supremum errors.
We note that the identity in \eqref{eq:main-projection} has the form of
the truncation bias of the HOIF estimators, which is the expectation of
a product of two projection residuals~\cite[Theorems~3.7--3.8]{hoif2008}.

The decomposition $\psi_J=\psi_0+\sum_{j=1}^J(\psi_j-\psi_{j-1})$ allows us
to use different approximation orders at different resolutions. Each
increment is a sum over parent cells of terms in integrals of $g$, $ag$,
and $bg$ against polynomials, since the projections act cell by cell.

\begin{lemma}[Polynomials for projection increments]
\label{lem:projection-polynomials}
Fix a level $j\ge1$, and set $K=2^{j-1}$ and $h=K^{-1/d}$. For each parent
cell $C\in\mathcal C_{j-1}$ there are a known observable vector $Z_C(O)$,
its moment vector $\mathsf m_C=\E_PZ_C(O)$, and a function $F_C$ such that
\[
 \psi_j-\psi_{j-1}=K^{-1}\sum_{C\in\mathcal C_{j-1}}F_C(\mathsf m_C),
 \qquad |F_C(\mathsf m_C)|\le C_0h^{\alpha+\beta},
\]
and
\begin{equation}\label{eq:cell-observable-bounds}
 \|Z_C(O)\|_2\le C_0K1_C(X),\qquad P(X\in C)\le C_0/K.
\end{equation}
For every integer $m\ge0$, there is a polynomial $F_{C,m}$ in these
moments of total degree at most $m+\nu+2$ such that
\begin{equation}\label{eq:increment-approximation}
 |F_C(\mathsf m_C)-F_{C,m}(\mathsf m_C)|
 \le C_0h^{\alpha+\beta}(m+1)^\nu\rho^m.
\end{equation}
At every population moment vector, that is, at $\mathsf m_C=\mathsf m_C(P)$
for every $P\in\cP(\alpha,\beta)$, it satisfies
\begin{align}
 |F_{C,m}(\mathsf t)|&\le C_0
   &&\text{for complex }\mathsf t\text{ with }
     \|\mathsf t-\mathsf m_C\|_\infty<1/C_0,\label{eq:complex-bound}\\
 \|\nabla F_{C,m}(\mathsf m_C)\|_2&\le C_0h^{s_0}.
       &&\label{eq:polynomial-gradient}
\end{align}
The functions $Z_C$ and the coefficients of $F_{C,m}$ depend only on the
class parameters, the cell, and $m$. The dimension of $Z_C$ and the
constant $C_0\ge1$ do not depend on $j$, $m$, $C$, or $P$. For the base
term there are a known observable vector $Z_0(O)$ and, for every $m\ge0$,
a polynomial $F_{0,m}$ of degree at most $m+2$ in $\mathsf m_0=\E_PZ_0(O)$
with $|\psi_0-F_{0,m}(\mathsf m_0)|\le C_0\rho^m$, for which
\eqref{eq:cell-observable-bounds}, \eqref{eq:complex-bound}, and
\eqref{eq:polynomial-gradient} hold with $K=h=1$.
\end{lemma}

\begin{proof}
\emph{Bases.} For an increment, let $C\in\mathcal C_{j-1}$ be a parent
cell, of volume $1/K$ and diameter at most $2\sqrt d\,h$. Its \emph{child
space} consists of the functions on $C$ whose restriction to each child is
a polynomial of total degree at most $\ell$. For $\eta=a,b$, its
\emph{$\eta$-parent space} consists of the polynomials on $C$ of total
degree at most $\lceil\alpha\rceil-1$ for $\eta=a$ and $\ell$ for
$\eta=b$, and has dimension $p_a=q_\alpha$ and $p_b=q_\beta$,
respectively. Since $\alpha\le\beta$, the $a$-parent space is contained in
the $b$-parent space, which is contained in the child space. For the base
term, we let $C=[0,1]^d$ and $K=h=1$, we let the child space consist of
the polynomials of total degree at most $\ell$, and we let both parent
spaces be $\{0\}$, so that $p_a=p_b=0$.

Let $z_C$ be a basis of the child space that is orthonormal for $K\dd x$ on
$C$, and for $\eta=a,b$, let $L_\eta^\top z_C$ be an orthonormal basis of
the $\eta$-parent space for $K\dd x$ on $C$. Then $L_\eta$ has $p_\eta$
columns, $L_\eta^\top L_\eta=I$, and the column space of $L_a$ is contained
in that of $L_b$. Each parent cell is the image of one of the $d$ reference
rectangles $[0,\frac12]^r\times[0,1]^{d-r}$, $0\le r<d$, under a dilation
followed by a translation, and it is split in coordinate $r+1$. We fix the
bases on the reference rectangles and on $[0,1]^d$, and we transport them
by these affine maps. The maps preserve polynomial degrees and the
normalized Lebesgue measures on the cell and on its children. In
particular, the
functions in $z_C$ are bounded uniformly over levels and cells, and
$L_a,L_b$ take finitely many values.

\emph{Moments.} Suppressing the cell index, we define
\begin{equation}\label{eq:main-moments}
 \begin{gathered}
 \Omega=K\E_P[1_C(X)Dz_C(X)z_C(X)^\top],\qquad
 \Gamma_\eta=L_\eta^\top\Omega L_\eta,\\
 \mathsf u=K\E_P[1_C(X)Uz_C(X)],\qquad \mathsf v=K\E_P[1_C(X)Vz_C(X)].
 \end{gathered}
\end{equation}
They equal $K\int_Cgz_Cz_C^\top$, $K\int_Cagz_C$, and $K\int_Cbgz_C$.
Since $z_C$ is orthonormal for $K\dd x$ on $C$ and $g_-\le g\le g_+$, we
have $g_-I\preceq\Omega\preceq g_+I$, and since $L_\eta^\top L_\eta=I$, the
same bounds hold for $\Gamma_\eta$ if $p_\eta\ge1$. By Bessel's inequality,
for every bounded function $f$ on $C$,
\begin{equation}\label{eq:bessel-moment}
 \Big\|K\int_Cfgz_C\Big\|_2\le\Big(K\int_Cf^2g^2\Big)^{1/2}
 \le g_+\sup_C|f|.
\end{equation}
With $f=a,b$, this gives $\|\mathsf u\|_2,\|\mathsf v\|_2\le Hg_+$. Let
$\mathsf m_C$ collect the upper-triangular entries of $\Omega$ and the
entries of $\mathsf u$ and $\mathsf v$. Then $\mathsf m_C=\E_PZ_C(O)$, where
$Z_C(O)$ collects the corresponding entries of
$K1_C(X)Dz_C(X)z_C(X)^\top$, $K1_C(X)Uz_C(X)$, and $K1_C(X)Vz_C(X)$. Since
the bases and the responses are bounded, $\|Z_C(O)\|_2\le C_0K1_C(X)$,
and $p=g/w\le g_+/\delta$ gives $P(X\in C)\le C_0/K$.

\emph{The increment.} We set
$B_\eta=L_\eta\Gamma_\eta^{-1}L_\eta^\top$ and $N_\eta=\Omega^{-1}-B_\eta$,
with $B_\eta=0$ and $\det\Gamma_\eta=1$ if the $\eta$-parent space is
$\{0\}$. For the inner product $\langle f_1,f_2\rangle_{g,C}=K\int_Cf_1f_2g$,
the Gram matrices of the bases $z_C$ and $L_\eta^\top z_C$ are $\Omega$ and
$\Gamma_\eta$, and the inner products of $z_C$ with $a$ and $b$ are
$\mathsf u$ and $\mathsf v$. By the normal equations, the projections of
$a$ onto the child space and onto the $\eta$-parent space therefore have
the coefficient vectors $\Omega^{-1}\mathsf u$ and $B_\eta\mathsf u$ in
the basis $z_C$. Thus $N_\eta\mathsf u$ represents $E_\eta a$, where
$E_\eta$ projects orthogonally onto the complement of the $\eta$-parent
space in the child space for $\langle\cdot,\cdot\rangle_{g,C}$, and
similarly for $b$ with $\mathsf v$. Since the ranges of $\Pi_j$ and
$\Pi_{j-1}$ on $C$ are the child space and the $b$-parent space, $K$
times the contribution of $C$ to $\psi_j-\psi_{j-1}$ is
\[
 F_C=\mathsf u^\top\Omega^{-1}\mathsf v-\mathsf u^\top B_b\mathsf v
    =\mathsf u^\top N_b\mathsf v.
\]
For the base term, the same formula gives
$F_C=\mathsf u^\top\Omega^{-1}\mathsf v=\psi_0$. As in the projection
kernels of the HOIF estimators~\cite[p.~357]{hoif2008}, we represent a
projection through the inverse of a Gram matrix; here the Gram matrices
are local, one for each parent cell.

\emph{Two small polynomial factors.} Both terms of $F_C$ are of order one.
For the inner product $x^\top\Omega y$ on coefficient vectors, $B_\eta\Omega$
is the orthogonal projection onto the column space of $L_\eta$, and
$N_\eta\Omega=I-B_\eta\Omega$ is the complementary projection. Since the
column space of $L_a$ is contained in that of $L_b$, we have
$(N_a\Omega)(N_b\Omega)=N_b\Omega$, that is, $N_a\Omega N_b=N_b$. We define
\[
 R_\eta=(\det\Gamma_\eta)I-\Omega L_\eta\operatorname{adj}(\Gamma_\eta)
        L_\eta^\top=(\det\Gamma_\eta)\Omega N_\eta,
\]
where $\operatorname{adj}(\Gamma_\eta)=(\det\Gamma_\eta)\Gamma_\eta^{-1}$
is the adjugate matrix, with $R_\eta=I$ if the $\eta$-parent space is
$\{0\}$, and
\begin{equation}\label{eq:main-residuals}
 r_a=R_a\mathsf u,\qquad r_b=R_b\mathsf v,\qquad
 S=\frac{\Omega^{-1}}{\det\Gamma_a\det\Gamma_b}.
\end{equation}
Since $N_a$ and $\Omega$ are symmetric,
$R_a^\top SR_b=N_a\Omega\Omega^{-1}\Omega N_b=N_b$, and therefore
$F_C=r_a^\top Sr_b$. Since $\det\Gamma_\eta$ and the entries of
$\operatorname{adj}(\Gamma_\eta)$ are polynomials of degrees $p_\eta$ and
$p_\eta-1$ in the entries of $\Omega$, the vectors $r_a$ and $r_b$ are
polynomials in $\mathsf m_C$ of degrees at most $p_a+1$ and $p_b+1$.

To see that $r_a$ is small, let $\widetilde a$ be the Taylor polynomial of
$a$ of degree $\lceil\alpha\rceil-1$ at a point of $C$, so that
$|a-\widetilde a|\le Ch^\alpha$ on $C$; for the base term, we let
$\widetilde a=0$, and then $|a-\widetilde a|\le H$. Since $\widetilde a$
lies in the $a$-parent space, $E_a\widetilde a=0$. We write
$\|x\|_\Omega^2=x^\top\Omega x$. Since $\Omega^2\preceq g_+\Omega$, we have
$\|\Omega x\|_2\le g_+^{1/2}\|x\|_\Omega$, and since
$\det\Gamma_a\le g_+^{p_a}$ and $E_a$ is a contraction,
\[
 \|r_a\|_2\le g_+^{p_a+1/2}\|N_a\mathsf u\|_\Omega
 =g_+^{p_a+1/2}\|E_a(a-\widetilde a)\|_{g,C}
 \le g_+^{p_a+1}\sup_C|a-\widetilde a|\le Ch^\alpha,
\]
where $\|f\|_{g,C}^2=\langle f,f\rangle_{g,C}$. In the same way, with
$\beta$ and the $b$-parent space, $\|r_b\|_2\le Ch^\beta$. Since
$\|S\|_{\rm op}\le g_-^{-(p_a+p_b+1)}$, we conclude that
$|F_C|\le Ch^{\alpha+\beta}$. Since we use the $a$-parent space, $S$
contains $1/\det\Gamma_a$ in place of a second copy of
$1/\det\Gamma_b$, and this keeps the power of $m+1$ below at
$\nu=q_\alpha+q_\beta$.

\emph{Polynomial approximation.} We apply
\cref{lem:matrix-approximation}(a) with $N=2$ and
$(M_0,M_1,M_2)=(\Omega,\Gamma_a,\Gamma_b)$, so that $\nu'=p_a+p_b$, and we
define
\begin{equation}\label{eq:polynomial-increment}
 F_{C,m}(\mathsf m_C)=r_a^\top
 S_m(\Omega,L_a^\top\Omega L_a,L_b^\top\Omega L_b)r_b.
\end{equation}
The substitutions $\Gamma_\eta=L_\eta^\top\Omega L_\eta$ are linear, and
the total degree of $F_{C,m}$ is therefore at most $m+p_a+p_b+2$. Since
$F_C-F_{C,m}=r_a^\top(S-S_m)r_b$, the bounds on $r_a,r_b$ and
\eqref{eq:matrix-approximation-error} imply that
$|F_C-F_{C,m}|\le Ch^{\alpha+\beta}(m+1)^{p_a+p_b}\rho^m$. For an
increment, $p_a+p_b=\nu$, which proves \eqref{eq:increment-approximation};
for the base term, $p_a=p_b=0$, which gives the degree $m+2$ and the error
$C\rho^m$.

\emph{The complex bound and the gradient.} For each of the finitely many
choices of $(L_a,L_b)$, and for the base term, let
$\mathcal V=\{\mathsf t:g_-I\preceq\Omega\preceq g_+I,\
\|\mathsf u\|_2,\|\mathsf v\|_2\le Hg_+\}$, where $\Omega,\mathsf u,\mathsf v$
are read off from $\mathsf t$ as from $\mathsf m_C$. This bounded set
contains $\mathsf m_C(P)$ for every level, cell, and
$P\in\cP(\alpha,\beta)$, and the affine maps $\Omega$, $L_a^\top\Omega L_a$,
and $L_b^\top\Omega L_b$ are real symmetric with spectra in $[g_-,g_+]$ on
it. By \eqref{eq:matrix-polynomial}, $F_{C,m}$ and the entries of $S_m$
are the truncations at order $m$ in $z$ of
\[
 \bar g^{-(\nu'+1)}r_a^\top\mathcal K(z;\Omega)r_b\,
 \det\mathcal K(z;L_a^\top\Omega L_a)\det\mathcal K(z;L_b^\top\Omega L_b)
\]
and of the entries of the corresponding matrix, in which each determinant
is a fixed sum of products of kernel entries, or one, and the entries of
$r_a,r_b$ are fixed polynomials in $\mathsf t$. By finitely many
applications of \cref{lem:matrix-approximation}(b), there exist
$\varepsilon\in(0,1]$ and $C$, uniform over levels and cells, such that
$F_{C,m}$ and the entries of $S_m$ are bounded by $C$ on the complex
$\varepsilon$-neighborhood of $\mathcal V$ for every $m$. This proves
\eqref{eq:complex-bound}. By \cref{lem:matrix-approximation}(c) with
$r=1$, the entries of $S_m$ have bounded gradients on $\mathcal V$, and so
do the fixed polynomials $r_a,r_b$. By the product rule, at
$\mathsf m_C(P)$ we therefore have
\[
 \|\nabla F_{C,m}\|_2
 \le C\bigl(\|r_a\|_2+\|r_b\|_2+\|r_a\|_2\|r_b\|_2\bigr)\le Ch^{s_0},
\]
which is \eqref{eq:polynomial-gradient}. Let $C_0\ge1$ be at least as
large as all constants above and the inverses of the radii of the complex
neighborhoods.
\end{proof}

For $0<\alpha,\beta\le1$, each projection is a $g$-weighted cell mean.
Let $C_1,C_2$ be the children of $C$, and let $x_i=2K\int_{C_i}g$,
$\mathsf u_i=2K\int_{C_i}ag$, $\mathsf v_i=2K\int_{C_i}bg$, and
$\bar x=(x_1+x_2)/2$. With $z_C=(\sqrt2\,1_{C_1},\sqrt2\,1_{C_2})^\top$ and
$L_a=L_b=(1,1)^\top/\sqrt2$, we have $\Omega=\operatorname{diag}(x_1,x_2)$,
$\mathsf u=(\mathsf u_1,\mathsf u_2)^\top/\sqrt2$,
$\Gamma_a=\Gamma_b=\bar x$, and $r_a=A_C(1,-1)^\top/(2\sqrt2)$, so that
\begin{equation}\label{eq:scalar-increment}
 F_C=\frac{A_CB_C}{4x_1x_2\bar x},\qquad
 A_C=x_2\mathsf u_1-x_1\mathsf u_2,\qquad
 B_C=x_2\mathsf v_1-x_1\mathsf v_2.
\end{equation}
The contrasts $A_C=x_1x_2(\mathsf u_1/x_1-\mathsf u_2/x_2)$ and $B_C$
multiply the two normalized masses by the differences of the $g$-weighted
child means of $a$ and $b$, respectively.

\section{Unbiased estimation and variance}\label{sec:unbiased}

Let $O_1,\ldots,O_n$ be the observations, let $Z$ be a bounded measurable
map from the observation space to $\mathbb R^p$, and let
$\mathsf m=\E_PZ(O)$. A real polynomial $F$ on $\mathbb R^p$ of degree at
most $R$ satisfies
$F(\mathsf m)=\sum_{q=0}^R D^qF(0)[\mathsf m,\ldots,\mathsf m]/q!$. For
$0\le q\le n$, we set $(n)_q=n(n-1)\cdots(n-q+1)$, with $(n)_0=1$. If
$R\le n$, we define the \emph{lift} of $F$ by
\begin{equation}\label{eq:lift}
 \widehat F=\sum_{q=0}^R\frac1{q!\,(n)_q}
   \sideset{}{^{\neq}}\sum_{i_1,\ldots,i_q\le n}
   D^qF(0)[Z(O_{i_1}),\ldots,Z(O_{i_q})],
\end{equation}
where $\sum^{\neq}$ runs over ordered $q$-tuples of pairwise distinct
indices in $\{1,\ldots,n\}$, and the term $q=0$ is $F(0)$. The lift
replaces each monomial by a U-statistic~\cite{halmos,hoeffding}, is
linear in $F$, and uses only its coefficients, $Z$, and the observations.

\begin{lemma}[Mean and covariance of the lift]\label{lem:lift}
Let $F_1,F_2$ be real polynomials on $\mathbb R^p$ of degree at most $n$,
and let $Y_l=Z(O_l')-\mathsf m$, where $O_1',O_2',\ldots$ are i.i.d.\ with
law $P$. Then $\E\widehat F_1=F_1(\mathsf m)$ and
\[
 \Cov(\widehat F_1,\widehat F_2)=\sum_{r\ge1}\frac1{r!\,(n)_r}
 \E\Bigl[D^rF_1(\mathsf m)[Y_1,\ldots,Y_r]\,
         D^rF_2(\mathsf m)[Y_1,\ldots,Y_r]\Bigr],
\]
where the sum stops at the smaller degree.
\end{lemma}
\begin{proof}
Let $Y_i=Z(O_i)-\mathsf m$, and let $F$ be $F_1$ or $F_2$. We substitute
$Z(O_i)=\mathsf m+Y_i$ in every argument in \eqref{eq:lift} and expand by
multilinearity. An ordered $r$-tuple $(k_1,\ldots,k_r)$ of pairwise
distinct indices then arises from $\binom qr(n-r)_{q-r}$ terms of order
$q$, and by the symmetry of $D^qF(0)$, each of them equals
$D^qF(0)[Y_{k_1},\ldots,Y_{k_r},\mathsf m,\ldots,\mathsf m]$. Since
$\binom qr(n-r)_{q-r}/\{q!(n)_q\}=1/\{r!(q-r)!(n)_r\}$, and since
$D^rF(\mathsf m)[y_1,\ldots,y_r]=\sum_{q\ge r}D^qF(0)[y_1,\ldots,y_r,\mathsf m,\ldots,\mathsf m]/(q-r)!$
by Taylor's formula, we obtain
\[
 \widehat F=F(\mathsf m)+\sum_{r=1}^{\deg F}\frac1{r!\,(n)_r}
 \sideset{}{^{\neq}}\sum_{k_1,\ldots,k_r\le n}
 D^rF(\mathsf m)[Y_{k_1},\ldots,Y_{k_r}].
\]
The $Y_i$ are independent with mean zero, and each kernel is linear in
each of its arguments. If an index appears in one of two tuples but not
in the other, then integrating first over that observation shows that the
product of the two corresponding terms has mean zero. In particular, every
term with $r\ge1$ has mean zero, so that $\E\widehat F_1=F_1(\mathsf m)$.
The surviving pairs have the same index set and length $r$, and there
are $(n)_r\,r!$ such pairs. By symmetry, each contributes
$\E[D^rF_1(\mathsf m)[Y_1,\ldots,Y_r]D^rF_2(\mathsf m)[Y_1,\ldots,Y_r]]/\{r!(n)_r\}^2$,
which gives the covariance identity.
\end{proof}

The expansion in this proof is the Hoeffding decomposition of the
U-statistics in \eqref{eq:lift} into degenerate U-statistics;
see~\cite[Sections~11.4 and~12.3]{vandervaart}
and~\cite[Section~5.1.5, Lemma~A]{serfling}.
Hoeffding~\cite[Sections~5--6]{hoeffding} derived exact variance and
covariance formulas by grouping pairs of index sets according to their
overlap. A covariance identity of the same form as that of
\cref{lem:lift}, for linear combinations of complete U-statistics
estimating trace polynomials, appears in~\cite[Theorem~6.2]{song2026}.

For a multilinear form $H$, the kernel $H[Y_1,\ldots,Y_r]$ is obtained
from $H[Z(O_1'),\ldots,Z(O_r')]$ by centering each argument, that is, by
applying $\prod_{l=1}^r(I-\mathcal E_l)$, where $\mathcal E_l$ integrates
the $l$th observation against $P$. These commuting orthogonal projections
in $L^2(P^{\otimes r})$ have an orthogonal projection as their product,
which implies that
\begin{equation}\label{up:centering-contraction}
 \E|H[Y_1,\ldots,Y_r]|^2
 \le\E|H[Z(O_1'),\ldots,Z(O_r')]|^2.
\end{equation}
This allows us to bound the variance using the cell support of $Z$; the
centered variables $Z_C(O)-\mathsf m_C$ themselves are not supported in
the cell $C$. The centered kernel $H[Y_1,\ldots,Y_r]$ is the component of
highest order in the Hoeffding decomposition of
$H[Z(O_1'),\ldots,Z(O_r')]$~\cite[Section~11.4]{vandervaart}, and
\eqref{up:centering-contraction} is the corresponding contraction
property of that decomposition.

\begin{lemma}[Variance of a polynomial in cell moments]\label{lem:variance}
Fix integers $n,K,R\ge1$ with $2R\le n$, a number $h\in(0,1]$, and a
partition of $[0,1]^d$ into $K$ cells. For every cell $C$, let $Z_C$ be a
bounded observable vector with $\mathsf m_C=\E_PZ_C(O)$, and let $F_C$ be
a real polynomial of degree at most $R$. Let $\widehat T$ be the lift of
$T=K^{-1}\sum_CF_C(\mathsf m_C)$, regarded as a polynomial in the moment
vector $(\mathsf m_C)_C$ with $Z=(Z_C)_C$. Suppose that, for some
$C_0\ge1$, $s_0>0$, and every cell,
\[
 \|Z_C(O)\|_2\le C_0K1_C(X),\qquad P(X\in C)\le C_0/K,
\]
that $|F_C(\mathsf t)|\le C_0$ for complex $\mathsf t$ with
$\|\mathsf t-\mathsf m_C\|_\infty<1/C_0$, and that
$\|\nabla F_C(\mathsf m_C)\|_2\le C_0h^{s_0}$. Then, with
$C_v=2e^2C_0^7$,
\begin{equation}\label{eq:variance-bound}
 \Var\widehat T\le\frac{C_vh^{2s_0}}n
   +\frac1n\sum_{r=2}^RC_v^r\,r!\left(\frac Kn\right)^{r-1}.
\end{equation}
If moreover $C_vRK/n\le1/2$, then
\begin{equation}\label{eq:small-resolution-variance}
 \Var\widehat T\le\frac{C_vh^{2s_0}}n+\frac{4C_v^2K}{n^2}.
\end{equation}
\end{lemma}
\begin{proof}
We write $\mathsf m=(\mathsf m_C)_C$ and
$H_r=D^rT(\mathsf m)[Z(O_1'),\ldots,Z(O_r')]$. By \cref{lem:lift} and
\eqref{up:centering-contraction}, the term of order $r$ in the variance
of $\widehat T$ is at most $\E|H_r|^2/\{r!(n)_r\}$, and the terms with
$r>R$ vanish. For observations $o_l=(x_l,\zeta_l)$, we have
\[
 D^rT(\mathsf m)[Z(o_1),\ldots,Z(o_r)]
 =K^{-1}\sum_CD^rF_C(\mathsf m_C)[Z_C(o_1),\ldots,Z_C(o_r)].
\]
Since the cells are disjoint and $Z_C$ is supported on $C$, the summand
for $C$ vanishes unless $x_1,\ldots,x_r\in C$, and at most one summand
is nonzero. Thus
\[
 |D^rT(\mathsf m)[Z(o_1),\ldots,Z(o_r)]|^2
 =K^{-2}\sum_C|D^rF_C(\mathsf m_C)[Z_C(o_1),\ldots,Z_C(o_r)]|^2.
\]

For $r\ge2$, by \cref{lem:reciprocal-bounds}(c) with
$\varepsilon=1/C_0$ and $A=C_0$, together with
$\|Z_C\|_\infty\le\|Z_C\|_2\le C_0K$, the summand for $C$ is at most
$C_0^2(r!)^2(eC_0)^{2r}(C_0K)^{2r}$ on the event that
$X_1',\ldots,X_r'\in C$, which has probability at most $(C_0/K)^r$. Summing
over the $K$ cells, we obtain
\[
 \E|H_r|^2
 \le C_0^2(r!)^2(e^2C_0^5)^rK^{r-1}
 \le(r!)^2(e^2C_0^7)^rK^{r-1}.
\]
For $r=1$, by the gradient bound, we have
$|\nabla F_C(\mathsf m_C)\cdot Z_C(o)|^2\le C_0^4h^{2s_0}K^21_C(x)$, and
since the cells partition $[0,1]^d$, we obtain
$\E|H_1|^2\le C_0^4h^{2s_0}\le C_vh^{2s_0}$. Since $2R\le n$, we have
$(n)_r\ge(n/2)^r$ for $r\le R$, and these bounds imply
\eqref{eq:variance-bound}. In \eqref{eq:variance-bound}, the ratio of the
$(r+1)$st to the $r$th term is $C_v(r+1)K/n\le C_vRK/n$ for $r+1\le R$. If
this is at most $1/2$, the sum over $r\ge2$ is at most twice its first
term $2C_v^2K/n^2$, which proves \eqref{eq:small-resolution-variance}.
\end{proof}

For a fixed order $r$, variances of order $n^{-1}\max\{1,(K/n)^{r-1}\}$,
with $K$ the number of basis functions, were obtained for the terms of
the HOIF estimators in~\cite[Theorems~3.20--3.21]{hoif2008}. In
\eqref{eq:variance-bound}, we keep the dependence on the order explicit,
including the factorial, since our orders grow with $n$.

\section{Choice of the degrees}\label{sec:tuning}

For a level $j\ge1$ and an order $m\ge0$, we set $K_j=2^{j-1}$, the number
of parent cells, and $h_j=K_j^{-1/d}$, and we let
$T_{j,m}(P)=K_j^{-1}\sum_{C\in\mathcal C_{j-1}}F_{C,m}(\mathsf m_C)$ be the
order-$m$ polynomial approximation of $\psi_j-\psi_{j-1}$. For the base
term, we let $K_0=h_0=1$ and $T_{0,m}(P)=F_{0,m}(\mathsf m_0)$. In
particular, $K_0=K_1=1$. We write $\widehat T_{j,m}$ for the lift of
$T_{j,m}$ and $C_v=2e^2C_0^7$, with $C_0$ from
\cref{lem:projection-polynomials}. By that lemma, if $R=m+\nu+2$
satisfies $2R\le n$, then the hypotheses of \cref{lem:variance} hold
for $T_{j,m}$ with $K=K_j$, $h=h_j$, and this $R$, including the base
term. Since
$r!\le R^{r-1}$ for $2\le r\le R$, the bound \eqref{eq:variance-bound}
then implies
\begin{equation}\label{up:variance-levels}
 \Var\widehat T_{j,m}\le\frac{C_v}n
 \Bigl(h_j^{2s_0}+\sum_{r=1}^{R-1}y^r\Bigr),\qquad
 y=\frac{C_vRK_j}n.
\end{equation}

For a terminal level $J$ and orders $m_0,\ldots,m_J$ with
$m_j+\nu+2\le n/2$, the estimator of $T(P)=\int abg$ is
$\widehat T=\sum_{j=0}^J\widehat T_{j,m_j}$, the lift of
$\sum_{j=0}^JT_{j,m_j}$. Its mean is $\sum_{j=0}^JT_{j,m_j}(P)$ by
\cref{lem:lift}. By the telescoping decomposition of $\psi_J$,
\eqref{eq:main-projection}, \eqref{eq:increment-approximation}, the
base-term bound, and $h_j^{\alpha+\beta}=K_j^{-\theta}$, we obtain
\begin{equation}\label{eq:bias-levels}
 |\E\widehat T-T(P)|\le C2^{-J\theta}
   +C_0\sum_{j=0}^Jf_j(m_j),\qquad f_j(m)=K_j^{-\theta}(m+1)^\nu\rho^m.
\end{equation}
By the $L^2$ triangle inequality, the standard deviation of $\widehat T$ is
at most the sum of the standard deviations of the $\widehat T_{j,m_j}$.

We choose the order on each level as the smaller of two values. The
first, a bias rule, makes the approximation errors decrease geometrically
from the terminal level toward the coarse levels. The second, a variance
rule, is needed only below the threshold. We keep high orders only at
coarse resolutions, and in this respect our estimator resembles that of
Robins et al.~\cite[Section~9, (9.6)]{robins}. Their kernel of each order
three or higher is a sum over the contiguous pairs of projection kernels
in its product; in each summand, that pair is truncated below a
hyperbola in the indices of the basis functions, and all other kernels
are truncated at $n$ basis functions.

On a level with $K$ parent cells
and degree $R$, the bound \eqref{up:variance-levels} is at most about
$n^{-2\theta}$ if $y^R\le n^{1-2\theta}$. If $x=\log y>0$, this limits the
degree to about $(1-2\theta)\log n/x$. Since $K=ne^x/(C_vR)$, the level
error $K^{-\theta}\rho^m$ at this degree is about
$n^{-\theta}(C_vR)^\theta\exp\{-\theta x-\tau(1-2\theta)\log n/x\}$. The
minimum of $\theta x+\tau(1-2\theta)\log n/x$ over $x>0$ is
$\kappa\sqrt{\log n}$, and the relevant levels have $R$ of order
$\sqrt{\log n}$. The three factors of $(\log n)^{\theta/2+\nu/2+1/4}$ in
$a_n(\log n)^{\nu/2+1/4}/(n^{-\theta}e^{-\kappa\sqrt{\log n}})$ come from
$(C_vR)^\theta$, which reflects the factorials in the variance, from the
factor $(m+1)^\nu$ in \eqref{eq:increment-approximation}, and from the sum
over levels. The next lemma bounds this sum. It is the only source of the
factor $(\log n)^{1/4}$ in the upper bound.

\begin{lemma}[Window sum]\label{lem:window-sum}
Let $\theta,\tau,B,\bar X>0$, and let $S\subset(0,\bar X]$ be a finite
set whose points are pairwise at distance at least $\log2$. Then
\[
 \sum_{x\in S}\exp\Bigl(-\theta x-\frac{\tau B}x\Bigr)
 \le\Bigl(2+\frac{(\pi\bar X/\theta)^{1/2}}{\log2}\Bigr)
     e^{-2\sqrt{\theta\tau B}}.
\]
\end{lemma}
\begin{proof}
Let $\xi=\sqrt{\tau B/\theta}$. Completing the square, we obtain
\[
 \theta x+\frac{\tau B}x=2\sqrt{\theta\tau B}+\frac{\theta(x-\xi)^2}x
 \ge2\sqrt{\theta\tau B}+\frac{\theta(x-\xi)^2}{\bar X},
 \qquad0<x\le\bar X.
\]
We list the points of $S$ in $[\xi,\infty)$ in increasing order and those
in $(0,\xi)$ in decreasing order. The $k$th point of either list, counted
from $k=0$, satisfies $|x-\xi|\ge k\log2$. With
$c=\theta(\log2)^2/\bar X$, the sum is therefore at most
$2e^{-2\sqrt{\theta\tau B}}\sum_{k\ge0}e^{-ck^2}$, and
$\sum_{k\ge0}e^{-ck^2}\le1+\int_0^\infty e^{-cu^2}\dd u
=1+\frac12(\pi/c)^{1/2}$.
\end{proof}

The construction below also gives estimators satisfying, for every
integer $n\ge3$,
\begin{equation}\label{eq:basic-upper}
 \sup_{P\in\cP(\alpha,\beta)}
 \|\widehat\psi_n-\psi(P)\|_{L^2(P^n)}
 \le C\{n^{-1/2}+n^{-\theta}\}.
\end{equation}
Moreover, fix $0\le M_W<\infty$ and a compact set
$\mathcal G\subset\{(x,y):0<x<y<\infty\}$.
With the remaining class parameters fixed, the bounds hold with
one constant $C$ for all $(g_-,g_+)\in\mathcal G$ and all admissible
observables with $\|W\|_\infty\le M_W$, using each pair's own
$\rho,\tau,\kappa$ and $a_n$ where applicable.
The tuning constants can be chosen common.
The same assertions hold when the observables depend on $X$.

\begin{proof}[Proof of \cref{thm:generic}(a)]
We write $L=\log n$, and we let $\bar r_n=n^{-1/2}$ if $\theta\ge1/2$ and
$\bar r_n=a_nL^{\nu/2+1/4}=n^{-\theta}L^{\theta/2+\nu/2+1/4}e^{-\kappa\sqrt L}$
if $0<\theta<1/2$. In both cases $n^{-1/2}\le C\bar r_n$; in the second,
$n^{-1/2}/\bar r_n\le e^{-(1/2-\theta)L+\kappa\sqrt L}$ is bounded. By the
reduction at the start of this section, it therefore suffices to
construct estimators $\widehat T$ of $T(P)=\int abg$ such that
$\E(\widehat T-T(P))^2\le C\bar r_n^2$ for every $P\in\cP(\alpha,\beta)$.
Below, $n_1\ge3$ is an integer, depending only on the parameters listed
in \cref{thm:generic}(a), that is large enough for every condition
that we impose for $n\ge n_1$. For $3\le n<n_1$ we use the estimator $0$,
whose error $|T(P)|\le H^2g_+$ is at most a constant multiple of
$\bar r_n$. Let $n\ge n_1$.

\emph{The two rules.} We fix $A>0$ such that $\tau A\ge\theta\log2+1$. We
write $K_*=2^J$ for the terminal resolution and $l_j=J-j+1$ for
$0\le j\le J$. Then $K_j=K_*2^{-l_j}$ for $j\ge1$ and $K_0=1=2K_*2^{-l_0}$,
so that $K_*2^{-l_j}\le K_j\le2K_*2^{-l_j}$ for every $j$. The \emph{bias
rule} is the order $\lceil Al_j\rceil$. Since
$(\lceil Al_j\rceil+1)^\nu\le(Al_j+2)^\nu$ and
$\rho^{\lceil Al_j\rceil}\le e^{-\tau Al_j}$, we have
\begin{equation}\label{eq:bias-rule-sum}
 \sum_{j=0}^Jf_j(\lceil Al_j\rceil)
 \le K_*^{-\theta}\sum_{l\ge1}(Al+2)^\nu e^{-l}=C_AK_*^{-\theta},
\end{equation}
because the $l_j$ are distinct positive integers and
$\tau A-\theta\log2\ge1$. In each case below, the order at level $j$ is
$m_j=\min\{\lceil Al_j\rceil,M_j\}$ for a \emph{variance rule}
$M_j\in\{0,1,2,\ldots\}\cup\{\infty\}$, and we call level $j$
\emph{capped} if $M_j<\lceil Al_j\rceil$. Then \eqref{eq:bias-levels},
\eqref{eq:bias-rule-sum}, and $2^{-J\theta}=K_*^{-\theta}$ imply that
\begin{equation}\label{eq:bias-two-rules}
 |\E\widehat T-T(P)|\le CK_*^{-\theta}
   +C\sum_{j\text{ capped}}f_j(M_j).
\end{equation}
The degree $R_j=m_j+\nu+2$ satisfies $R_j\le Al_j+\nu+3$. In both cases
below $J\le CL$, so that $R_j\le CL\le n/2$, and
\eqref{up:variance-levels} applies to level $j$ with
$y=y_j=C_vR_jK_j/n$.

\emph{The case $\theta\ge1/2$.} Let $c_K\in(0,1]$ satisfy
$c_K\le\{4C_v\sup_{l\ge1}(Al+\nu+3)2^{-l}\}^{-1}$, let
$J=\lfloor\log_2(c_Kn)\rfloor\ge0$, and let $M_j=\infty$ for all $j$.
Since $K_*\le c_Kn$, we have
$C_vR_jK_j/n\le2C_vc_K(Al_j+\nu+3)2^{-l_j}\le1/2$, and
\eqref{eq:small-resolution-variance} gives
$\Var\widehat T_{j,m_j}\le C_vh_j^{2s_0}/n+4C_v^2K_j/n^2$. Since
$\sum_{j\ge0}h_j^{s_0}\le1+\sum_{i\ge0}2^{-is_0/d}$ and
$\sum_{j=0}^JK_j^{1/2}\le CK_*^{1/2}\le Cn^{1/2}$, the standard deviation
of $\widehat T$ is at most $Cn^{-1/2}$. Since no level is capped,
\eqref{eq:bias-two-rules} implies that the bias is at most
$CK_*^{-\theta}\le C(c_Kn/2)^{-\theta}\le Cn^{-1/2}$.

\emph{The case $0<\theta<1/2$.} Let $\Delta=1-2\theta$, so that
$\kappa=2\sqrt{\theta\Delta\tau}$. We take
\[
 J=\Bigl\lceil\log_2n+\frac{\kappa\sqrt L}{\theta\log2}\Bigr\rceil,
 \qquad\text{so that}\qquad
 K_*^{-\theta}\le n^{-\theta}e^{-\kappa\sqrt L}\le\bar r_n,
 \quad\log\frac{K_*}n\le\frac\kappa\theta\sqrt L+\log2.
\]
We call level $j$ \emph{high} if $K_j>n/L^2$, and \emph{low} otherwise.
The levels $0$ and $1$ are low. On a high level,
$2^{l_j}=K_*/K_j<K_*L^2/n$, so that
\[
 R_j\le\bar R=A\log_2(K_*L^2/n)+\nu+3\le C\sqrt L,
\]
and there are at most $\log_2(K_*L^2/n)\le C\sqrt L$ high levels. We set
$c_L=\log(C_v\bar R)$, so that $0<c_L\le\frac12\log L+C$, and
\[
 B=\Delta L-2\kappa\sqrt L,\qquad
 x_j=\log\frac{K_j}n+c_L,\qquad
 \bar X=\log\frac{K_*}n+c_L\le C\sqrt L.
\]
We define the variance rule by
\[
 M_j=\lfloor B/x_j\rfloor-\nu-2\quad\text{if }x_j>0,\qquad
 M_j=\infty\quad\text{if }x_j\le0.
\]
If $x_j>0$, then $K_j>n/(C_v\bar R)\ge n/L^2$, so that level $j$ is high,
and $0<x_j\le\bar X$. Since $B\ge\Delta L/2$ and $B/\bar X\ge\nu+3$, we
then have $M_j\ge1$.

\emph{Bias.} By \eqref{eq:bias-two-rules} and $K_*^{-\theta}\le\bar r_n$,
it remains to bound $f_j(M_j)$ on the capped levels. Since a capped
level has $M_j<\infty$, it is high and has $x_j>0$, and we have
$M_j+1\le\lceil Al_j\rceil\le\bar R$, so that $(M_j+1)^\nu\le CL^{\nu/2}$.
Moreover $\rho^{M_j}\le e^{\tau(\nu+3)}e^{-\tau B/x_j}$ and
$K_j^{-\theta}=n^{-\theta}e^{\theta c_L}e^{-\theta x_j}
\le Cn^{-\theta}L^{\theta/2}e^{-\theta x_j}$, since
$e^{\theta c_L}=(C_v\bar R)^\theta\le CL^{\theta/2}$. Combining these
bounds, we find
\[
 f_j(M_j)\le Cn^{-\theta}L^{\theta/2+\nu/2}
     \exp\Bigl(-\theta x_j-\frac{\tau B}{x_j}\Bigr).
\]
The points $x_j$ of distinct capped levels lie in $(0,\bar X]$ and differ
by nonzero multiples of $\log2$. By \cref{lem:window-sum} and
$\bar X\le C\sqrt L$, the sum of these bounds is at most
$Cn^{-\theta}L^{\theta/2+\nu/2+1/4}e^{-2\sqrt{\theta\tau B}}$. Finally,
\[
 \kappa\sqrt L-2\sqrt{\theta\tau B}
 =\frac{2\sqrt{\theta\tau}\,(\Delta L-B)}{\sqrt{\Delta L}+\sqrt B}
 \le\frac{4\kappa\sqrt{\theta\tau}}{\sqrt\Delta}=8\theta\tau,
\]
so that the bias is at most $C\bar r_n$.

\emph{Variance.} Since $h_j\le1$, \eqref{up:variance-levels} implies that
$\Var\widehat T_{j,m_j}\le C_vR_jn^{-1}\max\{1,y_j\}^{R_j-1}$. On a low
level, we have $R_j\le CL$ and $K_j\le n/L^2$, so that $y_j\le CL^{-1}\le1$,
and the variance is at most $CL/n$. There are at most $J+1\le CL$ low
levels, and their standard deviations sum to at most $CL^{3/2}n^{-1/2}$.
On a high level, we have $R_j\le\bar R$, so that
$y_j\le C_v\bar RK_j/n=e^{x_j}$. If $x_j\le0$, then $y_j\le1$. If
$x_j>0$, then $R_j\le M_j+\nu+2\le B/x_j$, and therefore
$y_j^{R_j-1}\le e^{R_jx_j}\le e^B$. Since $e^B\ge1$, in both cases the
variance is at most
$C_v\bar Rn^{-1}e^B=C_v\bar Rn^{-2\theta}e^{-2\kappa\sqrt L}$, and the
standard deviations of the at most $C\sqrt L$ high levels sum to at most
$CL^{3/4}n^{-\theta}e^{-\kappa\sqrt L}$. Since $\nu\ge2$, this is at most
$C\bar r_n$, and so is $CL^{3/2}n^{-1/2}$, since
$L^{3/2}e^{-\Delta L/2+\kappa\sqrt L}$ is bounded.

With $\lambda$ and the observable mean restored as at the start of the
section, this proves
\cref{thm:generic}(a), and \eqref{eq:basic-upper} follows because
$L^{\theta/2+\nu/2+1/4}e^{-\kappa\sqrt L}$ is bounded below the threshold.
All constants depend only on the parameters listed in
\cref{thm:generic}(a). The estimator is an explicit function of the covariates $X_i$ and the
values of $D,U,V,W$ at the observations, and it involves no baseline law.

It remains to prove the uniformity assertions. For compact $\mathcal G$,
$g_-$ and $d_g$ are bounded below, $g_+$ is bounded above, and $\rho$ lies
in a compact subinterval of $(0,1)$. The approximation constants, the
radii of the complex neighborhoods, and the polynomial bounds in
\cref{lem:matrix-approximation,lem:projection-polynomials} can therefore be chosen uniformly,
which gives common $C_0$ and $C_v$. We take $A$ with
$A\inf_{\mathcal G}\tau\ge\theta\log2+1$, and then we choose $c_K$
uniformly. Below the threshold, $\Delta>0$ is fixed, and $\tau$ and
$\kappa$ have positive lower and finite upper bounds on $\mathcal G$, so
that $\bar R\le C\sqrt L$ and $\bar X\le C\sqrt L$ uniformly. Then every
large-sample condition, rate comparison, and finite-sample enlargement
above holds beyond one common $n_1$ and with one constant $C$. The
reduction at the start of the section costs at most $M_Wn^{-1/2}$.
Dependence of the observables on $X$ changes none of the moment
identities or observable bounds; equivalently, we may regard $(X,Z)$ as
the response.
\end{proof}
% ---------- 05-lower.tex
\chapter{Lower bounds}
\label{sec:lower-ideas}
\label{sec:generic-lower}

We prove \cref{thm:generic}(b) in this section and the next two.
The root-$n$ bound follows from a two-point argument
(\cref{sec:two-point}). Below the threshold, we use two fuzzy
hypotheses~\cite[Section~2.7.4]{tsybakov}: the prior mixtures are close in
total variation, while their target means differ by more than the
prior fluctuations. In a finite-response model, we bound the
Hellinger distance between Poissonized mixtures
(\cref{lem:poisson-comparison}) and reduce the lower bound to
coefficient and separation estimates (\cref{prop:reduction}).
\Cref{sec:lattice} constructs priors giving $c\,a_n$ with
block-constant phases and $c\,a_n\log n$ with lattice phases.
\Cref{sec:generic-embedding} embeds them in the generic model;
\cref{sec:hard-laws} establishes localization and completes the proof.

For the construction below the threshold, we follow the strategy of the
lower bound of Aronow and Lopatto~\cite[Section~1.2]{aronowlopatto} for
estimating a constant conditional variance under rough design. There,
the covariate density is random under both priors, approximate moment
matching makes the local laws of the data under the two priors close
simultaneously for every number $k\le M$ of observations in the same
local region, and the sample is Poissonized, with a reweighting of
the priors that cancels the exponential likelihood factor arising from
the random total mass of the perturbed density.
A matching order $M\asymp\sqrt{\log n}$, balanced against regions that
hold about $e^{-C(\log n)/M}$ observations on average, gives a factor
$e^{-C\sqrt{\log n}}$. Here, too, a random density hides the
difference between the priors from every configuration with fewer than
about $M$ observations in a pair of blocks, and we Poissonize the sample
and take $M\asymp\sqrt{\log n}$ to balance the two losses. Relative to~\cite{aronowlopatto}, the hiding and the separation take a
new form. Our priors share the law of a
random phase $\Theta$ of the density and differ only through sign
correlations at frequency $M$ in this phase. The two blocks of each pair
oscillate with opposite signs, so that the mass of each pair is fixed
and no reweighting is needed. The target sees the difference between the
priors through the $M$th Fourier coefficient of the reciprocal density,
which identifies the constant $\kappa$. The lattice phases of
\cref{sec:lattice} gain a further factor $\log n$.

\section{Divergences and a two-point bound}\label{sec:two-point}

For probability laws $\mathbb P,\mathbb Q$, we write
\[
 \TV(\mathbb P,\mathbb Q)=\sup_E|\mathbb P(E)-\mathbb Q(E)|,\qquad
 d_{\rm H}^2(\mathbb P,\mathbb Q)
 =\int(\sqrt{\dd\mathbb P}-\sqrt{\dd\mathbb Q})^2 .
\]
Any test between $\mathbb P$ and $\mathbb Q$ has error probabilities
that sum to at least $1-\TV(\mathbb P,\mathbb Q)$. We use the
inequality $\TV\le d_{\rm H}$ and the fact that the affinity
$1-d_{\rm H}^2/2$ is multiplicative over product
laws~\cite[Section~2.4 and Lemma~2.3]{tsybakov}. Since
$\prod_i(1-x_i)\ge1-\sum_ix_i$ for $x_i\in[0,1]$, the squared Hellinger
distance between two product laws is at most the sum of the squared
Hellinger distances between their factors. The next lemma applies the
two-point method~\cite[Sections~2.3 and~2.4.2]{tsybakov} along a
one-parameter path.

\begin{lemma}[A one-dimensional testing bound]\label{lem:parametric-path}
Suppose a class $\mathcal P$ contains probability laws $P_t$ for
$t$ in an open interval $J$, with
\[
 \frac{\dd P_t}{\dd P_*}=f_0+t f_1,\qquad
 f_0+t f_1\ge c_f>0,\qquad |f_1|\le C_f
 \quad P_*\text{-almost surely},
\]
where $P_*$ is a probability law and $c_f,C_f$ are independent
of $t\in J$. If $t\mapsto T(P_t)$ is differentiable at an
interior point $t_0$ and its derivative there is nonzero, then
$\risk(T,\mathcal P)\ge c n^{-1/2}$ for all sufficiently large
$n$. The constant $c>0$ is independent of $n$.
\end{lemma}

\begin{proof}
Let $f_t=f_0+tf_1$. Since
$(\sqrt{f_t}-\sqrt{f_s})^2=(t-s)^2f_1^2/(\sqrt{f_t}+\sqrt{f_s})^2$ and
$f_t,f_s\ge c_f$, we have
$d_{\rm H}^2(P_t,P_s)\le C_f^2(t-s)^2/(4c_f)$. Let
$t_\pm=t_0\pm a/\sqrt n$, where $a>0$ is fixed and small enough that
$C_fa/\sqrt{c_f}\le1/2$; for all large $n$, these points lie in $J$.
By the subadditivity of $d_{\rm H}^2$ over products, we have
$d_{\rm H}^2(P_{t_+}^{\otimes n},P_{t_-}^{\otimes n})\le C_f^2a^2/c_f\le1/4$,
so that the total variation distance between the product laws is at
most $1/2$. Since the derivative of the target at $t_0$ is nonzero, the
separation $D_n=|T(P_{t_+})-T(P_{t_-})|$ is at least $c'n^{-1/2}$ for a
constant $c'>0$ and all sufficiently large $n$. We threshold any
estimator at the midpoint of the two target values. The two error
probabilities of the resulting test sum to at least $1/2$. At one of
the laws, the estimator therefore has probability at least $1/4$ of an
estimation error of at least $D_n/2$, and its RMSE is at least $D_n/4$.
\end{proof}

\section{Phase priors and the Poisson comparison}
\label{sec:lower-preparation}

Throughout the lower-bound constructions, we fix $d\ge1$,
$\alpha,\beta>0$ and $0<r_-<1<r_+<\infty$, and we write
\[\begin{gathered}
 a_*=\alpha/d,\qquad b_*=\beta/d,\qquad
 \theta=a_*+b_*<\tfrac12,\\
 \rho=\frac{\sqrt{r_+}-\sqrt{r_-}}{\sqrt{r_+}+\sqrt{r_-}},\qquad
 \tau=-\log\rho,\qquad \delta_n=(\log n)^{-2}.
\end{gathered}\]
Here $r_-,r_+$ are the endpoints of the density range, and $\rho,\tau$
are as in \eqref{eq:sharp-parameters} for the interval $[r_-,r_+]$.
The constructions produce densities with the margin
$r_-+\delta_n\le p\le r_+-\delta_n$. This margin is small enough that
shrinking the density range by $O(\delta_n)$ changes $\rho^M$ only by
a factor $1+o(1)$ when $M=O(\sqrt{\log n})$. It is large enough to
absorb the change of $\E_P[D\mid X]$ caused by the perturbations
(\cref{sec:hard-laws}).
Let $\mathcal E$ be a finite response space, with conditional
probabilities, for $z\in\mathcal E$,
\begin{equation}\label{eq:finite-affine-law}
 Q_z(u,v)=q_{z0}+u q_{zu}+v q_{zv},\qquad
 q_{z0}>0,\quad \sum_zq_{z0}=1,\quad
 \sum_zq_{zu}=\sum_zq_{zv}=0.
\end{equation}
Here $u=u(x)$ and $v=v(x)$ are the perturbations. We call
$\mathsf s_u(z)=q_{zu}/q_{z0}$ and $\mathsf s_v(z)=q_{zv}/q_{z0}$ the
score functions of $u$ and $v$; they may coincide. Let $F$ be $C^4$ near
$(0,0)$ with $F_{uv}(0,0)\ne0$, and assume that the target is a
parameter that equals $T(P)=\int pF(u,v)$ on the laws that we
construct. We state the lower bounds in terms of
$a_n=a_n(\theta,\tau)$ from \eqref{eq:subcritical-scale}, with these
$\theta$ and $\tau$.

We next describe the priors. We fix a cube $\mathcal Q$ in the interior
of $[0,1]^d$, of volume $v_0$, and we partition it into $B$ congruent
subcubes, the \emph{blocks}, where $B$ is the $d$th power of an even
integer. We group the blocks into $B/2$ \emph{pairs} of adjacent
blocks, and we label the two blocks of each pair by $\ell\in\{1,2\}$.
In each block we use rescaled coordinates in $\mathcal U=[0,1]^d$, and
we identify a pair with $\mathcal U\times\{1,2\}$, which carries the
product of Lebesgue measure and counting measure. Let $M\ge2$ be an
integer, $C_0>0$, and $A_u,A_v>0$. We call two laws, the $+$ and $-$
\emph{priors}, of random functions $(p,u,v)$ on $[0,1]^d$
\emph{admissible with frequency $M$} if the following hold under both
of them.
\begin{enumerate}[label=(A\arabic*)]
\item The parameters of distinct pairs are independent.
\item The function $p$ is a smooth probability density with values in
$[r_-,r_+]$, it equals a fixed constant $p_0$ outside $\mathcal Q$, and
its mass on a pair is the same for every pair and in every realization.
\item We have $u=v=0$ outside $\mathcal Q$, $|u|\le C_0A_u$, and
$|v|\le C_0A_v$.
\item The parameters of a pair are $(\Theta,H,\sigma_u,\sigma_v)$,
where $(\Theta,H)$ has the same law under both priors, $\Theta$ is
uniform on $[0,2\pi)$ and independent of $H$, and
$\sigma_u,\sigma_v\in\{-1,1\}$. There are functions
$p_c,p_1,p_2,h_u,h_v$ on the pair, which depend on $H$ but
neither on $\Theta$ nor on the prior, such that, in block coordinates,
\[
 p=p_c+p_1\cos\Theta+p_2\sin\Theta,\qquad
 pu=A_u\sigma_uh_u,\qquad pv=A_v\sigma_vh_v.
\]
\item Conditional on $(\Theta,H)$, the signs satisfy
\[
 P_\pm\{\sigma_u=s,\sigma_v=t\mid\Theta,H\}
 =\tfrac14\{1\pm st\cos(M\Theta)\}\qquad(s,t\in\{-1,1\}).
\]
\end{enumerate}
We write $E_\pm$ for expectation under the $\pm$ prior, and $E$ for
expectation over the parameters $(\Theta,H)$ of one pair. Given a
realization, the law of one observation is that of $O=(X,Z)$, where $X$
has density $p$ and $P(Z=z\mid X=x)=Q_z(u(x),v(x))$.
In~\cite[Theorem~2.1 and Sections~3--4]{robins2009}, Robins et al.\ used
priors with independent parameters on disjoint cells whose masses do not
depend on the parameters, with a finite response space, to prove lower
bounds for the missing-data mean and the expected conditional
covariance. There the perturbations are bumps with independent random
signs, and in the missing-data case the covariate density changes with the parameters
without changing the mass of any cell.

By (A4), the unsigned products $pu$ and $pv$ do not depend on $\Theta$;
by (A5), only the sign correlation differs between the priors.
Summing over the four sign values, for integers $k_u,k_v\ge0$ and every
bounded function $f$ of $(\Theta,H)$, we obtain
\begin{equation}\label{lo:sign-difference}
 (E_+-E_-)\bigl[\sigma_u^{k_u}\sigma_v^{k_v}f(\Theta,H)\bigr]
 =2\,\mathbf 1_{\{k_u,k_v\text{ odd}\}}\,E\bigl[\cos(M\Theta)f(\Theta,H)\bigr].
\end{equation}
In particular, each prior is invariant under
$(\sigma_u,\sigma_v)\mapsto(-\sigma_u,-\sigma_v)$, which maps $(p,u,v)$
to $(p,-u,-v)$, and the $-$ prior is the image of the $+$ prior under
$\sigma_v\mapsto-\sigma_v$, which maps $(p,u,v)$ to $(p,u,-v)$.

As in~\cite{jiaovenkathanweissman,wuyangentropy,aronowlopatto}, we
Poissonize the sample: we compare the priors through the marked Poisson
process on $[0,1]^d\times\mathcal E$ whose intensity, with respect to $\dd x$
times counting measure, is $2np(x)Q_z(u(x),v(x))$. Given the
parameters, its restrictions to distinct pairs are independent. By (A2)
and since $\int p=1$, the mean value of $p$ over a pair in block
coordinates is $\bar p=\{1-p_0(1-v_0)\}/v_0$ in every realization. Let
$\Lambda=2nv_0/B$. As reference law on a pair we use the marked Poisson
process with intensity measure $\mu(\dd x,z)=\Lambda\bar p\,q_{z0}\dd x$
on $\mathcal Y=(\mathcal U\times\{1,2\})\times\mathcal E$. In block
coordinates, the process on the pair has intensity
$\Lambda p(x)Q_z(u(x),v(x))$, whose density with respect to $\mu$ is
$1+\varphi$, where
\begin{equation}\label{eq:poisson-point-factor}
 \varphi(x,z)=e_0(x)+e_u(x)\mathsf s_u(z)+e_v(x)\mathsf s_v(z),\qquad
 e_0=\frac p{\bar p}-1,\quad e_u=\frac{pu}{\bar p},\quad e_v=\frac{pv}{\bar p}.
\end{equation}
Since $\sum_zq_{zu}=\sum_zq_{zv}=0$ and the pair has mass $2\bar p$ in
block coordinates, we have $\int\varphi\dd\mu=0$ in every realization.

The next lemma is in the spirit of the affinity bound of Robins et
al.~\cite[Theorem~2.1]{robins2009} for two mixtures of product laws
under a product prior over the cells of a partition. We use the
Poisson process in place of the sample, and we expand in the factors that
the points carry.

\begin{lemma}[Comparison of the Poisson mixtures]
\label{lem:poisson-comparison}
Let two priors be admissible with frequency $M$. Assume $B\ge n$,
$A_u,A_v\le1$, and
$\sup_{x,z}(|u(x)q_{zu}|+|v(x)q_{zv}|)/q_{z0}\le1/2$.
Let $\mathbb P^{\rm Poi}_{B,\pm}$ be the mixtures, under the $\pm$
prior, of the laws of the marked Poisson process restricted to
$\mathcal Q$. For $j\ge0$, define on $(\mathcal U\times\{1,2\})^{j+2}$,
in the block coordinates of a given pair,
\[
 \Gamma_j(x,y,\zeta_1,\ldots,\zeta_j)=(E_+-E_-)
 \Bigl[(pu)(x)(pv)(y)\prod_{i=1}^j\{p(\zeta_i)-\bar p\}\Bigr],
\]
and let $\bar\Gamma_j^2$ be the maximum over the pairs of
$\|\Gamma_j\|_2^2=\|\Gamma_j\|_{L^2((\mathcal U\times\{1,2\})^{j+2})}^2$.
Then $\Gamma_j=0$ for $j<M$, and
\begin{equation}\label{eq:poisson-comparison-bound}
\begin{split}
 d_{\rm H}^2(\mathbb P^{\rm Poi}_{B,+},\mathbb P^{\rm Poi}_{B,-})
 &\le CB\Lambda^2\sum_{j\ge M}\frac{(C\Lambda)^j}{j!}\bar\Gamma_j^2\\
 &\quad+CB\Lambda^4A_u^2A_v^2(A_u^2+A_v^2)^2\frac{(C\Lambda)^M}{M!}.
\end{split}
\end{equation}
Here $C\ge1$ depends only on $r_\pm$, $v_0$, $C_0$, and the
probabilities $q_{z0},q_{zu},q_{zv}$. In particular, it is independent
of $n$, $B$, $M$, $A_u$, and $A_v$.
\end{lemma}

\begin{proof}
By (A1), each mixture is a product of mixtures over the pairs, and the
squared Hellinger distance is at most the sum of the distances for the
pairs. We fix a pair. Recall that a Poisson process with finite
intensity measure consists of a Poisson number of independent points~\cite[Propositions~3.2
and~3.5]{lastpenrose}. Conditioning on the number of points and using
$\int\varphi\dd\mu=0$, we find that the likelihood ratio
of the process on the pair with respect to the reference process, given
the parameters, is $\prod_{y\in\eta}\{1+\varphi(y)\}$, where $\eta$ is
the configuration of marked points. The mixture likelihood ratios are
therefore $f_\pm(\eta)=E_\pm\prod_{y\in\eta}\{1+\varphi(y)\}$. By the
amplitude assumption and $p/\bar p\ge r_-/r_+$, we have
$1+\varphi\ge c_\varphi$, where $c_\varphi=r_-/(2r_+)$, so that
$f_\pm\ge c_\varphi^N$ for a configuration with $N$ points. Since
$(\sqrt f-\sqrt g)^2=(f-g)^2/(\sqrt f+\sqrt g)^2$, we obtain
\begin{equation}\label{eq:hellinger-weighted}
 d_{\rm H}^2\le\tfrac14E_{\mathbb R}\bigl[c_\varphi^{-N}(f_+-f_-)^2\bigr],
\end{equation}
where $E_{\mathbb R}$ is the expectation under the reference process.

We compute the right side exactly. Let $\nu=\Pi_+-\Pi_-$ be the
difference of the laws of the parameters of the pair under the two
priors, and let $\varphi_\omega$ be the point factor at the parameters
$\omega$. Conditioning on the number of points gives
$E_{\mathbb R}\prod_{y\in\eta}h(y)=\exp\{\int(h-1)\dd\mu\}$ for bounded
$h$; for $0\le h\le1$, this is the formula for the probability
generating functional of a Poisson
process~\cite[Section~3.3 and Exercise~3.6]{lastpenrose}. With
$h=c_\varphi^{-1}(1+\varphi_\omega)(1+\varphi_{\omega'})$ and
$\int\varphi_\omega\dd\mu=\int\varphi_{\omega'}\dd\mu=0$, this gives
\[
 E_{\mathbb R}\bigl[c_\varphi^{-N}(f_+-f_-)^2\bigr]
 =e^{|\mu|(1/c_\varphi-1)}\iint
 \exp\bigl\{c_\varphi^{-1}
 \langle\varphi_\omega,\varphi_{\omega'}\rangle_{L^2(\mu)}\bigr\}
 \dd\nu(\omega)\dd\nu(\omega').
\]
We expand the exponential and use
$\langle\varphi_\omega,\varphi_{\omega'}\rangle_{L^2(\mu)}^k
=\langle\varphi_\omega^{\otimes k},\varphi_{\omega'}^{\otimes k}
\rangle_{L^2(\mu^k)}$. Since $\nu$ has total mass zero, the term $k=0$
vanishes, and we obtain the exact identity
\begin{equation}\label{eq:chaos-identity}
 E_{\mathbb R}\bigl[c_\varphi^{-N}(f_+-f_-)^2\bigr]
 =e^{|\mu|(1/c_\varphi-1)}\sum_{k\ge1}\frac{c_\varphi^{-k}}{k!}
 \|G_k\|^2_{L^2(\mu^k)},\qquad
 G_k=(E_+-E_-)\varphi^{\otimes k}.
\end{equation}
The interchanges are justified because the point factors are bounded
and $\nu$ has finite total variation. The prefactor is bounded, since
$|\mu|=2\Lambda\bar p\le4v_0r_+$ when $B\ge n$. We note that
$c_\varphi^{-N}$ is, up to the factor $e^{|\mu|(1/c_\varphi-1)}$, the
density of the Poisson process with intensity measure $\mu/c_\varphi$
relative to the reference process, and that the kernels of
$\prod_{y\in\eta}\{1+\varphi_\omega(y)\}$ under that process are the
tensor powers $\varphi_\omega^{\otimes k}$. Thus
\eqref{eq:chaos-identity} is also a special case of the Fock-space
representation of Poisson functionals~\cite[Theorem~1.1]{lastpenrose2011}.

For a labeling $\omega\in\{0,u,v\}^k$, let $k_u$ and $k_v$ be the
numbers of labels $u$ and $v$, let $j=k-k_u-k_v$, and let
$\Gamma_\omega(x_1,\ldots,x_k)=(E_+-E_-)\prod_{i=1}^ke_{\omega_i}(x_i)$.
Then $G_k=\sum_\omega s_\omega\Gamma_\omega$, where
$s_\omega(z_1,\ldots,z_k)$ is the product of $\mathsf s_{\omega_i}(z_i)$
over the $i$ with $\omega_i\ne0$. Let $\bar{\mathsf s}\ge1$ bound
$|\mathsf s_u|$ and $|\mathsf s_v|$. By the Cauchy--Schwarz inequality
over the $3^k$ labelings, and since $\sum_zq_{z0}=1$,
\begin{equation}\label{eq:labelings}
 \|G_k\|^2_{L^2(\mu^k)}\le(3\bar{\mathsf s}^2\Lambda\bar p)^k
 \sum_\omega\|\Gamma_\omega\|^2_{L^2((\mathcal U\times\{1,2\})^k)}.
\end{equation}

Only few labelings contribute. By (A4), we have
$e_u=A_u\sigma_uh_u/\bar p$ and $e_v=A_v\sigma_vh_v/\bar p$, and all
other factors are functions of $(\Theta,H)$. By
\eqref{lo:sign-difference}, $\Gamma_\omega=0$ unless $k_u$ and $k_v$
are both odd, and then $\Gamma_\omega$ is
$2A_u^{k_u}A_v^{k_v}\bar p^{-k}$ times the expectation over
$(\Theta,H)$ of the product of $\cos(M\Theta)$, of the $k_u+k_v$ factors
$h_u(x_i)$ and $h_v(x_i)$, and of the $j$ factors $p(x_i)-\bar p$.
Conditional on $H$, the factors $h_u,h_v$ do not depend on $\Theta$, and
the product of the $j$ density factors is a trigonometric polynomial of
degree at most $j$ in $\Theta$. Since $\Theta$ is uniform and
independent of $H$, we conclude that $\Gamma_\omega=0$ when $j<M$,
whatever the value of $k_u+k_v$.

If $k_u=k_v=1$, then $\Gamma_\omega$ equals $\bar p^{-k}\Gamma_j$ with
permuted arguments; in particular, $\Gamma_j=0$ for $j<M$. There are
$k(k-1)$ such labelings, and $k(k-1)/k!=1/j!$. Since $\bar p\ge r_-$,
\eqref{eq:hellinger-weighted}--\eqref{eq:labelings} imply that their
contribution to the distance for the pair is at most
$C\sum_{j\ge M}(C_2\Lambda)^{j+2}\|\Gamma_j\|_2^2/j!$, where
$C_2=3\bar{\mathsf s}^2r_+/(c_\varphi r_-^2)$. Summing over the $B/2$
pairs gives the first term of \eqref{eq:poisson-comparison-bound}.

For the remaining labelings, $q=k_u+k_v\ge4$ and $j\ge M$. Since
$|e_0|\le r_+/r_-$, $|e_u|\le C_0r_+A_u/r_-$, $|e_v|\le C_0r_+A_v/r_-$,
and $(\mathcal U\times\{1,2\})^k$ has measure $2^k$, we have
$\|\Gamma_\omega\|_2^2\le4C_3^kA_u^{2k_u}A_v^{2k_v}$ with
$C_3=2\max(1,C_0)^2(r_+/r_-)^2$. Let $\mathcal A=A_u^2+A_v^2\le2$.
Since $\binom q{k_u}\le q(q-1)\binom{q-2}{k_u-1}$ for
$1\le k_u\le q-1$, we have
\[
 \sum_{k_u=1}^{q-1}\binom q{k_u}A_u^{2k_u}A_v^{2(q-k_u)}
 \le q(q-1)A_u^2A_v^2\mathcal A^{q-2}.
\]
There are $\binom{j+q}q$ choices of the points labeled $u$ or $v$, and
$\binom{j+q}q/(j+q)!=1/(j!\,q!)$. With $z=C_2C_3r_-^2\Lambda$, these
labelings therefore contribute at most
\[
 CA_u^2A_v^2\Bigl(\sum_{j\ge M}\frac{z^j}{j!}\Bigr)
 \Bigl(\sum_{q\ge4}\frac{z^q\mathcal A^{q-2}}{(q-2)!}\Bigr)
 \le CA_u^2A_v^2\mathcal A^2\frac{z^{M+4}}{M!}e^{z(1+\mathcal A)}
\]
to the distance for the pair, where we used $(M+i)!\ge M!\,i!$ and
$(i+2)!\ge i!$. Since $\Lambda\le2v_0$, the exponential factor is
bounded, and summing over the pairs gives the second term of
\eqref{eq:poisson-comparison-bound}.
\end{proof}

\section{Reduction to two estimates}\label{sec:lower-reduction}

The priors of \cref{sec:lattice} divide each block into $R\ge1$
fine cells and multiply the amplitudes $(BR)^{-a_*}$ and $(BR)^{-b_*}$
by $m\ge1$; block-constant phases have $R=m=1$.
In \eqref{eq:coefficient-hypothesis}, $R^{1-M}$ records overlap after
integration over the fine cells, $R^{(j-M)\sqrt M}$ allows a cost
$R^{\sqrt M}$ per additional density factor, and $e^{C_EM}$ collects
the remaining losses.

\begin{proposition}[Reduction to two estimates]\label{prop:reduction}
Under the standing assumptions of \cref{sec:lower-preparation},
let $L=\log n$ and $\Delta_\theta=1-2\theta$, let $M$ be the nearest
even integer to $\sqrt{\theta\Delta_\theta L/\tau}$, and let $\rho_n$
be the parameter $\rho$ of \eqref{eq:sharp-parameters} computed for
$[r_-+\delta_n,r_+-\delta_n]$. Let $R=R_n\ge1$ and $m=m_n\ge1$ satisfy
\[
 m\le c_*M,\qquad m\le\min(R^{a_*},R^{b_*}),\qquad
 \log R=O(\log M)
\]
for a constant $c_*>0$, and fix $\epsilon_u,\epsilon_v\in(0,1]$.
Suppose that there are constants $C_0,C_\Gamma\ge1$, $C_E\ge0$, and
$c_2>0$ such that the following holds for all large $n$ and every
$B\ge n$ that is the $d$th power of an even integer. There are two
priors that are admissible with frequency $M$, with the constant $C_0$
and the amplitudes
\[
 A_u=\epsilon_u(BR)^{-a_*}m,\qquad A_v=\epsilon_v(BR)^{-b_*}m,
\]
such that, for every pair and every $j\ge0$,
\begin{equation}\label{eq:coefficient-hypothesis}
 \|\Gamma_j\|_2^2\le C_\Gamma^{j+1}A_u^2A_v^2e^{C_EM}
 R^{1-M+(j-M)\sqrt M}\mathbf 1_{j\ge M},
\end{equation}
and
\begin{equation}\label{eq:separation-hypothesis}
 \Bigl|(E_+-E_-)\int puv\Bigr|\ge c_2A_uA_v\rho_n^M.
\end{equation}
Then, for all large $n$, there is such a $B$ with $B\ge n$, for which
$A_u\le\epsilon_uB^{-a_*}$ and $A_v\le\epsilon_vB^{-b_*}$, with the
following property. Let $\mathcal P$ be any class that contains the
law of one observation for every realization of either prior with
this $B$, and let $T$ be a parameter on $\mathcal P$ with
$T(P)=\int pF(u,v)$ on these laws. Then there is $c_1>0$ such that, for
all large $n$,
\begin{equation}\label{eq:lower-testing-transfer}
 \inf_{\widehat T}\sup_{P\in\mathcal P}
 P\{|\widehat T-T(P)|\ge c_1m^2a_n\}\ge\frac38 .
\end{equation}
The infimum includes randomized estimators. Writing $B_n$ for this
choice of $B$, we also have $B_n/\{n(\log n)^q\}\to\infty$ for every
fixed $q>0$. The constant $c_1$ depends only on the fixed quantities,
on $\epsilon_u\epsilon_v$, $F$, the response law, and the constants of
the hypotheses.
\end{proposition}

\begin{proof}
\emph{Scales.} Let $C\ge1$ be the constant of
\cref{lem:poisson-comparison} for the constant $C_0$, let
$C_1=2v_0CC_\Gamma$, and fix $c_0>1+\log C_1+C_E$. We define
\[
 t_n=\frac{\Delta_\theta L}M-\log M+c_0,
\]
we let $B$ be the $d$th power of the smallest even integer $k$ with
$k^dR\ge ne^{t_n}$, and we set $x=BR/n$. Since $M$ is within one of
$M_*=\sqrt{\theta\Delta_\theta L/\tau}$ and
$\Delta_\theta L=(\tau/\theta)M_*^2$, we have
$\Delta_\theta L/M=(\tau/\theta)M+O(1)$, so that
$t_n=(\tau/\theta)M+O(\log M)$. As $\log R=O(\log M)$, the number
$k_*=(ne^{t_n}/R)^{1/d}$ tends to infinity, and $k_*\le k<k_*+2$
yields $1\le xe^{-t_n}\le(1+2/k_*)^d\to1$. In particular,
$x\ge e^{t_n}$, and
\begin{equation}\label{eq:scales-basic}
 \log x=t_n+o(1),\qquad
 \log(B/n)=\frac\tau\theta M+O(\log M),\qquad
 \log\Lambda=-\frac\tau\theta M+O(\log M),
\end{equation}
where we used $\log(B/n)=\log x-\log R$ and $\Lambda=2v_0n/B$. In
particular, $B\ge n$ for all large $n$, and since $M\asymp\sqrt L$, we
have $B/\{n(\log n)^q\}\to\infty$ for every fixed $q>0$. Since
$m\le R^{a_*}$, we have $A_u\le\epsilon_uB^{-a_*}\le n^{-a_*}$, and
similarly $A_v\le\epsilon_vB^{-b_*}\le n^{-b_*}$. Since $|u|\le C_0A_u$
and $|v|\le C_0A_v$, the amplitude assumptions of
\cref{lem:poisson-comparison} hold for all large $n$.

\emph{Indistinguishability.} Let $y=CC_\Gamma\Lambda=C_1n/B$. Since
$\log y=-(\tau/\theta)M+O(\log M)$ by \eqref{eq:scales-basic}, and
$\sqrt M\log R=O(\sqrt M\log M)=o(M)$, we have $yR^{\sqrt M}\le1/2$
for all large $n$. By \eqref{eq:coefficient-hypothesis}, the first term
of \eqref{eq:poisson-comparison-bound} is at most
\[
 CC_\Gamma B\Lambda^2A_u^2A_v^2e^{C_EM}R^{1-M}
 \sum_{j\ge M}\frac{y^jR^{(j-M)\sqrt M}}{j!}
 \le2CC_\Gamma\mathcal I_n,\qquad
 \mathcal I_n=B\Lambda^2A_u^2A_v^2R^{1-M}e^{C_EM}\frac{y^M}{M!},
\]
where we wrote $j=M+i$ and used $M!/(M+i)!\le1$ and
$yR^{\sqrt M}\le1/2$. Since $B\Lambda^2=4v_0^2nR/x$,
$(BR)^{-2\theta}=(nx)^{-2\theta}$, and $R^{1-M}y^M=R(C_1/x)^M$, we have
\[
 \mathcal I_n=4v_0^2(\epsilon_u\epsilon_v)^2n^{\Delta_\theta}R^2m^4
 e^{C_EM}\frac{C_1^M}{M!\,x^{M+1+2\theta}}.
\]
Using $x\ge e^{t_n}$,
$Mt_n=\Delta_\theta L-M\log M+c_0M$, and $\log M!\ge M\log M-M$, we
obtain
\[
 \log\mathcal I_n\le\log(4v_0^2)+2\log R+4\log m-(1+2\theta)t_n
 -(c_0-1-\log C_1-C_E)M.
\]
The terms $2\log R$ and $4\log m$ are $O(\log M)$, since
$m\le c_*M$, and $t_n>0$ for all large $n$. By the choice of $c_0$, we
conclude that $\mathcal I_n\to0$. Let $s_*=\min(a_*,b_*)$. The second
term of \eqref{eq:poisson-comparison-bound}, divided by $\mathcal I_n$,
equals
\[
 C\Lambda^2(A_u^2+A_v^2)^2R^{M-1}e^{-C_EM}C_\Gamma^{-M}
 \le Cn^{-4s_*}e^{O(M\log M)},
\]
since $\Lambda\le2v_0$, $A_u^2+A_v^2\le2n^{-2s_*}$, and
$\log R=O(\log M)$. This tends to zero, since $M\log M=o(L)$. It
follows that
$d_{\rm H}(\mathbb P^{\rm Poi}_{B,+},\mathbb P^{\rm Poi}_{B,-})\to0$.

\emph{Separation.} Let $D_n=|E_+T-E_-T|$, and let
$F^{\rm oo}(u,v)=\tfrac14\{F(u,v)-F(-u,v)-F(u,-v)+F(-u,-v)\}$ be the
part of $F$ that is odd in each argument. By the two symmetries that
follow from (A5), we have
\[
 E_+T-E_-T=E_+\int p\{F(u,v)-F(u,-v)\}=2E_+\int pF^{\rm oo}(u,v),
\]
and in the same way $(E_+-E_-)\int puv=2E_+\int puv$. Moreover,
$4F^{\rm oo}(u,v)=\int_{-u}^u\int_{-v}^vF_{uv}(s,t)\dd t\dd s$ for small
$u,v$. The linear Taylor terms of $F_{uv}$ at $(0,0)$ integrate to zero
over this symmetric rectangle, and since $F\in C^4$ the remainder is
at most $C(s^2+t^2)$. This gives
\[
 F^{\rm oo}(u,v)=F_{uv}(0,0)uv+O\{|uv|(u^2+v^2)\}.
\]
Since $|u|\le C_0A_u$, $|v|\le C_0A_v$, and $\int p=1$, it follows from
\eqref{eq:separation-hypothesis} that
\[
 D_n\ge|F_{uv}(0,0)|c_2A_uA_v\rho_n^M-CA_uA_v(A_u^2+A_v^2).
\]
Since $A_u^2+A_v^2\le2n^{-2s_*}$, while $\rho_n^{-M}=e^{O(\sqrt L)}$, the
second term is at most half the first for all large $n$. Since
$\log\rho_n-\log\rho=O(\delta_n)$ and $M\delta_n\to0$, we also have
$\rho_n^M\ge\rho^M/2$ for all large $n$. With
$f(z)=\theta\Delta_\theta L/z+\tau z$, we have
\[
 f(M)-f(M_*)=\frac{\tau(M-M_*)^2}M=O(M^{-1}),\qquad
 f(M_*)=2\sqrt{\theta\Delta_\theta\tau L}=\kappa\sqrt L,
\]
with $\kappa$ as in \eqref{eq:subcritical-scale}. By
\eqref{eq:scales-basic}, we have $\log(BR)=L+t_n+o(1)$, and
$\theta t_n+\tau M=f(M)-\theta\log M+\theta c_0$. Then
\[
 \log(A_uA_v\rho^M)=\log(\epsilon_u\epsilon_vm^2)-\theta L-f(M)
 +\theta\log M-\theta c_0+o(1)=\log(m^2a_n)+O(1),
\]
where we used $\theta\log M=(\theta/2)\log L+O(1)$. Thus
$D_n\ge c_3m^2a_n$ for a constant $c_3>0$ and all large $n$.

\emph{Testing.} The target $T$ is a constant plus a sum of independent
contributions of the pairs, each bounded by $C/B$, since $F$ is bounded
near the origin. Then $\Var_\pm T\le C/B$, and since $B\ge n$, $m\ge1$,
and $\theta<1/2$, we have
$BD_n^2\ge c_3^2na_n^2=c_3^2n^{1-2\theta}L^\theta e^{-2\kappa\sqrt L}\to\infty$.
Outside $\mathcal Q$ the Poisson process has the same law in every
realization, and it is independent of the process on $\mathcal Q$.
Given the parameters, the number $N$ of points has law
$\operatorname{Poisson}(2n)$, and given $N$ the points are independent
observations from the law $P$. We order the points randomly, keep the first $n$ when $N\ge n$, and
return a fixed sample otherwise. This defines a kernel that does not
depend on the parameters. Its output
mixtures are $(1-\pi_n)\mathbb Q_\pm^{(n)}+\pi_n\delta_*$, where
$\mathbb Q_\pm^{(n)}$ are the mixtures of the laws of $n$ observations,
$\delta_*$ is the law of the fixed sample, and
$\pi_n=P\{\operatorname{Poisson}(2n)<n\}\le e^{-(1-\log2)n}$. Since a
kernel does not increase total variation and $\TV\le d_{\rm H}$, we
obtain
$\TV(\mathbb Q_+^{(n)},\mathbb Q_-^{(n)})
\le d_{\rm H}(\mathbb P^{\rm Poi}_{B,+},\mathbb P^{\rm Poi}_{B,-})/(1-\pi_n)
=o(1)$.

Given an estimator $\widehat T$, we decide for the prior whose mean
$E_\pm T$ is closer to $\widehat T$. Under the $\pm$ prior, the decision
is correct when $|T-E_\pm T|<D_n/4$ and $|\widehat T-T(P)|<D_n/4$, and
by Chebyshev's inequality the first event fails with prior probability
at most $16\Var_\pm T/D_n^2$. Since the two testing errors sum to at
least $1-\TV(\mathbb Q_+^{(n)},\mathbb Q_-^{(n)})$, we obtain
\[
 \sup_{P\in\mathcal P}P\{|\widehat T-T(P)|\ge D_n/4\}
 \ge\frac12\Bigl\{1-\TV(\mathbb Q_+^{(n)},\mathbb Q_-^{(n)})
 -\frac{16(\Var_+T+\Var_-T)}{D_n^2}\Bigr\}=\frac12-o(1).
\]
Since $D_n/4\ge c_1m^2a_n$ with $c_1=c_3/4$, this proves
\eqref{eq:lower-testing-transfer} for all large $n$. This argument is
the total-variation version of the method of two fuzzy
hypotheses~\cite[Theorem~2.15]{tsybakov}, and we control the prior
concentration by Chebyshev's inequality. Since an independent
random seed appended to the observations does not change the total
variation distance, the same argument applies to randomized
estimators.
\end{proof}
% ---------- 06-lattice.tex
\chapter{Lattice priors}\label{sec:lattice}

We construct priors satisfying \cref{prop:reduction}
(\cref{lem:lattice-priors}). In each pair, the density on the bump
supports is $r_0\pm h_n\cos(\Theta+\phi(x))$, where $\Theta$ is uniform
and $\phi$ is a random block function. The perturbations are bumps
divided by the density, multiplied by a gate $S\in[0,1]$, with sign
correlation $\pm\cos(M\Theta)$.
The phase $\phi$ sums independent coefficients $\xi_{a,q}$ times
smoothed binary digits over levels $a=1,\ldots,J$ and coordinates
$q=1,\ldots,d$. For $J=0$, we have $\phi=0$ and $S=1$, and the bound is
$c\,a_n$; for $J\ge1$, three properties gain the factor
$\log n$.
\begin{enumerate}[label=(\roman*)]
\item Outside a set of small measure, $\phi\in(2\pi/M)\mathbb Z$,
so the $M$th Fourier coefficient of $1/p$ in $\Theta$ retains the
separation factor $\rho_n^M$ of block-constant phases.
\item The gated characteristic functions of $\xi_{a,q}$ vanish outside
$\bigcup_{l\in\mathbb Z}[lM-K_a,lM+K_a]$. For a likelihood term with
one factor each of $u,v$ and $j$ density factors, averaging over
$\Theta$ gives zero for $j<M$. For $M\le j<M+\sqrt M$, the band
restriction implies the digit mismatch bounds \eqref{lat:mismatch},
whose integrals yield $R^{1-M}e^{C_EM}$, where $R$ counts fine cells.
For larger $j$, a crude bound suffices.
\item The gate controls weighted coefficient powers: $S$ times the
change of $\phi$ across a fine cell is $O(1/m)$, with $m\asymp M$.
Dividing a bump by the density costs a factor $1/m$ of a fine-scale
derivative, allowing amplitudes $m(BR)^{-a_*}$ and $m(BR)^{-b_*}$.
\end{enumerate}

\section{Lattice laws with compact Fourier support}\label{lat:sec-gate}

Throughout this section we fix the integer
$Q=\lceil\max(\alpha,\beta)\rceil+2$, and we define
\[
 \Pi(y)=c_Q^{-1}\Bigl(\frac{\sin y}{y}\Bigr)^{2Q},\qquad
 c_Q=\int_{\mathbb R}\Bigl(\frac{\sin y}{y}\Bigr)^{2Q}\dd y,
\]
where $\sin y/y$ takes its continuous value one at $y=0$. For
$Y\sim\Pi$ we write $\mu_2=EY^2$, which is finite because
$\Pi(y)\le c_Q^{-1}\min(1,|y|^{-2Q})$, and we write
$\widehat f(w)=\int_{\mathbb R}f(x)e^{iwx}\dd x$.

\begin{lemma}[Gated lattice laws]\label{lat:gate}
Let $\eta>0$ and $\lambda\ge1$, let $M\ge2$ be an integer, and let
$h_M=2\pi/M$, $\pi_\eta(x)=\eta^{-1}\Pi(x/\eta)$, and $K=6Q/\eta$. We
define
\[
 G(x)=\Bigl(\frac{\sin(x/(\lambda\eta))}{x/(\lambda\eta)}\Bigr)^{2Q}
 \qquad(x\in\mathbb R),\qquad G(0)=1.
\]
If $4K\le M$, then the numbers $h_M\pi_\eta(x)$, $x\in h_M\mathbb Z$,
sum to one. If $\xi$ has the law $P\{\xi=x\}=h_M\pi_\eta(x)$ on
$h_M\mathbb Z$, then
\begin{equation}\label{lat:gate-bounds}
 0\le G\le1,\qquad
 1-EG(\xi)\le\frac{Q\mu_2}{3\lambda^2},\qquad
 G(x)|x|^r\le(\lambda\eta)^r
 \quad(x\in\mathbb R,\ 0\le r\le2Q-2).
\end{equation}
Moreover, the function $\widehat\Xi(w)=E[G(\xi)^2e^{iw\xi}]$ of
$w\in\mathbb R$ satisfies $|\widehat\Xi|\le1$, and $\widehat\Xi(w)=0$
unless $|w-lM|\le K$ for some $l\in\mathbb Z$.
\end{lemma}

\begin{proof}
The identity $\sin y/y=\frac12\int_{-1}^1e^{i\omega y}\dd\omega$ shows
that $(\sin y/y)^{2Q}$ is the Fourier transform of the $2Q$-fold
convolution of $\frac12\mathbf 1_{[-1,1]}$, which is supported in
$[-2Q,2Q]$. By Fourier inversion, $\widehat\pi_\eta$ and $\widehat G$
are supported in $[-2Q/\eta,2Q/\eta]$ and
$[-2Q/(\lambda\eta),2Q/(\lambda\eta)]$, respectively. Since
multiplication convolves Fourier transforms and adds their supports, and
since $\lambda\ge1$, $\widehat{G^k\pi_\eta}$ is supported in $[-K,K]$ for
$k=0,1,2$.

Let $f=G^k\pi_\eta$ with $k\in\{0,1,2\}$, and let $w\in\mathbb R$.
Since $0\le f(x)\le C_\eta(1+|x|)^{-2Q}$, the $h_M$-periodic function
$y\mapsto\sum_{x\in h_M\mathbb Z}f(y+x)e^{iw(y+x)}$ is continuous, and
its Fourier coefficients are $h_M^{-1}\widehat f(w-lM)$,
$l\in\mathbb Z$. Since only finitely many of them are nonzero, the
function is a trigonometric polynomial and equals its Fourier series
everywhere. At $y=0$ we obtain an instance of the Poisson summation
formula~\cite[Chapter~5, Section~3]{steinshakarchi},
\begin{equation}\label{lat:sampling}
 h_M\sum_{x\in h_M\mathbb Z}f(x)e^{iwx}
       =\sum_{l\in\mathbb Z}\widehat f(w+lM).
\end{equation}
At $w=0$ only the term $l=0$ remains, since $K<M$. With $k=0$, the
numbers $h_M\pi_\eta(x)$ sum to $\widehat\pi_\eta(0)=1$, and with $k=1$
we obtain $EG(\xi)=\int G\pi_\eta$. With $k=2$, \eqref{lat:sampling}
gives $\widehat\Xi(w)=\sum_l\widehat f(w+lM)$. Since $\widehat f$
vanishes outside $[-K,K]$ and $2K<M$, at most one term is nonzero, and
it vanishes unless $|w+lM|\le K$; moreover
$|\widehat\Xi|\le EG(\xi)^2\le1$.

The bound $|\sin z/z|\le\min(1,|z|^{-1})$ implies that $0\le G\le1$ and
$G(x)|x|^r\le(\lambda\eta)^r$ for $0\le r\le2Q$. For independent $U,V$
uniform on $[-1,1]$, we have
$(\sin z/z)^2=E\cos(z(U+V))\ge1-z^2E(U+V)^2/2=1-z^2/3$. Then
$1-(\sin z/z)^{2Q}\le Q\{1-(\sin z/z)^2\}\le Qz^2/3$. Since
$\int G\pi_\eta=E(\sin(Y/\lambda)/(Y/\lambda))^{2Q}$ for $Y\sim\Pi$, we
obtain $1-EG(\xi)\le Q\mu_2/(3\lambda^2)$.
\end{proof}

Lattice laws obtained in this way, by sampling a probability density
whose Fourier transform has compact support, are constructed
in~\cite[Section~V, Proposition~5 and Corollary~6]{vatedka2015}. There,
the sampled masses are normalized by the same formula, the
characteristic function is the sum of the translated copies of the
Fourier transform, and Example~1 uses powers of $\sin y/y$.

\section{Construction of the priors}\label{lat:sec-construction}
\label{sec:block-phases}

We use the cube $\mathcal Q$ of volume $v_0$, its blocks and pairs, and
the block coordinates $x\in\mathcal U=[0,1]^d$ of
\cref{sec:lower-preparation}. We fix smooth functions
$\iota_{\rm in},\iota_{\rm out}:\mathcal U\to[0,1]$ with compact
supports in the interior of $\mathcal U$, such that $\iota_{\rm in}$ is
not identically zero and $\iota_{\rm out}=1$ on the support of
$\iota_{\rm in}$. We extend them by zero to $\mathbb R^d$, and we write
$\bar\iota=\int_{\mathcal U}\iota_{\rm out}$ and
$I_{\rm in}=\int_{\mathcal U}\iota_{\rm in}^2>0$. We set
\[
 r_0=\frac{r_-+r_+}2,\qquad h_n=\frac{r_+-r_-}2-\delta_n,\qquad
 p_0=\frac{1-v_0\bar\iota r_0}{1-v_0\bar\iota},
\]
and we choose $v_0$ so small that $p_0$ lies in a fixed interior
subinterval of $(r_-,r_+)$. Throughout this section, $n$ is large
enough that
\[
 0<\delta_n\le\tfrac12\min\{r_0-r_-,\ p_0-r_-,\ r_+-p_0\}.
\]
We write $s_0=\min(\alpha,\beta)$ and $\alpha_0=\min(s_0,1)/2$; we use
only $0<\alpha_0<1$ and $\alpha_0<s_0$. Let $J\ge0$ be an integer, and
let $R=2^{Jd}$ and $m=2^{Js_0}$. The fine cells of a block are its $R$
dyadic subcubes of side $2^{-J}$. Sums over an empty set of levels are
zero, and products over it are one.

\emph{Digit selectors.} For $y\in\mathbb R$ and $a\ge1$ we write
$e_a(y)=\lfloor2^ay\rfloor\bmod2$. If $x$ is uniform on $\mathcal U$,
the digits $e_a(x_q)$, $1\le a\le J$, $1\le q\le d$, are independent and
uniform on $\{0,1\}$, and the string of these $dJ$ digits determines the
fine cell of $x$. To smooth the digits, we fix a nondecreasing
$C^\infty$ function $\Upsilon$ with $\Upsilon=0$ on $(-\infty,-1]$,
$\Upsilon=1$ on $[1,\infty)$, and $\Upsilon(0)=1/2$. For
$0<\gamma\le1/4$ and $y\in\mathbb R$ we let $z$ be a nearest integer to
$y$, and we define
\[
 \varpi_\gamma(y)=\Upsilon\bigl((-1)^{z+1}(y-z)/\gamma\bigr).
\]
At distance at least $\gamma$ from the integers, $\varpi_\gamma(y)$
equals $\lfloor y\rfloor\bmod2$; in particular it is locally constant
near the half-integers, where the two choices of $z$ agree. This
implies that $\varpi_\gamma$ is $C^\infty$, with values in $[0,1]$ and
$|\varpi_\gamma^{(N)}|\le C_N\gamma^{-N}$ for $N\ge1$. If
$\lfloor y\rfloor$ is odd, then either $z=\lfloor y\rfloor$ and
$y-z\ge0$, or $z=\lfloor y\rfloor+1$ is even and $z-y>0$; in both cases
$\varpi_\gamma(y)\ge\Upsilon(0)=1/2$. In the same way
$\varpi_\gamma(y)\le1/2$ if $\lfloor y\rfloor$ is even. Therefore
\begin{equation}\label{lat:soft-digit}
 |\varpi_\gamma(y)-e|\ge\tfrac12\qquad
 \text{if }e\in\{0,1\}\text{ and }e\ne\lfloor y\rfloor\bmod2 .
\end{equation}

\emph{Scales.} For a level $1\le a\le J$, we write $b=J-a$ and define
\[
 \gamma_a=\frac{\gamma_*}{(b+1)^2},\qquad
 \lambda_a=\lambda_*(b+1),\qquad
 \eta_a=\frac{\gamma_a2^{b\alpha_0}}{\lambda_am},\qquad
 K_a=\frac{6Q}{\eta_a},
\]
where $\gamma_*=\min\{1/4,I_{\rm in}/(16d)\}$ and
$\lambda_*=\max\{1,\sqrt{8Qd\mu_2/3}\}$. Thus
$\lambda_a\eta_a=\gamma_a2^{b\alpha_0}/m$, so that coarse levels carry
larger coefficients, and $K_a=K(b)m$ with
$K(b)=6Q\lambda_*(b+1)^32^{-b\alpha_0}/\gamma_*$. We set
$\overline K=\sup_{b\ge0}K(b)$ and $c_*=1/(4\overline K)$, and we
require that $m\le c_*M$, so that
\begin{equation}\label{lat:band}
 4K_a\le M\qquad(1\le a\le J).
\end{equation}
All these constants depend only on $d,\alpha,\beta$ and
$\iota_{\rm in}$, and they are fixed before $M$ is chosen.

\emph{The phase and the gate.} Let $M\ge2$ be even with $m\le c_*M$,
and let $h_M=2\pi/M$. In each pair, let $\xi_{a,q}$, for $1\le a\le J$
and $1\le q\le d$, be independent, with the law of $\xi$ in
\cref{lat:gate} for $\eta=\eta_a$ and $\lambda=\lambda_a$; by
\eqref{lat:band}, the lemma applies. We write $G_a$ and $\widehat\Xi_a$
for the corresponding functions $G$ and $\widehat\Xi$, and we define
\[
 S=\prod_{a,q}G_a(\xi_{a,q}),\qquad
 s_{a,q}(x)=\varpi_{\gamma_a}(2^ax_q),\qquad
 \phi(x)=\sum_{a=1}^J\sum_{q=1}^d\xi_{a,q}s_{a,q}(x)
\]
for $x\in\mathbb R^d$. By independence, \eqref{lat:gate-bounds}, and the
inequality $\prod_i(1-t_i)\ge1-\sum_it_i$ for $t_i\in[0,1]$, we have
\[
 ES\ge1-\frac{Qd\mu_2}{3\lambda_*^2}\sum_{b\ge0}(b+1)^{-2}\ge\frac34,
 \qquad ES^2\ge(ES)^2\ge\frac9{16},
\]
where we used $\sum_{b\ge0}(b+1)^{-2}<2$ and the choice of $\lambda_*$.
Since $S\le G_a(\xi_{a,q})$ for every $(a,q)$ and $0\le G_a\le1$, the
last bound in \eqref{lat:gate-bounds} implies that, for distinct pairs
$(a_i,q_i)$ and integers $r_i\ge1$ with $\sum_ir_i\le2Q-2$,
\begin{equation}\label{lat:gate-product}
 S\prod_i|\xi_{a_i,q_i}|^{r_i}
 \le\prod_i(\lambda_{a_i}\eta_{a_i})^{r_i}.
\end{equation}

Let $\mathcal U^\circ$ be the set of $x\in\mathcal U$ such that
$2^ax_q$ has distance at least $\gamma_a$ from the integers for all $a$
and $q$. On $\mathcal U^\circ$ we have $s_{a,q}(x)=e_a(x_q)\in\{0,1\}$,
so that $\phi(x)$ is a sum of elements of $h_M\mathbb Z$, and therefore
\begin{equation}\label{lat:coherence}
 \cos(M\phi(x))=1\qquad(x\in\mathcal U^\circ).
\end{equation}
For each $a$ and $q$, the set of $x\in\mathcal U$ for which $2^ax_q$ has
distance less than $\gamma_a$ from the integers has measure
$2\gamma_a$. Since $\sum_a\gamma_a<2\gamma_*$, a union bound gives
\begin{equation}\label{lat:collar}
 |\mathcal U\setminus\mathcal U^\circ|<4d\gamma_*\le I_{\rm in}/4 .
\end{equation}

\emph{The priors.} In each pair, let $\Theta$ be uniform on $[0,2\pi)$
and independent of $\xi=(\xi_{a,q})_{a,q}$. We call $(\Theta,\xi)$ the
density parameters of the pair; both blocks of the pair use them, and
their law is the same under both priors. On the block of the pair with
label $\ell$, in block coordinates, we define
\[
 r(x)=r_0+(-1)^{\ell-1}h_n\cos(\Theta+\phi(x)),\qquad
 p=p_0+\iota_{\rm out}(r-p_0),
\]
and $p=p_0$ outside $\mathcal Q$. Every realization of $\phi$ is
$C^\infty$, since it is a finite sum with finite coefficients. Since
both blocks of a pair use the same function $\phi$ in block coordinates,
the oscillations of the two blocks cancel in $\int p$, and the mass of
$p$ on each pair is $2(v_0/B)\bar p$ in every realization, where
\[
 \bar p=p_0+\bar\iota(r_0-p_0)
\]
is the mean value of $p$ over a pair. The total mass is therefore
$p_0(1-v_0)+v_0\bar p=p_0+v_0\bar\iota(r_0-p_0)=1$. Since
$\iota_{\rm out}$ vanishes near the block boundaries, $p$ is a
$C^\infty$ probability density. On the support of $\iota_{\rm out}$ it
is a convex combination of $p_0$ and $r\in[r_-+\delta_n,r_+-\delta_n]$,
so that $r_-+\delta_n\le p\le r_+-\delta_n$.

Conditional on the density parameters, the signs
$\sigma_u,\sigma_v\in\{-1,1\}$ of a pair have the law
\[
 P_\pm\{\sigma_u=s,\sigma_v=t\mid\Theta,\xi\}
 =\tfrac14\{1\pm st\cos(M\Theta)\}\qquad(s,t\in\{-1,1\})
\]
under the $\pm$ prior, as in (A5) with $H=\xi$. As in
\cref{sec:lower-preparation}, $E$ averages the density parameters
of one pair. Let
$\epsilon_u,\epsilon_v\in(0,1]$. With the amplitudes
\begin{equation}\label{eq:lattice-amplitudes}
 A_u=\epsilon_u(BR)^{-a_*}m,\qquad A_v=\epsilon_v(BR)^{-b_*}m,
\end{equation}
we define
\begin{equation}\label{lat:profiles}
 u=A_u\sigma_uS\iota_{\rm in}/p,\qquad
 v=A_v\sigma_vS\iota_{\rm in}/p
\end{equation}
on both blocks of each pair, and $u=v=0$ outside $\mathcal Q$. The
parameters of distinct pairs are independent. Without the gate $S$ and
the division by $p$, these are smooth bumps on the blocks with random
signs. Classical lower-bound constructions use such localized bumps
with binary coefficients~\cite[Section~2, p.~1045]{stone} or random
signs~\cite[Section~3.5.2]{ingster},~\cite[Sections~3--4]{robins2009};
here the signs are correlated with the phase of the density.

\begin{lemma}[Lattice priors]\label{lem:lattice-priors}
\label{lat:smoothness}\label{lat:coefficients}\label{lat:separation}
Let $J\ge0$, let $M\ge2$ be even with $m\le c_*M$, and let $B$ be the
$d$th power of an even integer. There are constants $C,C_\Gamma\ge1$
and $C_E\ge0$, which depend only on $d,\alpha,\beta,r_\pm$ and the
fixed choices of the construction, and not on
$n,B,J,M,\epsilon_u,\epsilon_v$, such that the two priors above have
the following properties.
\begin{enumerate}[label=(\alph*)]
\item They are admissible with frequency $M$, with $C_0=1/r_-$ and the
amplitudes \eqref{eq:lattice-amplitudes}. Every realization of $p$ is
a $C^\infty$ probability density with
\begin{equation}\label{eq:lower-density-margin}
 r_-+\delta_n\le p\le r_+-\delta_n,
\end{equation}
and the functions $u$ and $v$ vanish outside $\mathcal Q$.
\item In every realization, $u$ and $v$ are $C^\infty$, and
\begin{equation}\label{eq:lower-norms}
 \|u\|_{C^\alpha}\le C\epsilon_u,\qquad
 \|v\|_{C^\beta}\le C\epsilon_v,\qquad
 \|u\|_\infty\le A_u/r_-,\qquad \|v\|_\infty\le A_v/r_-.
\end{equation}
Let $K$ be smooth on a neighborhood $N$ of $(0,0)$. If $A_u+A_v$ is
below a threshold that depends on $K$, then
\begin{equation}\label{eq:lower-smooth-map}
\begin{aligned}
 K(0,v)=0\text{ on }N&\ \Longrightarrow\
       \|K(u,v)\|_{C^\alpha}\le C_K\epsilon_u,\\
 K(u,0)=0\text{ on }N&\ \Longrightarrow\
       \|K(u,v)\|_{C^\beta}\le C_K\epsilon_v,
\end{aligned}
\end{equation}
where $C_K$ depends only on $K$ and on the quantities on which $C$
depends.
\item For every pair and every $j\ge0$, the coefficients $\Gamma_j$ of
\cref{lem:poisson-comparison} satisfy
\eqref{eq:coefficient-hypothesis}, that is,
\[
 \|\Gamma_j\|_2^2\le C_\Gamma^{j+1}A_u^2A_v^2
 e^{C_EM}R^{1-M+(j-M)\sqrt M}\mathbf 1_{j\ge M},
\]
where the norm is that of $L^2((\mathcal U\times\{1,2\})^{j+2})$.
\item With $\rho_n$ as in \cref{prop:reduction}, we have
\[
 (E_+-E_-)\int_{[0,1]^d}puv\ge c_2A_uA_v\rho_n^M,
 \qquad c_2=\frac{9v_0I_{\rm in}}{16r_0}.
\]
\end{enumerate}
\end{lemma}

\emph{Block-constant phases.} For $J=0$, we have $\phi=0$, $S=1$,
$R=m=1$, and $\mathcal U^\circ=\mathcal U$, with $\Theta$ the only
density parameter. On the bump supports, $p=r_0\pm h_n\cos\Theta$,
$pu=A_u\sigma_u\iota_{\rm in}$, and $pv=A_v\sigma_v\iota_{\rm in}$. After
removing their signs, these products do not depend on $\Theta$. The $M$th
Fourier coefficient of $1/p$ is of order $\rho_n^M$
by \eqref{eq:cosine-reciprocal}. The scale conditions of
\cref{prop:reduction} reduce to $1\le c_*M$, which holds
for all large $n$. That proposition, applied with
\cref{lem:lattice-priors}(a), (c), and~(d), yields, for every
class $\mathcal P$ and target $T$ as in the proposition, a constant
$c_1>0$ such that, for all large $n$,
\begin{equation}\label{eq:lattice-testing}
 \inf_{\widehat T}\sup_{P\in\mathcal P}
 P\{|\widehat T-T(P)|\ge c_1m^2a_n\}\ge\frac38
\end{equation}
with $m=1$. The proof of \cref{prop:local-lower} uses the
largest $J$ with $m\le c_*M$, for which $m^2\asymp\log n$.

\section{Proof of \texorpdfstring{\cref{lem:lattice-priors}}{the lemma on lattice priors}}
\label{lat:sec-properties}\label{lat:sec-proof}

We use the H\"older norm $\|\cdot\|_{C^t}$ of the model class on
$[0,1]^d$, and we write $\|\cdot\|_{C^t(\mathcal U)}$ for the norm
defined in the same way on $\mathcal U$, in block coordinates.

\begin{proof}[Proof of \cref{lem:lattice-priors}(a)]
We showed above that $p$ is a $C^\infty$ probability density that
satisfies \eqref{eq:lower-density-margin}, equals $p_0$ outside
$\mathcal Q$, and has the same mass on every pair in every
realization, which is (A2). By construction, the parameters of
distinct pairs are independent, $u$ and $v$ vanish outside
$\mathcal Q$, and the signs have the law in (A5). Since $p\ge r_-$ and
$0\le S,\iota_{\rm in}\le1$, we have $|u|\le A_u/r_-$ and
$|v|\le A_v/r_-$, which gives (A3) with $C_0=1/r_-$. In (A4), we have
$H=\xi$, $pu=A_u\sigma_uS\iota_{\rm in}$,
$pv=A_v\sigma_vS\iota_{\rm in}$, and, on the block with label $\ell$,
\[
 p_{\rm c}=p_0+\iota_{\rm out}(r_0-p_0),\qquad
 p_1=(-1)^{\ell-1}h_n\iota_{\rm out}\cos\phi,\qquad
 p_2=-(-1)^{\ell-1}h_n\iota_{\rm out}\sin\phi,
\]
and $h_u=h_v=S\iota_{\rm in}$. These functions depend on $\xi$ but
neither on $\Theta$ nor on the prior.
\end{proof}

\begin{proof}[Proof of \cref{lem:lattice-priors}(b)]
We work in block coordinates on a block with label $\ell$, and we set
$k_{\max}=\lceil\max(\alpha,\beta)\rceil=Q-2$ and
\[
 f(y)=\frac1{r_0+(-1)^{\ell-1}h_n\cos(\Theta+y)},\qquad
 \mathsf H=S\iota_{\rm in}f(\phi).
\]
Since $p=r$ on the support of $\iota_{\rm in}$, we have
$\mathsf H=S\iota_{\rm in}/p$ and $|\mathsf H|\le1/r_-$. The
derivatives of $f$ through order $k_{\max}+1$ are bounded by constants
that depend only on $r_\pm$ and $k_{\max}$, since its denominator lies in
$[r_-,r_+]$. Below, constants may depend on
$d,\alpha,\beta,r_\pm,\iota_{\rm in},\Upsilon$ and the fixed constants
of the construction, but not on $n,B,J,M,\epsilon_u,\epsilon_v$, or the
realization.

\emph{Compositions.} Let $\kappa$ be a $C^\infty$ function on
$[-1/r_-,1/r_-]$ with $\kappa(0)=0$, and let
$s_0\le t\le\max(\alpha,\beta)$. We claim that
\begin{equation}\label{lat:composition}
 \|\kappa(\mathsf H)\|_{C^t(\mathcal U)}\le C_\kappa(1+2^{Jt}/m),
\end{equation}
where $C_\kappa$ depends on $\kappa$ only through bounds on its
derivatives of order at most $\lceil t\rceil+1$. We define
\[
 w_a=\lambda_a\eta_a=\frac{\gamma_a2^{(J-a)\alpha_0}}m,\qquad
 L_a=\frac{2^a}{\gamma_a},\qquad
 \phi_a=\sum_{c=1}^a\sum_{q=1}^d\xi_{c,q}s_{c,q}\quad(0\le a\le J).
\]
Since $\gamma_a\le1/4$ and $2^{(J-a)\alpha_0}\le2^{Js_0}=m$, we have
$w_a\le1$ and $L_a\ge1$. Since $s_{a,q}$ depends only on $x_q$, we
also have $|D^{\mathbf j}s_{a,q}|\le CL_a^{|\mathbf j|}$ for every
multi-index $\mathbf j$, including $\mathbf j=0$. For an integer
$1\le k\le\lceil t\rceil$ and $b=J-a$, we have
$w_{a-i}L_{a-i}^k/(w_aL_a^k)=\{(b+i+1)/(b+1)\}^{2(k-1)}2^{-i(k-\alpha_0)}$
for $0\le i<a$, and hence
\begin{equation}\label{lat:earlier-levels}
 \sum_{c=1}^a\sum_{q=1}^dw_cL_c^k
 \le dw_aL_a^k\sum_{i\ge0}(i+1)^{2(k-1)}2^{-i(k-\alpha_0)}
 \le Cw_aL_a^k,
\end{equation}
where the series converges because $k\ge1>\alpha_0$.

Let $F_a=\kappa(S\iota_{\rm in}f(\phi_a))$, $V_a=F_a-F_{a-1}$, and
$W_a=\phi_a-\phi_{a-1}$ for $1\le a\le J$. By the fundamental theorem
of calculus,
\[
 V_a=S\iota_{\rm in}W_a\int_0^1
 f'(z_{a,u})\kappa'\bigl(S\iota_{\rm in}f(z_{a,u})\bigr)\dd u,
 \qquad z_{a,u}=\phi_{a-1}+uW_a,
\]
and the arguments of $\kappa$ lie in $[0,1/r_-]$. Let
$0\le|\mathbf j|\le\lceil t\rceil$. By the chain and product rules,
$D^{\mathbf j}V_a$ is a sum, with a number of terms that depends only
on $|\mathbf j|$ and $d$, of products of a power $S^\varrho$ with
$\varrho\ge1$, of derivatives of $\iota_{\rm in}$, $f$, and $\kappa$ of
order at most $|\mathbf j|+1$, of $D^{\mathbf j_0}W_a$, and of
$D^{\mathbf j_l}z_{a,u}$ for $1\le l\le r$, where $|\mathbf j_l|\ge1$
for $l\ge1$ and $|\mathbf j_0|+\sum_l|\mathbf j_l|\le|\mathbf j|$. In
particular $r\le|\mathbf j|$. We expand $W_a$ and $z_{a,u}$ over levels
and coordinates. The coefficients then enter through products
\[
 \xi_{a,q_0}D^{\mathbf j_0}s_{a,q_0}
 \prod_{l=1}^{r}\xi_{c_l,q_l}D^{\mathbf j_l}s_{c_l,q_l}
 \qquad(c_l\le a),
\]
with at most $|\mathbf j|+1\le k_{\max}+1=Q-1\le2Q-2$ coefficients,
counted with multiplicity. Since $S^\varrho\le S$, we group repeated
coefficients and apply \eqref{lat:gate-product}, which bounds $S^\varrho$
times such a product by
$Cw_aL_a^{|\mathbf j_0|}\prod_lw_{c_l}L_{c_l}^{|\mathbf j_l|}$.
Summing over $q_0$, and over the levels $c_l\le a$ and the coordinates
$q_l$ by \eqref{lat:earlier-levels}, and using $w_a\le1$ and
$L_a\ge1$, we obtain
\[
 \|D^{\mathbf j}V_a\|_\infty\le C_\kappa w_aL_a^{|\mathbf j|}
 \qquad(0\le|\mathbf j|\le\lceil t\rceil).
\]
With $N=\lceil t\rceil-1$ and $\vartheta=t-N\in(0,1]$, we find from
the mean value theorem and this bound that, for $|\mathbf j|=N$,
\[
 |D^{\mathbf j}V_a(x)-D^{\mathbf j}V_a(y)|
 \le C_\kappa w_aL_a^N\min\{1,L_a|x-y|\}
 \le C_\kappa w_aL_a^t|x-y|^\vartheta,
\]
where we used $\min(1,z)\le z^\vartheta$ for $z\ge0$. Since $L_a\ge1$,
the derivative suprema in the norm are also at most
$C_\kappa w_aL_a^t$, and therefore
$\|V_a\|_{C^t(\mathcal U)}\le C_\kappa w_aL_a^t$. The function
$F_0=\kappa(S\iota_{\rm in}f(0))$ has norm at most $C_\kappa$, since
$Sf(0)\in[0,1/r_-]$ is a number. Since
$\kappa(\mathsf H)=F_J=F_0+\sum_{a=1}^JV_a$, we obtain
\[
 \|\kappa(\mathsf H)\|_{C^t(\mathcal U)}
 \le C_\kappa\Bigl\{1+\frac{2^{Jt}}m\gamma_*^{1-t}
 \sum_{b\ge0}(b+1)^{2(t-1)}2^{-b(t-\alpha_0)}\Bigr\}
 \le C_\kappa(1+2^{Jt}/m),
\]
where the series is bounded uniformly in $s_0\le t\le\max(\alpha,\beta)$
because $s_0>\alpha_0$. For $J=0$ the sum over $a$ is empty. This
proves \eqref{lat:composition}.

\emph{Conclusion.} On a block, we have $u=A_u\sigma_u\mathsf H$ and
$v=A_v\sigma_v\mathsf H$ in block coordinates, which gives the
sup-norm bounds. Let $g$ be a $C^\infty$ function on $[0,1]^d$ that
vanishes outside $\mathcal Q$ and, on each block, equals $g_k$ in block
coordinates, where every $g_k$ vanishes near $\partial\mathcal U$. The
blocks have side $\ell_B=(v_0/B)^{1/d}\le1$. Rescaling multiplies a
derivative of order $|\mathbf j|\le t$ by $\ell_B^{-|\mathbf j|}\le\ell_B^{-t}$
and a seminorm of order $\vartheta$ of a derivative of order
$\lceil t\rceil-1$ by $\ell_B^{-t}$. If $x$ and $y$ do not lie in the
same block, the segment from $x$ to $y$ meets the boundary of the block
of $x$ at a point $x'$ and the boundary of the block of $y$ at a point
$y'$, where all derivatives of $g$ vanish; the same holds if one of the
points lies outside $\mathcal Q$. It follows that a difference quotient
of order
$\vartheta\le1$ between $x$ and $y$ is at most twice the largest such
quotient within one block. Since the norm adds the derivative suprema
and the H\"older seminorm, whose maxima may occur in different blocks,
we obtain
\begin{equation}\label{lat:rescale}
 \|g\|_{C^t([0,1]^d)}\le3(B/v_0)^{t/d}\max_k\|g_k\|_{C^t(\mathcal U)}.
\end{equation}
We apply \eqref{lat:composition} with $\kappa(z)=z$ and $t=\alpha$, and
\eqref{lat:rescale}. Since $2^{J\alpha}=R^{a_*}$,
$A_uB^{a_*}=\epsilon_uR^{-a_*}m$, and $m=2^{Js_0}\le R^{a_*}$, we obtain
\[
 \|u\|_{C^\alpha}\le CA_uB^{a_*}(1+R^{a_*}/m)
 =C\epsilon_u(mR^{-a_*}+1)\le2C\epsilon_u .
\]
The bound for $v$ is the same with $\beta$, $b_*$, and $\epsilon_v$.

Let $K$ be smooth on a neighborhood of $(0,0)$ with $K(0,v)=0$ there,
and fix a closed ball $N'$ of radius at most one around $(0,0)$ inside
this neighborhood. On $N'$ we write $K(u,v)=u\widetilde K(u,v)$ with
$\widetilde K(u,v)=\int_0^1\partial_uK(su,v)\dd s$. If
$(A_u+A_v)/r_-$ is at most the radius of $N'$, then on each block
$K(u,v)=A_u\sigma_u\kappa(\mathsf H)$ with
$\kappa(z)=z\widetilde K(A_u\sigma_uz,A_v\sigma_vz)$, and $K(u,v)=0$
outside the blocks. The derivatives of $\kappa$ of order at most
$k_{\max}+1$ on $[-1/r_-,1/r_-]$ are bounded by a constant that depends
only on $K$ and $N'$, because $A_u\le\epsilon_uB^{-a_*}\le1$ and
$A_v\le\epsilon_vB^{-b_*}\le1$. The argument above, with this
$\kappa$ and $t=\alpha$, yields
$\|K(u,v)\|_{C^\alpha}\le C_K\epsilon_u$. The case $K(u,0)=0$ is
symmetric, with $\beta$ and $\epsilon_v$.
\end{proof}

\begin{proof}[Proof of \cref{lem:lattice-priors}(c)]
We fix a pair. Since $pu=A_u\sigma_uS\iota_{\rm in}$ and
$pv=A_v\sigma_vS\iota_{\rm in}$, \eqref{lo:sign-difference} gives
\[
 \Gamma_j(x,y,\zeta_1,\ldots,\zeta_j)
 =2A_uA_v\iota_{\rm in}(x)\iota_{\rm in}(y)
 E\Bigl[S^2\cos(M\Theta)\prod_{i=1}^j\{p(\zeta_i)-\bar p\}\Bigr].
\]
Since $p$ and $\bar p$ lie in $[r_-,r_+]$, we have $|p-\bar p|\le r_+$,
so that $\|\Gamma_j\|_\infty\le2A_uA_vr_+^j$. Since
$\|\Gamma_j\|_2^2\le\|\Gamma_j\|_\infty\|\Gamma_j\|_1$, it suffices to
prove that
\begin{equation}\label{lat:L1-bound}
 \|\Gamma_j\|_1\le8A_uA_v(6\mathrm er_+)^je^{C_EM}
 R^{1-M+(j-M)\sqrt M}\mathbf 1_{j\ge M},
\end{equation}
where $\mathrm e=\exp(1)$; then (c) holds with
$C_\Gamma=16(1+6\mathrm er_+^2)$. As $(\mathcal U\times\{1,2\})^{j+2}$
has measure $2^{j+2}$, we have
\begin{equation}\label{lat:crude}
 \|\Gamma_j\|_1\le8A_uA_v(2r_+)^j .
\end{equation}
If $j\ge M+\sqrt M$, then $(j-M)\sqrt M\ge M$, so that
\eqref{lat:crude} implies \eqref{lat:L1-bound}.

We expand each factor in Fourier modes in $\Theta$. On the block with
label $\ell$,
\[
 \begin{gathered}
 p(\zeta)-\bar p=c(\zeta)+w(\zeta)\sum_{o=\pm1}e^{io(\Theta+\phi(\zeta))},\\
 c=(\iota_{\rm out}-\bar\iota)(r_0-p_0),\quad
 w=(-1)^{\ell-1}\iota_{\rm out}h_n/2,
 \end{gathered}
\]
and $|c|\le r_+$ and $|w|\le r_+/2$. We write
$\cos(M\Theta)=\frac12\sum_{\varepsilon=\pm1}e^{i\varepsilon M\Theta}$,
and we fix $\varepsilon$ and a mode $o_i\in\{-1,0,1\}$ for each
$i\in\{1,\ldots,j\}$. We denote the corresponding term of $\Gamma_j$ by
$F_{\varepsilon,o}$; it is
$A_uA_v\iota_{\rm in}(x)\iota_{\rm in}(y)
\prod_{o_i=0}c(\zeta_i)\prod_{o_i\ne0}w(\zeta_i)$ times the expectation
of $S^2e^{i\varepsilon M\Theta}\prod_{o_i\ne0}e^{io_i(\Theta+\phi(\zeta_i))}$.
Since
$S^2e^{i\sum_io_i\phi(\zeta_i)}=\prod_{a,q}G_a(\xi_{a,q})^2e^{i\xi_{a,q}\Phi_{a,q}}$,
integrating over the uniform angle and the independent coefficients,
we find that this expectation equals
\begin{equation}\label{lat:alias-factor}
 \mathbf 1\Bigl\{\sum_io_i=-\varepsilon M\Bigr\}
 \prod_{a,q}\widehat\Xi_a(\Phi_{a,q}),\qquad
 \Phi_{a,q}=\sum_{i:o_i\ne0}o_is_{a,q}(\zeta_i).
\end{equation}
Since $|\sum_io_i|\le j$, it follows that $\Gamma_j=0$ for $j<M$.

Now let $M\le j<M+\sqrt M$. Since there are $2\cdot3^j$ terms
$F_{\varepsilon,o}$, \eqref{lat:L1-bound} follows from the bound
\begin{equation}\label{lat:pattern-bound}
 \|F_{\varepsilon,o}\|_1\le4A_uA_v(2\mathrm er_+)^jR^{1-M}e^{C_EM},
\end{equation}
which we prove for a fixed pattern with $\sum_io_i=-\varepsilon M$. Let
$I_-=\{i:o_i=-\varepsilon\}$ and $I_+=\{i:o_i=\varepsilon\}$, and let
$k=|I_+|$. Then $|I_-|=M+k$, and since $|I_-|+|I_+|\le j$, we have
$2k\le j-M<\sqrt M\le M$. For a digit string
$e^*=(e^*_{a,q})\in\{0,1\}^{J\times d}$ and $i\in I_-\cup I_+$, we
define the target digits $t_{i,a,q}=e^*_{a,q}$ if $i\in I_-$ and
$t_{i,a,q}=1-e^*_{a,q}$ if $i\in I_+$, and we let
$\chi_{i,a,q}=|s_{a,q}(\zeta_i)-t_{i,a,q}|$. We claim that if the
product in \eqref{lat:alias-factor} does not vanish, then there is a
string $e^*$ such that
\begin{equation}\label{lat:mismatch}
 \sum_{i\in I_-\cup I_+}\chi_{i,a,q}\le k+K_a
 \qquad\text{for all }a,q.
\end{equation}
Indeed, we fix $a$ and $q$, and we let
$X=-\varepsilon\Phi_{a,q}=\sum_{I_-}s_{a,q}(\zeta_i)-\sum_{I_+}s_{a,q}(\zeta_i)$.
Since $0\le s_{a,q}\le1$, we have $-k\le X\le M+k$, so that
$-M/2<X<3M/2$. Since $\widehat\Xi_a(\Phi_{a,q})\ne0$,
\cref{lat:gate} implies that there exists an integer $l'$ with
$|\Phi_{a,q}-l'M|\le K_a$. With $l=-\varepsilon l'$, we have
$|X-lM|\le K_a\le M/4$ by \eqref{lat:band}, and hence $l\in\{0,1\}$. We
set $e^*_{a,q}=l$. Since $|I_-|=M+k$ and $|I_+|=k$, the sum in
\eqref{lat:mismatch} equals $X+k$ if $l=0$ and $M+k-X$ if $l=1$. In
both cases it is at most $k+K_a$, which proves the claim. By
\eqref{lat:soft-digit}, we also have
$\chi_{i,a,q}\ge\frac12\mathbf 1\{e_a(\zeta_{i,q})\ne t_{i,a,q}\}$,
where $\zeta_{i,q}$ is the $q$th block coordinate of $\zeta_i$.

Since the product in \eqref{lat:alias-factor} has modulus at most one,
its modulus is at most the sum, over the $R$ strings $e^*$, of the
indicator
that the $dJ$ inequalities \eqref{lat:mismatch} hold. We define
$\Omega_a=2\log d+4\log(J-a+2)$ for $1\le a\le J$ and
$\Omega=d\sum_a\Omega_a$. We bound the indicator of each inequality
$X'\le c'$ by $e^{\Omega_a(c'-X')}$ and multiply over $a$ and $q$. The
indicator for a given $e^*$ is therefore at most
\[
 e^{\Omega k}\exp\Bigl\{d\sum_a\Omega_aK_a\Bigr\}
 \prod_{i\in I_-\cup I_+}\exp\Bigl\{-\sum_{a,q}\Omega_a\chi_{i,a,q}\Bigr\}.
\]
With $b=J-a$, we have $\Omega_aK_a=\{2\log d+4\log(b+2)\}K(b)m$, and
since $K(b)$ decays geometrically in $b$, we obtain
$d\sum_a\Omega_aK_a\le C_{\rm gate}m$ with
$C_{\rm gate}=d\sum_{b\ge0}\{2\log d+4\log(b+2)\}K(b)<\infty$. Moreover,
$\Omega\le dJ\{2\log d+4\log(J+1)\}$ and $m=2^{Js_0}$, so that
$\Omega\le D\sqrt m$ with
$D=\sup_{J'\ge1}dJ'\{2\log d+4\log(J'+1)\}2^{-J's_0/2}<\infty$. Since
$\overline K\ge K(0)=6Q\lambda_*/\gamma_*>1$, we have $c_*<1$ and
$m\le M$. With $2k<\sqrt M$, we obtain
$\Omega k+d\sum_a\Omega_aK_a\le C_EM$, where $C_E=C_{\rm gate}+D$. For
$i\in I_-\cup I_+$, by the lower bound on $\chi_{i,a,q}$, the
independence of the digits under Lebesgue measure, and
$de^{-\Omega_a/2}=(b+2)^{-2}$, we have
\[
 \int_{\mathcal U\times\{1,2\}}|w(\zeta_i)|
     e^{-\sum_{a,q}\Omega_a\chi_{i,a,q}}\dd\zeta_i
 \le r_+\prod_{a,q}\frac{1+e^{-\Omega_a/2}}2
 \le\frac{r_+}R\exp\Bigl\{\sum_{b\ge0}\frac1{(b+2)^2}\Bigr\}
 \le\frac{\mathrm er_+}R .
\]
Each $i$ with $o_i=0$ contributes at most $2r_+$, and each of $x$ and
$y$ at most $2$. Summing over the $R$ strings $e^*$, we obtain
\[
 \|F_{\varepsilon,o}\|_1\le4A_uA_vR\,e^{C_EM}
 \Bigl(\frac{\mathrm er_+}R\Bigr)^{M+2k}(2r_+)^{j-M-2k}
 \le4A_uA_v(2\mathrm er_+)^jR^{1-M}e^{C_EM},
\]
which proves \eqref{lat:pattern-bound}.
\end{proof}

\begin{proof}[Proof of \cref{lem:lattice-priors}(d)]
Let $I=[r_-+\delta_n,r_+-\delta_n]$, so that $\rho_I=\rho_n$ in the
notation of \cref{lem:reciprocal}, with $c_I=r_0$, $d_I=h_n$, and
$\bar g_I=\sqrt{r_0^2-h_n^2}\le r_0$. On the block with label $\ell$,
we have $puv=A_uA_v\sigma_u\sigma_vS^2\iota_{\rm in}^2/r$
and $r=c_I+d_I\cos(\Theta+\phi+(\ell-1)\pi)$. Since the series
\eqref{eq:cosine-reciprocal} converges absolutely, multiplying
it by $\cos(M\Theta)$ and averaging over $\Theta$ with $\xi$ fixed keeps
only its term of order $M$. We write $E_\Theta$ for this average. Since
$M$ is even, we obtain on both blocks
\[
 E_\Theta\frac{\cos(M\Theta)}{r}=\frac{\rho_n^M\cos(M\phi)}{\bar g_I}.
\]
By \eqref{lo:sign-difference}, and since $\Theta$ is independent of
$\xi$, each of the $B$ blocks has volume $v_0/B$, and the density
parameters of all pairs have the same law, we obtain
\[
 (E_+-E_-)\int_{[0,1]^d}puv
 =\frac{2v_0A_uA_v\rho_n^M}{\bar g_I}
 E\Bigl[S^2\int_{\mathcal U}\iota_{\rm in}^2\cos(M\phi)\Bigr].
\]
By \eqref{lat:coherence}, $\cos(M\phi)=1$ on $\mathcal U^\circ$. Since
$0\le\iota_{\rm in}\le1$, \eqref{lat:collar} implies that the integral
is at least $I_{\rm in}-2|\mathcal U\setminus\mathcal U^\circ|\ge I_{\rm in}/2$
in every realization. The expectation is therefore at least
$(I_{\rm in}/2)ES^2\ge9I_{\rm in}/32$, although $S$ and $\phi$ are
dependent. With $\bar g_I\le r_0$, this gives (d) with
$c_2=9v_0I_{\rm in}/(16r_0)$.
\end{proof}
% ---------- 07-embedding.tex
\chapter{Embedding and proof of the lower bounds}
\label{sec:lower-proof}\label{emb:lower-proof}

\section{Embedding the priors in the generic model}
\label{sec:generic-embedding}

We perturb $\pi_0$ along scores that shift $a,b$ by $u/w,v/w$
(both by $(u+v)/w$ in the diagonal case). We apply
\cref{prop:reduction} to the resulting finite affine law.
The applications use \cref{prop:local-lower} with arbitrary
density intervals and scores. As in the lower bounds
of~\cite[Sections~3--4]{robins2009}, where the responses are binary, we
use a finite response space whose conditional laws are perturbed through
the nuisance functions.

We write a finite-support law for $Z$ as
$\pi_0=\sum_{k=1}^{m_0}\pi_{0,k}\delta_{z_k}$, with every
$\pi_{0,k}>0$, we let $w_0,a_0,b_0$ be as in
\eqref{eq:generic-baseline-moments}, with $w_0>0$, and we define
$\mathsf r=(U-a_0D,\,V-b_0D)^\top$, which has mean zero under $\pi_0$.
We call bounded functions $\mathsf s_u,\mathsf s_v$ on the support of
$\pi_0$ \emph{scores} for $\pi_0$ if
$\E_{\pi_0}\mathsf s_u=\E_{\pi_0}\mathsf s_v=0$ and either
\begin{equation}\label{eq:generic-scores}
 \E_{\pi_0}[\mathsf r(\mathsf s_u,\mathsf s_v)]=\mathrm{Id},
 \quad\text{or}\quad
 U=V,\ \ \mathsf s_u=\mathsf s_v,\ \ \E_{\pi_0}[(U-a_0D)\mathsf s_u]=1,
\end{equation}
where $\mathrm{Id}$ is the $2\times2$ identity matrix. We call the
second alternative the \emph{diagonal case}; in it we set
$\mathsf r_{\rm d}=U-a_0D$, and we call the common score
$\mathsf s_{\rm d}$ a \emph{diagonal score}. If the residual covariance
$\Sigma=\E_{\pi_0}[\mathsf r\mathsf r^\top]$ is positive definite, then
$(\mathsf s_u,\mathsf s_v)^\top=\Sigma^{-1}\mathsf r$ are scores. If
$U=V$ and $\E_{\pi_0}\mathsf r_{\rm d}^2>0$, then
$\mathsf s_{\rm d}=\mathsf r_{\rm d}/\E_{\pi_0}\mathsf r_{\rm d}^2$ is a
diagonal score. Given scores, we define
\begin{equation}\label{eq:generic-affine-response}
 \pi_{u,v}(\{z_k\})
 =\pi_{0,k}\{1+u\mathsf s_u(z_k)+v\mathsf s_v(z_k)\}.
\end{equation}
These probabilities sum to one because the scores have mean zero, and,
since the support is finite, they are at least $\pi_{0,k}/2$ when
$|u|+|v|$ is below a fixed threshold. In the notation of
\eqref{eq:finite-affine-law}, we have $q_{z_k0}=\pi_{0,k}$,
$q_{z_ku}=\pi_{0,k}\mathsf s_u(z_k)$, and
$q_{z_kv}=\pi_{0,k}\mathsf s_v(z_k)$.

Let $(\ell_a,\ell_b)=(u,v)$, and $(\ell_a,\ell_b)=(u+v,u+v)$ in the
diagonal case, and let $d_u=\E_{\pi_0}[D\mathsf s_u]$ and
$d_v=\E_{\pi_0}[D\mathsf s_v]$. For every $f$, we have
$\E_{\pi_{u,v}}f=\E_{\pi_0}f+u\E_{\pi_0}[f\mathsf s_u]+v\E_{\pi_0}[f\mathsf s_v]$.
By $\E_{\pi_0}\mathsf r=0$ and \eqref{eq:generic-scores}, we obtain
$\E_{\pi_{u,v}}\mathsf r=(\ell_a,\ell_b)^\top$; in the diagonal case,
$a_0=b_0$ and both entries of $\mathsf r$ equal $\mathsf r_{\rm d}$.
As $U=\mathsf r_1+a_0D$ and $V=\mathsf r_2+b_0D$, we find
\begin{equation}\label{eq:generic-affine-moments}
 w(u,v)=w_0+d_uu+d_vv,\qquad
 \E_{\pi_{u,v}}U=a_0w(u,v)+\ell_a,\qquad
 \E_{\pi_{u,v}}V=b_0w(u,v)+\ell_b,
\end{equation}
where $w(u,v)=\E_{\pi_{u,v}}D$. The regression ratios are therefore
\begin{equation}\label{eq:generic-lower-regressions}
 a(u,v)=a_0+\frac{\ell_a}{w(u,v)},\qquad
 b(u,v)=b_0+\frac{\ell_b}{w(u,v)}.
\end{equation}

Suppose that $X$ has density $p$ and the response law at $x$ is
$\pi_{u(x),v(x)}$. With $m_U=\E_{\pi_{u,v}}U$, $m_V=\E_{\pi_{u,v}}V$,
and $w=w(u,v)$, we have the identity
\[
 \frac{m_Um_V}w=a_0m_V+b_0m_U-a_0b_0w
 +\frac{(m_U-a_0w)(m_V-b_0w)}w .
\]
By \eqref{eq:generic-functional} and
\eqref{eq:generic-affine-moments}, this identity implies that
$\psi(P)=\int pF(u,v)$,
where
\begin{equation}\label{eq:generic-lower-integrand}
 F(u,v)=\E_{\pi_{u,v}}\widetilde W+\lambda\frac{\ell_a\ell_b}{w(u,v)},
 \qquad \widetilde W=W+\lambda(a_0V+b_0U-a_0b_0D).
\end{equation}
The first term is affine in $(u,v)$, and $F$ is smooth on a fixed
neighborhood of the origin on which $w(u,v)\ge w_0/2$. Since
$\ell_a\ell_b$ vanishes to second order at the origin, we have
$F_{uv}(0,0)=\lambda\,\partial_u\partial_v(\ell_a\ell_b)/w_0$, which is
$\lambda/w_0$ or $2\lambda/w_0$ in the diagonal case, and is nonzero.

\begin{proposition}[Local lower bounds]\label{prop:local-lower}
Fix $d,\alpha,\beta$ with $\theta=(\alpha+\beta)/d$, a finite-support
law $\pi_0$ for $Z$ with $w_0>0$, the functions $D,U,V,W$, and
$\lambda\ne0$. Suppose that $\mathsf s_u,\mathsf s_v$ are scores for
$\pi_0$, with $\alpha=\beta$ in the diagonal case, and let $\pi_{u,v}$
and $F$ be as above. Fix $0<r_-<1<r_+$ and $\epsilon>0$, and let $\tau$
be computed from $[r_-,r_+]$ as in \cref{sec:lower-preparation}.
\begin{enumerate}
\item[(a)] Suppose that $\theta<1/2$. For all sufficiently large $n$
there are two priors on laws $P$ of $O=(X,Z)$, under which $X$ has a
$C^\infty$ density $p$ and $Z$ given $X=x$ has law
$\pi_{u(x),v(x)}$, such that every realization satisfies
\eqref{eq:lower-density-margin}, $\|u\|_\infty\le Cn^{-\alpha/d}$,
$\|v\|_\infty\le Cn^{-\beta/d}$ (so that
$\|u\|_\infty+\|v\|_\infty=o(\delta_n)$), and, for every fixed smooth
function $K$ on a neighborhood $N$ of $(0,0)$, for all sufficiently
large $n$ depending on $K$,
\begin{equation}\label{eq:generic-angular-ratio-split}
 \begin{gathered}
 K(0,v)=0\text{ on }N\ \Longrightarrow\
     \|K(u,v)\|_{C^\alpha}\le C_K\epsilon,\\
 K(u,0)=0\text{ on }N\ \Longrightarrow\
     \|K(u,v)\|_{C^\beta}\le C_K\epsilon,
 \end{gathered}
\end{equation}
where $C_K$ does not depend on $\epsilon$ or $n$. We call the laws
under these priors the \emph{hard laws}, and we write $\cP^{\rm h}_n$
for their set. On these laws $\psi(P)=\int pF(u,v)$. If $\mathcal P$ is
a class that contains $\cP^{\rm h}_n$ and $T$ is a parameter with
$T(P)=\psi(P)+t_*$ on $\cP^{\rm h}_n$ for a constant $t_*$, then
there is $c>0$ such that, for all sufficiently large $n$,
\[
 \inf_{\widehat T}\sup_{P\in\mathcal P}
 P\{|\widehat T-T(P)|\ge c\,a_n\log n\}\ge\frac38,
 \qquad \risk(T,\mathcal P)\ge\sqrt{3/8}\,c\,a_n\log n,
\]
with $a_n=a_n(\theta,\tau)$. The infimum includes randomized
estimators.
\item[(b)] For every $u_*$ in a neighborhood of zero, the laws
$P_t=\operatorname{Unif}([0,1]^d)\otimes\pi_{u_*,t}$, for $t$ near
zero, have $\psi(P_t)=F(u_*,t)$. For every sufficiently small
$u_*\ne0$ there is $c>0$, depending also on $u_*$, such that, on every
class $\mathcal P$ that contains $P_t$ for all $t$ in an interval about
zero, a parameter $T$ with $T(P_t)=\psi(P_t)+t_*$ for a constant $t_*$
satisfies $\risk(T,\mathcal P)\ge cn^{-1/2}$ for all sufficiently large
$n$.
\end{enumerate}
The constants may depend on the fixed quantities above and on $\epsilon$.
\end{proposition}

\begin{proof}
For (a), we apply \cref{prop:reduction} to the response law
\eqref{eq:generic-affine-response}, which has the form
\eqref{eq:finite-affine-law}, with the integrand $F$ and the interval
$[r_-,r_+]$. Let $M$ be as in that proposition, let $c_*$ be the
constant of \cref{sec:block-phases}, and let
$s_0=\min(\alpha,\beta)$. For all large $n$ we have $c_*M\ge1$, and we
set
\[
 J=\Bigl\lfloor\frac{\log_2(c_*M)}{s_0}\Bigr\rfloor\ge0,\qquad
 m=2^{Js_0},\qquad R=2^{Jd},
\]
so that $2^{-s_0}c_*M<m\le c_*M$. Since $s_0\le\alpha,\beta$, we have
$m=R^{s_0/d}\le\min(R^{a_*},R^{b_*})$, and
$\log R=(d/s_0)\log m=O(\log M)$. Let
$\epsilon_u=\epsilon_v=\min(\epsilon,1)$. By
\cref{lem:lattice-priors}(a), (c), and~(d), for all large $n$ and
every $B\ge n$ that is the $d$th power of an even integer, the priors of
\cref{sec:lattice} with these $M$, $J$, $B$, $\epsilon_u$, and
$\epsilon_v$ satisfy the hypotheses of
\cref{prop:reduction}, with constants that do
not depend on $n$ or $B$. By \cref{prop:reduction}, there
exists such a $B$ with $A_u\le\epsilon_uB^{-a_*}\le n^{-\alpha/d}$ and
$A_v\le\epsilon_vB^{-b_*}\le n^{-\beta/d}$, and the hard laws are the
laws of one observation under the realizations of the two priors for
this $B$. By \cref{lem:lattice-priors}(a) and~(b), every
realization satisfies \eqref{eq:lower-density-margin},
$\|u\|_\infty\le A_u/r_-$, and $\|v\|_\infty\le A_v/r_-$. Since
$A_u+A_v\to0$, the response probabilities are positive for all large
$n$, and \eqref{eq:lower-smooth-map} implies
\eqref{eq:generic-angular-ratio-split} for all large $n$ depending on
$K$, since $\epsilon_u,\epsilon_v\le\epsilon$. Since
$n^{-\min(\alpha,\beta)/d}(\log n)^2\to0$, we have
$\|u\|_\infty+\|v\|_\infty=o(\delta_n)$.

By \eqref{eq:generic-lower-integrand}, $\psi(P)=\int pF(u,v)$ on the
hard laws. We apply \cref{prop:reduction} to $T-t_*$, whose
estimation errors equal those of $T$, and this yields the threshold
$c_1m^2a_n$ in \eqref{eq:lower-testing-transfer}. Since
$m>2^{-s_0}c_*M$ and
$M\ge\sqrt{\theta(1-2\theta)\log n/\tau}-1$, we have
\[
 m^2\ge2^{-2s_0-2}c_*^2\,\frac{\theta(1-2\theta)}{\tau}\log n
\]
for all large $n$, which gives the testing bound with the threshold
$c\,a_n\log n$ for $c=2^{-2s_0-2}c_*^2\theta(1-2\theta)c_1/\tau$.
The mean squared error is at least the squared threshold times the error
probability, which gives the RMSE bound. With $J=0$ in place of this
choice, so that $m=R=1$, these bounds hold with $a_n$ in place of
$a_n\log n$, by the same argument.

For (b), the constant perturbations $u=u_*$ and $v=t$, with $p=1$, give
the law $P_t$, so that $\psi(P_t)=F(u_*,t)$. Let $\varphi(u)=F_v(u,0)$.
Since $\varphi'(0)=F_{uv}(0,0)\ne0$, we have $\varphi(u_*)\ne0$ for
every sufficiently small $u_*\ne0$, by continuity if $\varphi(0)\ne0$,
and since $\varphi(u)=u\{\varphi'(0)+o(1)\}$ otherwise. Relative to
$P_*=\operatorname{Unif}([0,1]^d)\otimes\pi_0$, the law $P_t$ has the
density $1+u_*\mathsf s_u(Z)+t\mathsf s_v(Z)$, which is at least $1/2$
when $|u_*|+|t|$ is below a fixed threshold and has the bounded slope
$\mathsf s_v(Z)$ in $t$. On an open interval about zero on which
$P_t\in\mathcal P$, the map $t\mapsto T(P_t)=F(u_*,t)+t_*$ has the
derivative $\varphi(u_*)\ne0$ at $t=0$. The bound now follows from
\cref{lem:parametric-path}.
\end{proof}

\section{Localization}\label{sec:hard-laws}

Every conditional moment of a hard law is affine in $(u,v)$:
$\E_P[f(Z)\mid X]=\E_{\pi_0}f+u\E_{\pi_0}[f\mathsf s_u]+v\E_{\pi_0}[f\mathsf s_v]$.
Every conditional-response nuisance used below, such as a conditional
mean, a ratio of conditional means whose denominator is positive at the
baseline, a propensity, or its reciprocal, is therefore of the form
$K_P(x)=K(u(x),v(x))$ for a $C^\infty$ function $K$ near the origin.

\begin{lemma}[Localization and membership]
\label{lem:local}\label{cor:membership}
Adopt the setting of \cref{prop:local-lower}. Let $K$ be a
$C^\infty$ function on a neighborhood of $(0,0)$, let $k_0=K(0,0)$, and
for a law $P$ with perturbations $u,v$ let $K_P(x)=K(u(x),v(x))$.
Let $\gamma_K=\alpha$ if $K(0,v)=k_0$ for all small $v$, let
$\gamma_K=\beta$ if $K(u,0)=k_0$ for all small $u$, and let
$\gamma_K=\min(\alpha,\beta)$ otherwise. If both identities hold,
either value may be used. There is $C_K$, independent of $n$ and
$\epsilon$, such that, for all sufficiently large $n$, every hard law
$P$ satisfies
\[
 \|K_P-k_0\|_\infty
 \le C_K(\|u\|_\infty+\|v\|_\infty)=o(\delta_n),\qquad
\|K_P-k_0\|_{C^t}\le C_K\epsilon\quad(0<t\le\gamma_K),
\]
and, if $k_0>0$, also $r_-k_0\le pK_P\le r_+k_0$. For the laws $P_t$
of \cref{prop:local-lower}(b), $K_{P_t}$ is the constant
$K(u_*,t)$, which tends to $k_0$ as $(u_*,t)\to0$.

Consider finitely many conditions on $(p,u,v)$ of the following three
types, each with its own $C^\infty$ function $K$ and $k_0=K(0,0)$:
\begin{enumerate}
\item[(H)] $\|K_P\|_{C^t}\le H'$, with $0<t\le\gamma_K$ and $|k_0|<H'$;
\item[(I)] $k_-\le K_P\le k_+$, with $k_-<k_0<k_+$;
\item[(W)] $g_-'\le pK_P\le g_+'$, with $k_0>0$,
$g_-'\le r_-k_0$, and $r_+k_0\le g_+'$.
\end{enumerate}
If $\epsilon$ is small enough, depending only on these conditions, then
for all large $n$ every hard law satisfies all of them. Moreover, for
every sufficiently small $u_*$, the laws $P_t$, which have $p=1$,
satisfy all of them for $t$ in an interval about zero; this assertion
does not use $\theta<1/2$.
\end{lemma}

\begin{proof}
Shrinking the domain of $K$, we may assume that it is an open square $N$
centered at the origin, so that $(u,0)$ and $(0,v)$ lie in $N$ whenever
$(u,v)$ does, and that $K$ is Lipschitz on $N$ with a constant $C_K$.
By \cref{prop:local-lower}(a), we have
$\|u\|_\infty+\|v\|_\infty\le Cn^{-\min(\alpha,\beta)/d}=o(\delta_n)$,
uniformly over the hard laws, which gives the bound in the supremum
norm.

For the H\"older bounds, we write
\[
 K(u,v)-k_0=\{K(u,v)-K(0,v)\}+\{K(0,v)-k_0\}
\]
on $N$. The first brace vanishes when $u=0$, and the second when $v=0$.
By \eqref{eq:generic-angular-ratio-split}, for all large $n$, their
compositions with $(u,v)$ have norms at most $C\epsilon$ in $C^\alpha$
and in $C^\beta$, respectively, where $C$ does not depend on $\epsilon$.
If $K(0,v)=k_0$, only the first brace remains, which gives the bound in
$C^\alpha$. If $K(u,0)=k_0$, then $K-k_0$ itself vanishes when $v=0$,
and the bound in $C^\beta$ follows from
\eqref{eq:generic-angular-ratio-split}.
In general, both braces are bounded in $C^{\min(\alpha,\beta)}$, since
$\|f\|_{C^t}\le(1+d)\|f\|_{C^{t'}}$ for $0<t\le t'$ and $f$ on
$[0,1]^d$. This follows from $|x-y|\le\sqrt d$ on $[0,1]^d$ and, if
$\lceil t\rceil<\lceil t'\rceil$, from the mean value theorem, which
shows that the derivatives of order $\lceil t\rceil-1$ are Lipschitz
with constant $\sqrt d\|f\|_{C^{t'}}$ on the convex set $[0,1]^d$.
Using the same comparison and enlarging $C_K$, we find that the bounds
hold for every $0<t\le\gamma_K$.

Let $k_0>0$ and $\eta=\|K_P-k_0\|_\infty$. Since $\eta=o(\delta_n)$, for
all large $n$ we have $K_P>0$ and $r_+\eta\le k_0\delta_n$, and
\eqref{eq:lower-density-margin} gives
\[
 pK_P\ge(r_-+\delta_n)(k_0-\eta)
 \ge r_-k_0+k_0\delta_n-r_+\eta\ge r_-k_0 ,
\]
and in the same way
$pK_P\le(r_+-\delta_n)(k_0+\eta)\le r_+k_0-k_0\delta_n+r_+\eta\le r_+k_0$.

Since a constant has H\"older norm equal to its absolute value, we
obtain (H) from $\|K_P\|_{C^t}\le|k_0|+C_K\epsilon\le H'$, once
$C_K\epsilon\le H'-|k_0|$. The bound in the supremum norm implies (I),
and the bounds on $pK_P$ imply (W). For the laws $P_t$, we have $p=1$
and $K_{P_t}=K(u_*,t)$, a constant. The three conditions hold strictly
at $(u_*,t)=(0,0)$, for (W) because
$g_-'\le r_-k_0<k_0<r_+k_0\le g_+'$. By continuity, they hold for
$(u_*,t)$ near zero.
\end{proof}

\begin{proof}[Proof of \cref{thm:generic}(b)]
Since the minimax risk and the minimax error probability do not
decrease when the class grows, it suffices to work on $\cP_r(\pi_0)$.
We use the scores $\Sigma^{-1}\mathsf r$ in the positive definite case
of \eqref{eq:generic-nondegeneracy}, and the diagonal score
$\mathsf r_{\rm d}/\E_{\pi_0}\mathsf r_{\rm d}^2$ in the diagonal case;
both exist by \eqref{eq:generic-nondegeneracy}. Let $r_-=g_-/w_0$ and
$r_+=g_+/w_0$, so that $0<r_-<1<r_+$ by
\eqref{eq:generic-nondegeneracy}. Since scaling both endpoints of an
interval by the same factor leaves $\rho$ unchanged, and since
$a_*+b_*=\theta$, the $\tau$ and $a_n$ of \cref{prop:local-lower} are those of
\eqref{eq:sharp-parameters} and \eqref{eq:subcritical-scale}.

Under a hard law or a law $P_t$, the functions $w,a,b$ of
\cref{sec:models} are $K_P$ for the functions
$w(u,v),a(u,v),b(u,v)$ of \eqref{eq:generic-affine-moments} and
\eqref{eq:generic-lower-regressions}, with $k_0=w_0,a_0,b_0$, and
$g=pw$. Membership in $\cP_r(\pi_0)$ consists of the conditions
$\max(\delta,w_0-r)\le w\le w_0+r$, $|a-a_0|\le r$, and
$|b-b_0|\le r$, of type (I); $g_-\le pw\le g_+$, of type (W) with
$r_\pm w_0=g_\pm$; and $\|a\|_{C^\alpha}\le H$ and
$\|b\|_{C^\beta}\le H$, of type (H). By
\eqref{eq:generic-nondegeneracy}, the baseline values satisfy the
conditions of types (I) and (H) strictly. For (H), we may take
$\gamma_K=\alpha$ for $a$ and $\gamma_K=\beta$ for $b$, since
$a(0,v)=a_0$ and $b(u,0)=b_0$ in the positive definite case, and
$\alpha=\beta$ in the diagonal case. We fix $\epsilon$ as in
\cref{lem:local} for these conditions. Then, for all large $n$,
every hard law belongs to $\cP_r(\pi_0)$, and, for every sufficiently
small $u_*\ne0$, so does $P_t$ for every $t$ in an interval about zero.

If $0<\theta<1/2$, we apply \cref{prop:local-lower}(a) with
$\mathcal P=\cP_r(\pi_0)$, $T=\psi$, and $t_*=0$, and we find $c_0>0$
such that, for all large $n$, the minimax probability of an error of at least
$c_0a_n\log n$ is at least $3/8$, and
$\risk(\psi,\cP_r(\pi_0))\ge\sqrt{3/8}\,c_0a_n\log n$. With $c=c_0/2$,
this gives \eqref{eq:sharp-probability} and the lower bound in
\eqref{eq:sharp-bracket}, since $3/8\ge1/4$ and $\sqrt{3/8}\ge1/2$. For
every $\theta>0$, \cref{prop:local-lower}(b), with
$\mathcal P=\cP_r(\pi_0)$ and $T=\psi$, implies that there exists
$c'>0$ such that
$\risk(\psi,\cP_r(\pi_0))\ge c'n^{-1/2}$ for all large $n$, which is the
lower bound in \eqref{eq:critical-rate} when $\theta\ge1/2$. For $\theta\ge1/2$, we take $c=c'$ and the threshold of the
root-$n$ argument. For $0<\theta<1/2$, we take $c=\min\{c_0/2,c'\}$ and
the larger of the two thresholds for $n$. This gives all the
assertions.
\end{proof}
% ---------- 08-conclusion.tex
\chapter{Conclusion}\label{sec:conclusion}

First, when $s<d/4$, the bounds of \cref{thm:generic} differ by the
factor $(\log n)^{\nu/2-3/4}$, which is $(\log n)^{1/4}$ when
$0<\alpha,\beta\le1$. In the estimator, the factor $(\log n)^{1/4}$
comes from the number of resolutions near the dominant one
(\cref{sec:ideas}). Each of our lower-bound constructions, in
contrast, places its bumps at a single resolution. Closing this gap
would determine the minimax rate up to constants. The power of
$\log n$ is also left open in~\cite{park}.

Second, our estimator uses $\alpha,\beta$, the bounds $g_-,g_+$, and
the remaining class parameters. McClean et al.~\cite{mcclean} and
Park~\cite[Corollary~1]{park} attain the rough-design exponent for the
expected conditional covariance without knowing the nuisance
smoothness, and adaptive results for related functionals under smooth
density assumptions were obtained in~\cite{adaptive}. Park also
reports an estimator that does not use
$\alpha,\beta$ and attains his upper bound up to a larger power of
$\log n$, uniformly over smoothness indices in a known compact subset
of $\{\alpha+\beta<d/2\}$~\cite[Section~3.4]{park}. This estimator is
analyzed in Theorem~S13 of the supplement of~\cite{park}. For the
general class, it is open whether the upper bound of
\cref{thm:generic}, including the correction, can be attained without
knowledge of $\alpha,\beta$. In the setting of~\cite{park}, it is
open whether this is possible without the larger power of $\log n$.

Third, minimax rates remain unresolved when $g$ belongs to a H\"older
class of positive smoothness that fails the condition stated before
\eqref{eq:known-weight-rate}. Park~\cite[Section~6.2]{park} bounds the
risk of his selected estimator of the expected conditional covariance
at each fixed H\"older density, and he notes that this is not a minimax
result over a class of smooth densities.
The question of how the minimax rate
depends on the smoothness of the covariate density is also described as
open in~\cite[Remark~5]{generalhoif}. When the two nuisance functions
have equal smoothness indices, the conjecture following equation~(1.3)
of Robins et al.~\cite{hoif2008} addresses this case; our results
confirm the exponent it suggests in the rough-design limit.
\appendix
% cleveref's own \appendix hook is inert under the current kernel, whose label hook takes
% the reference type from \crefalias; this files the appendices as "Appendix A", "Appendix A.1".
\crefalias{chapter}{appendix}
\crefalias{section}{subappendix}
% ---------- A-reductions.tex
\chapter{Reductions}\label{sec:application-reductions}\label{sec:application-proofs}

Most targets of \cref{sec:applications} are related to a target
of \cref{thm:generic} through an observable mean, a Lipschitz map,
or a transformation of the data that does not depend on the unknown law.
We record the lower-bound constructions and these three reductions.

\section{Hard laws}\label{sec:app-hard-laws}

\begin{definition}[Hardness]\label{def:hard}
Let $\cP_n$ be classes of laws of one observation, let $T$ be a real
parameter defined on each $\cP_n$, and let $t_n>0$. We say that $T$ is
\emph{hard at scale $t_n$ on $(\cP_n)$} if
\[
 \liminf_{n\to\infty}\ \inf_{\widehat T}\ \sup_{P\in\cP_n}
 \Pp_P\{|\widehat T-T(P)|\ge t_n\}\ge\frac38,
\]
where, for each $n$, the infimum is over all estimators, including
randomized ones, based on $n$ independent observations.
\end{definition}

Hardness at scale $t_n$ implies hardness at scale $ct_n$ for every fixed
$c\in(0,1]$, and on every sequence of classes $\cP_n'\supseteq\cP_n$.

\begin{lemma}[From hardness to risk]
\label{lem:hard-risk}\label{lem:hardness-bracket}
Let $T$ be hard at scale $t_n$ on $(\cP_n)$ and a real parameter on a
class $\cP$. Suppose that $\cP_n\subseteq\cP$ for all large $n$.
Then, for all large $n$,
\begin{equation}\label{eq:hard-consequence}
 \inf_{\widehat T}\sup_{P\in\cP}\Pp_P\{|\widehat T-T(P)|\ge t_n\}\ge\frac14,
 \qquad \risk(T,\cP)\ge\frac{t_n}2 .
\end{equation}
The infimum includes randomized estimators. In particular, if
$0<\theta<1/2$ and $t_n=c_1\underline a_n(\theta,\tau_I)$ for a fixed
$c_1>0$ and an interval $I$ as in \cref{def:bracket}, then
$(T,\cP)$ satisfies the lower bracket $(\theta,I)$ with $c=c_1/2$.
\end{lemma}
\begin{proof}
By the $\liminf$ bound, the minimax probability is at least $1/4$ over
$\cP_n$, and hence over $\cP\supseteq\cP_n$, for all large $n$. For every estimator, we have
$\sup_{P\in\cP}\E_P(\widehat T-T(P))^2\ge t_n^2\sup_{P\in\cP}\Pp_P\{|\widehat T-T(P)|\ge t_n\}$,
which gives the RMSE bound and, with $t_n=c_1\underline a_n=2c\,\underline a_n$,
the lower bracket.
\end{proof}

We fix the inputs of \cref{prop:local-lower}: a finite-support
law $\pi_0$ for the response, the functions $D,U,V,W$ and $\lambda$,
scores for $\pi_0$ (in the diagonal case, one function $\mathsf s_{\rm d}$
with $\E_{\pi_0}\mathsf s_{\rm d}=0$ and
$\E_{\pi_0}[(U-a_0D)\mathsf s_{\rm d}]=1$, used for both perturbations),
smoothness indices $\alpha,\beta$ with $\theta=(\alpha+\beta)/d<1/2$, an
interval $[r_-,r_+]$ with $0<r_-<1<r_+$, and $\epsilon>0$. By
\cref{prop:local-lower}(a), there exists a fixed $c_1>0$ such
that, for every constant $t_*$, the parameter $\psi+t_*$ is hard at
scale $c_1a_n(\theta,\tau)\log n$ on the hard laws $(\cP^{\rm h}_n)$,
where $\tau$ is computed from $[r_-,r_+]$ and $\psi$ is the target of
\cref{sec:models} for $(D,U,V,W,\lambda)$. On a hard law, the
response law at $x$ is the perturbation of $\pi_0$ by $(u(x),v(x))$, and
$\|u\|_\infty\le Cn^{-\alpha/d}$, $\|v\|_\infty\le Cn^{-\beta/d}$.
For a fixed finite family of restrictions of the types (H), (I),
and (W) of \cref{lem:local}, we take $\epsilon$ sufficiently small.
Then, by that lemma, every hard law satisfies them for
all large $n$, and the laws $P_t$ of
\cref{prop:local-lower}(b) satisfy them for every
sufficiently small $u_*$ and all $t$ in an interval about zero.
The reductions preserve hardness, and the lower brackets on the
classes of \cref{sec:applications} follow from
\cref{lem:hard-risk}.

\section{Three reductions}\label{sec:transfers}

Let $h$ be a bounded measurable map from the observation space into
$\mathbb R^k$, with $k\ge0$, and write $\mu(P)=\E_Ph(O)$. Let
$\mathcal R$ be a compact rectangle that contains $\mu(P)$ for every
law $P$ that we consider, and let
$b_h^2=\sup_P\E_P\|h(O)-\mu(P)\|_2^2$. When $k=0$, the terms with $h$
and $\mu$ are absent and $b_h=0$. Let $I_1,\ldots,I_m$ be compact
intervals. A map $\Phi$ on $I_1\times\cdots\times I_m\times\mathcal R$ is
\emph{$L$-Lipschitz} if it is Lipschitz with constant $L\ge1$ for the
sum of the absolute differences of the real arguments plus the
Euclidean distance in $\mathcal R$.

\begin{lemma}[Reductions]\label{lem:reductions}
\leavevmode
\begin{enumerate}[label=(\alph*)]
\item (Combination.) Suppose that $T_1,\ldots,T_m$ are real parameters
on a class $\cP$ with $T_i(P)\in I_i$, and that
$T_0(P)=\Phi(T_1(P),\ldots,T_m(P),\mu(P))$ for all $P\in\cP$, where
$\Phi$ is $L$-Lipschitz. Then, for every $n\ge1$,
\[
 \risk(T_0,\cP)\le L\sum_{i=1}^m\risk(T_i,\cP)+L\,b_h\,n^{-1/2}.
\]
\item (Reverse relations.) Let $m=1$, let $\cP'$ be a class, and
suppose that $T_0(P)\in I_1$ and
\begin{equation}\label{eq:transfer-relation}
 |T_1(P)-\Phi(T_0(P),\mu(P))|\le e\qquad(P\in\cP'),
\end{equation}
where $\Phi$ is $L$-Lipschitz and $e\ge0$. Then, for every $n\ge1$ and
every $t>0$ with $t\ge2e$,
\begin{gather}
 \risk(T_1,\cP')\le L\risk(T_0,\cP')+L\,b_h\,n^{-1/2}+e,
 \label{eq:transfer-rmse}\\
 \inf_{\widehat T_0}\sup_{P\in\cP'}
 \Pp_P\Bigl\{|\widehat T_0-T_0(P)|\ge\frac t{4L}\Bigr\}
 \ge\inf_{\widehat T_1}\sup_{P\in\cP'}
 \Pp_P\{|\widehat T_1-T_1(P)|\ge t\}-\frac{16L^2b_h^2}{nt^2}.
 \label{eq:moment-probability-transfer}
\end{gather}
Consequently, suppose that for all large $n$ these hypotheses hold
with $\cP'=\cP'_n$ and $e=e_n\le t_n/2$, where $h$, $\mathcal R$,
$I_1$, and $\Phi$ do not depend on $n$. If $T_1$ is hard at scale $t_n$
on $(\cP'_n)$, and if $k=0$ or $nt_n^2\to\infty$, then $T_0$ is hard at
scale $t_n/(4L)$ on $(\cP'_n)$.
\item (Kernels.) Let $K$ be a Markov kernel from the observation space
of a class $\cP_0$ to that of a class $\cP_1$ that does not depend on
the unknown law, and let $KP$ be the law of the output when the input
has law $P$. Suppose that $KP\in\cP_1$ and $T_1(KP)=T_0(P)$ for every
$P\in\cP_0$. Then, for every $n\ge1$ and $t>0$,
\[
 \inf_{\widehat T_0}\sup_{P\in\cP_0}\Pp_P\{|\widehat T_0-T_0(P)|\ge t\}
 \le\inf_{\widehat T_1}\sup_{Q\in\cP_1}\Pp_Q\{|\widehat T_1-T_1(Q)|\ge t\},
\]
and $\risk(T_0,\cP_0)\le\risk(T_1,\cP_1)$. In particular, upper
bounds for $T_1$ on $\cP_1$ hold for $T_0$ on $\cP_0$, lower bounds
and probability bounds for $T_0$ on $\cP_0$ hold for $T_1$ on $\cP_1$,
and if $T_0$ is hard at scale $t_n$ on $(\cP_{0,n})$, then $T_1$ is
hard at scale $t_n$ on $(K\cP_{0,n})$.
\end{enumerate}
\end{lemma}

\begin{proof}
Let $\widehat\mu$ be the coordinatewise projection onto $\mathcal R$ of
the empirical mean of $h(O_1),\ldots,h(O_n)$, and let $\Pi_i$ be the
projection onto $I_i$. Since these projections do not increase distance to
points of $\mathcal R$ and $I_i$, we have
$\E_P\|\widehat\mu-\mu(P)\|_2^2\le b_h^2/n$. Suppose that real
parameters $S,S_1,\ldots,S_m$ satisfy $S_i(P)\in I_i$ and
$|S(P)-\Phi(S_1(P),\ldots,S_m(P),\mu(P))|\le e$. Given estimators
$\widehat S_i$, the estimator
$\widehat S=\Phi(\Pi_1\widehat S_1,\ldots,\Pi_m\widehat S_m,\widehat\mu)$
satisfies
\[
 |\widehat S-S(P)|
 \le L\Bigl\{\sum_{i=1}^m|\widehat S_i-S_i(P)|
 +\|\widehat\mu-\mu(P)\|_2\Bigr\}+e .
\]
With $S=T_0$, $S_i=T_i$, and $e=0$, (a) follows from Minkowski's
inequality in $L^2(P^n)$ by taking suprema and infima, without
independence between estimators. With $m=1$, $S=T_1$, and $S_1=T_0$, the same argument on
$\cP'$ proves \eqref{eq:transfer-rmse}. If moreover
$|\widehat S-T_1(P)|\ge t\ge2e$, the display implies that
$|\widehat T_0-T_0(P)|\ge t/(4L)$ or
$\|\widehat\mu-\mu(P)\|_2\ge t/(4L)$, and by Chebyshev's inequality the
latter event has probability at most $16L^2b_h^2/(nt^2)$. Since
$\widehat S$ is an estimator of $T_1$,
\eqref{eq:moment-probability-transfer} follows by taking the supremum
over $P\in\cP'$ and then the infimum over $\widehat T_0$. Applying it
with $t=t_n$ and taking the $\liminf$ gives the hardness assertion,
since the error term
vanishes when $k=0$ and otherwise tends to zero by assumption.

For (c), we apply $K$ independently to the $n$ observations, with
randomness independent of the sample, and evaluate $\widehat T_1$ on the
outputs. Since, under $P$, the outputs have law $(KP)^n$ and
$T_1(KP)=T_0(P)$, the resulting error has the law of the error of
$\widehat T_1$ under $KP\in\cP_1$. The suprema and infima give
both inequalities and the remaining assertions.
\end{proof}

\Cref{lem:reductions}(c) is the elementary direction of Le Cam's
comparison of experiments by randomization~\cite[Chapter~2]{lecam}: an
experiment is at least as informative as its image under a Markov kernel
that does not depend on the unknown law. Composing an estimator for the
image with the kernel, we obtain an estimator whose error has the same
law. This gives both the inequality for error probabilities and the
inequality for risks.

For a ratio we use (a) with $\Phi(x,y)=x/y$ on
$[-C_0,C_0]\times[\delta_{\mathcal D},C_0]$, which is $L$-Lipschitz with
$L=\max\{1,\delta_{\mathcal D}^{-1},C_0\delta_{\mathcal D}^{-2}\}$, since
$|x/y-x'/y'|\le|x-x'|/\delta_{\mathcal D}+C_0|y-y'|/\delta_{\mathcal D}^2$.
A target that differs from another by a sign and an additive constant
has the same estimation errors, so that all bounds transfer between
the two.

\begin{remark}[Combining rates]\label{rem:combining-rates}
For fixed $\tau$ and $\nu$, and $0<\theta<1/2$, we have
$\overline a_n(\theta,\tau,\nu)=n^{-\theta}e^{O(\sqrt{\log n})}$, since
$(\log n)^c=e^{o(\sqrt{\log n})}$ for every fixed $c$. Hence, if
$\theta_1<\min(\theta_2,1/2)$, then
$\overline a_n(\theta_2,\tau,\nu_2)+n^{-1/2}=o(\overline a_n(\theta_1,\tau,\nu_1))$,
and if $\theta_1=\theta_2$ and $\nu_1\le\nu_2$, then
$\overline a_n(\theta_1,\tau,\nu_1)\le\overline a_n(\theta_2,\tau,\nu_2)$.
For $\theta\ge1/2$ all these rates are $n^{-1/2}$. It follows from
\cref{lem:reductions}(a) that if each $(T_i,\cP)$ satisfies the
upper bracket $(\theta_i,I,\nu_i)$ with a common interval $I$, then
$(T_0,\cP)$ satisfies the upper bracket $(\theta,I,\nu)$ with
$\theta=\min_i\theta_i$ and $\nu=\max\{\nu_i:\theta_i=\theta\}$.
\end{remark}
% ---------- B-missing.tex
\chapter{Missing-data means and treatment effects}\label{sec:causal-extension}

The treatment effects inherit lower bounds below the threshold from one
MAR family through a kernel with the observed regression in one arm and
the constant $1/2$ in the other. Their arm means inherit upper bounds
for the MAR mean by keeping one arm.

\section{The MAR family and two maps}\label{sec:binary-embedding}

We take $Z=(R,RY)$ with $R,Y\in\{0,1\}$ and
$(D,U,V,W,\lambda)=(R,1,RY,0,1)$, and we let $\pi_0$ give
$(R,RY)=(0,0),(1,0),(1,1)$ the probabilities $1/2,1/4,1/4$. Then
$(w_0,a_0,b_0)=(1/2,2,1/2)$, and the residual
$\mathsf r=(1-2R,\,RY-R/2)^\top$ takes the values $(1,0)$,
$(-1,-1/2)$, and $(-1,1/2)$ at the three support points, so that
$\Sigma=\Cov_{\pi_0}(\mathsf r)=\operatorname{diag}(1,1/8)$. Since
$\Sigma$ is positive definite, the functions
$(\mathsf s_u,\mathsf s_v)^\top=\Sigma^{-1}\mathsf r=(1-2R,\,8RY-4R)^\top$
are scores for $\pi_0$, as in \cref{sec:generic-embedding}. The perturbed response law
$\pi_{u,v}(\{z\})=\pi_0(\{z\})\{1+u\mathsf s_u(z)+v\mathsf s_v(z)\}$ is
\begin{equation}\label{eq:binary-embedding}
 \begin{gathered}
 P(R=0\mid X)=\frac{1+u}2,\\
 P(R=1,Y=0\mid X)=\frac{1-u}4-v,\quad
 P(R=1,Y=1\mid X)=\frac{1-u}4+v,
 \end{gathered}
\end{equation}
so that
\begin{equation}\label{eq:binary-moments}
 w=\frac{1-u}2,\qquad \frac1w=\frac2{1-u},\qquad
 \frac1{1-w}=\frac2{1+u},\qquad b=\frac12+\frac{2v}{1-u},
\end{equation}
and $\psi(P)=\int abg=\E_P[b(X)]$ is the MAR mean. These observed laws
are those of full-data laws in which $Y$ is Bernoulli with mean $b(X)$
and independent of $R$ given $X$. The lower bound for the missing-data
mean in~\cite[Section~3]{robins2009} uses binary models of this kind.

For indices $(\alpha,\beta')$ with $(\alpha+\beta')/d<1/2$ and an
interval $[r_-,r_+]$ with $0<r_-<1<r_+$, we call the hard laws of
\cref{sec:app-hard-laws} for this $\pi_0$ and these scores the
\emph{MAR family}. The MAR mean is hard on it at scale
$c_1\underline a_n((\alpha+\beta')/d,\tau)$, with $\tau$ computed from
$[r_-,r_+]$. In the notation of \cref{lem:local}, the functions
$w$, $1-w$, $1/w$, and $1/(1-w)$ are of the form $K_P$ with $K(0,v)$
constant, so that $\gamma_K=\alpha$, and $b$ is of the form $K_P$ with
$K(u,0)=1/2$, so that $\gamma_K=\beta'$.

For $j\in\{0,1\}$, the map that keeps arm $j$ sends $(X,A,Y)$ to
$(X,R,RY)$ with $R=\mathbf 1_{\{A=j\}}$. If the input has law $P$, then
the output has covariate density $p$, observation probability $w_j$,
and regression $b_j$, so that its MAR mean is $\mu_j(P)$. The
kernel $K_j$ sends $(X,R,RY)$ to $(X,A,Y')$ with
\[
 A=R\ \text{ if }j=1,\qquad A=1-R\ \text{ if }j=0,\qquad
 Y'=RY+(1-R)\xi,
\]
where $\xi\sim\operatorname{Bernoulli}(1/2)$ is drawn independently. If
the input has covariate density $p$, observation probability $w$, and
regression $b$, then the output has covariate density $p$,
$P(A=j\mid X)=w$, $P(A=1-j\mid X)=1-w$, $\E[Y'\mid A=j,X]=b$, and
$\E[Y'\mid A=1-j,X]=1/2$. In particular, $\mu_j(K_jP)$ is the MAR mean
of $P$, $\mu_{1-j}(K_jP)=1/2$, and
\begin{equation}\label{eq:kernel-ate}
 \mu_j(K_jP)=\frac12+(2j-1)\,T_{\rm ATE}(K_jP).
\end{equation}
If $w$ and $b$ are constant and $p=1$, then all regressions of
$K_jP$ are constant, and
\begin{equation}\label{eq:kernel-constant}
 T_{\rm ATE}(K_jP)=T_{\rm ATT}(K_jP)=T_{\rm ATU}(K_jP)
 =(2j-1)\Bigl(\psi(P)-\frac12\Bigr).
\end{equation}

On every law of treatment observations, we write $p_A=\E_PA$. Since
$\E_PY-\mu_0=\int w(b_1-b_0)p=p_AT_{\rm ATT}$, and similarly for the
ATU, we have
\begin{equation}\label{eq:att-atu-identities}
 \begin{aligned}
 T_{\rm ATT}&=\frac{\E_PY-\mu_0}{p_A},
 &\mu_0&=\E_PY-p_AT_{\rm ATT},\\
 T_{\rm ATU}&=\frac{\mu_1-\E_PY}{1-p_A},
 &\mu_1&=\E_PY+(1-p_A)T_{\rm ATU}.
 \end{aligned}
\end{equation}
On $\mathcal T_\eta$ and $\mathcal T^{\rm sep}$, the overlap condition
implies that $p_A\in[\eta,1-\eta]$ and $p_A\in[\epsilon,1-\epsilon]$,
respectively. With $h(O)=(Y,A)$ and $\mathcal R=[0,1]\times[\eta,1-\eta]$
or $[0,1]\times[\epsilon,1-\epsilon]$, the maps on the left of
\eqref{eq:att-atu-identities} are Lipschitz on $[0,1]\times\mathcal R$,
as functions of $\mu_j$ and $(\E_PY,p_A)$, and those on the right are
Lipschitz on $[-1,1]\times\mathcal R$, as functions of the effect and
$(\E_PY,p_A)$.

\section{A fixed-accuracy pilot}\label{sec:broad-inclusions}

The model $\mathcal M^{\rm sep}$ bounds $g=wp$ only by $\epsilon p_-$
and $(1-\epsilon)p_+$. We divide the responses by a pilot estimate
$\widehat w$ of $w$ from an independent half of the sample. Estimating
nuisance functions on a separate subsample is standard in this
literature; see the training samples of~\cite{hoif2008} and the
cross-fitting of~\cite{neweyrobins}. Let
$\mathcal W=\{f\in\cH^\alpha(H_0):\epsilon\le f\le1-\epsilon\}$.

\begin{lemma}[A fixed-accuracy pilot]\label{lem:smooth-propensity-pilot}
Fix $d,\alpha,H_0$, $0<\epsilon<1/2$, $p_->0$, and $\zeta_0\in(0,1)$.
There are $H_1<\infty$, $c>0$, and, for each $n_1\ge1$, a map
$\widehat w$ from $n_1$ observations of $(X,R)$ to $C^\infty$ functions
on $[0,1]^d$, such that $(x,\text{data})\mapsto\widehat w(x)$ is
measurable, $\epsilon\le\widehat w\le1-\epsilon$,
$\|\widehat w\|_{C^\alpha}\le H_1$, and
\[
 \Pp_P\{\|\widehat w-w\|_\infty>\zeta_0\}\le c^{-1}e^{-cn_1}
\]
for every law of $(X,R)$ such that $R\in\{0,1\}$, the covariate density
is at least $p_-$, and $w=\E[R\mid X]\in\mathcal W$.
\end{lemma}

\begin{proof}
Let $\gamma=\min(\alpha,1)$ and $C_0=\sqrt d\,H_0$. Every
$f\in\mathcal W$ satisfies $|f(x)-f(y)|\le C_0\|x-y\|_2^\gamma$, by its
H\"older bound if $\alpha\le1$ and its gradient bound if $\alpha>1$. We
fix an integer $k$ with $C_0(2\sqrt d/k)^\gamma\le\zeta_0/4$, we
partition $[0,1]^d$ into $k^d$ congruent cubes $Q$, and we let $Q^+$ be
the set of points of $[0,1]^d$ at $\ell^\infty$-distance at most
$1/(2k)$ from $Q$, so that $\operatorname{diam}Q^+\le2\sqrt d/k$. We
fix a $C^\infty$ partition of unity $(\varphi_Q)$ of $[0,1]^d$ with
$\varphi_Q\ge0$ and $\supp\varphi_Q\subseteq Q^+$; for instance,
products of one-dimensional ones. Let $N_Q$ be the number of
observations with $X_i\in Q$, let $\widetilde w_Q$ be the average of
their $R_i$, with $\widetilde w_Q=1/2$ if $N_Q=0$, let
$\widehat w_Q$ be $\widetilde w_Q$ clipped to $[\epsilon,1-\epsilon]$,
and let $\widehat w=\sum_Q\widehat w_Q\varphi_Q$. Then $\widehat w$ is
$C^\infty$, depends measurably on the data, takes values in
$[\epsilon,1-\epsilon]$ as a convex combination, and satisfies
$\|\widehat w\|_{C^\alpha}\le H_1=\sum_Q\|\varphi_Q\|_{C^\alpha}$,
which depends only on $k$, $d$, and $\alpha$.

Let $\pi_Q=\Pp(X\in Q)\ge p_-k^{-d}$ and $m_Q=\E[R\mid X\in Q]$. Since
$m_Q$ averages $w$ over $Q$, we have
$m_Q\in[\epsilon,1-\epsilon]$ and $|m_Q-w(x)|\le\zeta_0/4$ for
$x\in Q^+$. Since clipping does not increase $|\widetilde w_Q-m_Q|$,
and $\varphi_Q(x)=0$ unless $x\in Q^+$, we have
\[
 \|\widehat w-w\|_\infty\le\max_Q|\widetilde w_Q-m_Q|+\zeta_0/4 .
\]
Conditionally on $N_Q=s\ge1$, the sum of the $R_i$ with $X_i\in Q$ is
$\operatorname{Binomial}(s,m_Q)$, and Hoeffding's
inequality~\cite[Theorem~1]{hoeffding1963} implies that
$\Pp\{|\widetilde w_Q-m_Q|>\zeta_0/2\mid N_Q=s\}\le2e^{-s\zeta_0^2/2}$,
a bound that also holds for $s=0$. Since
$N_Q\sim\operatorname{Binomial}(n_1,\pi_Q)$, we obtain
\[
 \Pp\{|\widetilde w_Q-m_Q|>\zeta_0/2\}
 \le2\bigl(1-\pi_Q+\pi_Qe^{-\zeta_0^2/2}\bigr)^{n_1}
 \le2e^{-n_1\pi_Q(1-e^{-\zeta_0^2/2})}.
\]
A union bound over the $k^d$ cells gives
$\Pp\{\|\widehat w-w\|_\infty>\zeta_0\}
\le2k^de^{-n_1p_-k^{-d}(1-e^{-\zeta_0^2/2})}$, which is at most
$c^{-1}e^{-cn_1}$ for a suitable $c>0$.
\end{proof}

\begin{proof}[Proof of the upper bracket on $\mathcal M^{\rm sep}$]
We fix $\zeta\in(0,1)$, and we let $\theta=(\alpha+\beta)/d$ and
$\nu=q_\alpha+q_\beta$. We apply
\cref{lem:smooth-propensity-pilot} with $\zeta_0=\zeta\epsilon$ to
the first $n_1=\lfloor n/2\rfloor$ observations. Conditionally on
these, the remaining $N=n-n_1$ observations are independent with law
$P$, and we use them with the observables
\[
 D'=\frac R{\widehat w(X)},\qquad U'=1,\qquad
 V'=\frac{RY}{\widehat w(X)},\qquad W'=0,\qquad\lambda'=1,
\]
which depend on $X$, as the assertions stated before the proof of
\cref{thm:generic}(a) in \cref{sec:tuning} allow. These
are bounded by $M_0'=\epsilon^{-1}$, and
\[
 w'=\frac w{\widehat w}\ge\frac\epsilon{1-\epsilon},\qquad
 a'=\frac{\widehat w}w,\qquad b'=b,\qquad g'=\frac{pw}{\widehat w},\qquad
 \int a'b'g'=\int bp .
\]
Since $\widehat w\in\cH^\alpha(H_1)$ and $w\in\mathcal W$, the function
$a'=\widehat w\cdot(1/w)$ lies in a ball $\cH^\alpha(H_2)$, with $H_2$
depending only on $d,\alpha,H_0,H_1,\epsilon$. Indeed, by the Leibniz
and chain rules, each derivative of $a'$ of order at most $\lceil\alpha\rceil-1$
is a polynomial in the derivatives of $\widehat w$ and $w$ of these
orders and in $1/w\le\epsilon^{-1}$, and each factor is bounded and
H\"older continuous with exponent $\alpha-\lceil\alpha\rceil+1$. On the
event $\mathcal A=\{\|\widehat w-w\|_\infty\le\zeta\epsilon\}$, we have
$|w/\widehat w-1|\le\zeta$, hence $g'\in I_\zeta$.

Conditionally on a pilot outcome in $\mathcal A$, the remaining observations
have law in $\cP(\alpha,\beta)$ for
$(D',U',V',W',\lambda')$ with the constants $M_0'$,
$\delta'=\epsilon/(1-\epsilon)$, $H'=\max(H_0,H_2)$, and
$[g_-,g_+]=I_\zeta$, and its target is the MAR mean. We let $\widehat\psi$ be
the estimator of \cref{thm:generic}(a) from \cref{sec:tuning},
applied with these constants and $\widehat w$ to the last $N$ observations
and clipped to $[0,1]$. Its tuning depends only on these constants, so that
$\widehat\psi$ is measurable, and on $\mathcal A$ its conditional mean
squared error is at most
$C\,\overline a_N(\theta,\tau_{I_\zeta},\nu)^2$, with $C$ independent
of the pilot. Since its squared error is at most one off $\mathcal A$,
its mean squared error is at most
$C\,\overline a_N(\theta,\tau_{I_\zeta},\nu)^2+c^{-1}e^{-cn_1}$. Since
$n/2\le N\le n$, the factors $N^{-\theta}$, the powers of $\log N$, and
$e^{-\kappa\sqrt{\log N}}$ differ from their values at $n$ by bounded
factors, the last because
$\sqrt{\log n}-\sqrt{\log N}\le\log2/\sqrt{\log N}$. We conclude that
$\overline a_N\le C'\overline a_n$. Since $e^{-cn_1}=o(n^{-1})$ and
$n^{-1/2}\le\overline a_n$ for large $n$, we obtain
$\risk(\psi,\mathcal M^{\rm sep})\le
C_\zeta\overline a_n(\theta,\tau_{I_\zeta},\nu)$.
\end{proof}

\section{Proof of \texorpdfstring{\cref{cor:applications}}{the corollary} for these targets}\label{sec:broad-lower}

\emph{Upper bounds.} We take the choices of
\cref{sec:mar-implications}, with the response space
$\{(0,0)\}\cup(\{1\}\times[0,1])$, $M_0=1$, $\delta=\eta$, and
$[g_-,g_+]=[\eta,\eta^{-1}]$. Then
$\mathcal M_\eta(\alpha,\beta,H)\subseteq\cP(\alpha,\beta)$, and
\cref{thm:generic}(a) implies the upper bracket for the MAR mean
on $\mathcal M_\eta$. The upper bracket on $\mathcal M^{\rm sep}$ with
$I_\zeta$, for every $H_0>1$, was proved in
\cref{sec:broad-inclusions}.
The map that keeps arm $j$ sends
$\mathcal T_\eta(\alpha,\beta_0,\beta_1,H)$ into
$\mathcal M_\eta(\alpha,\beta_j,H)$, since its output has observation
probability $w_j$ with $\eta\le w_j\le1-\eta$,
$1/w_j=a_j\in\cH^\alpha(H)$, and $\eta\le w_jp=g_j\le\eta^{-1}$, and
regression $b_j\in\cH^{\beta_j}(H)$. It sends $\mathcal T^{\rm sep}$
into the class $\mathcal M^{\rm sep}$ with indices $(\alpha,\beta_j)$
and radius $H_0+1$, since $\epsilon\le1-w\le1-\epsilon$ and
$\|1-w\|_{C^\alpha}\le1+\|w\|_{C^\alpha}$. Since the MAR mean of the
output is $\mu_j$, \cref{lem:reductions}(c) shows that $\mu_j$
satisfies the upper bracket $((\alpha+\beta_j)/d,I,q_\alpha+q_{\beta_j})$
on $\mathcal T_\eta$, with $I=[\eta,\eta^{-1}]$, and on
$\mathcal T^{\rm sep}$ with $I_\zeta$. \Cref{lem:reductions}(a),
applied to $T_{\rm ATE}=\mu_1-\mu_0$ and to the relations on the left
of \eqref{eq:att-atu-identities}, together with
\cref{rem:combining-rates}, gives the upper brackets of the
three effects with the indices \eqref{eq:treatment-smoothness} and
$\nu=q_\alpha+q_{\beta_*}$.

\emph{Lower bounds.} We treat two cases at once. In the first,
$\mathcal M=\mathcal M_\eta$, $\mathcal T=\mathcal T_\eta$,
$I=[\eta,\eta^{-1}]$, and $[r_-,r_+]=[2\eta,2\eta^{-1}]$. In the second,
$\mathcal M=\mathcal M^{\rm sep}$, $\mathcal T=\mathcal T^{\rm sep}$, and
$I=[r_-,r_+]=[p_-,p_+]$. Then $r_-<1<r_+$, since $\eta<1/4$ or by
\eqref{eq:design}, and since $\rho$ depends only on $r_+/r_-$, the
value of $\tau$ computed from $[r_-,r_+]$ is $\tau_I$. For the MAR mean,
let $\beta'=\beta$ and $j=1$. For the effects, let $\beta'=\beta_j$,
where $\beta_j=\min(\beta_0,\beta_1)$ for the ATE, $j=0$ for the ATT,
and $j=1$ for the ATU, so that $s=(\alpha+\beta')/2$ by
\eqref{eq:treatment-smoothness}. If $P$ is a law of missing-data
observations with $Y\in\{0,1\}$, then $K_jP$ has propensity $w$ or
$1-w$, regression $b$ in arm $j$ and the constant $1/2$ in the other
arm, outcomes in $\{0,1\}$, and covariate density $p$. Since the
constant $1/2$ lies in every H\"older ball of radius $H_0$ or $H$, the
law $P$ lies in $\mathcal M$, with $\beta'$ in place of $\beta$, and $K_jP$ lies
in $\mathcal T$, if
$1/w,1/(1-w)\in\cH^\alpha(H)$, $b\in\cH^{\beta'}(H)$,
$\eta\le w\le1-\eta$, and $\eta\le wp,(1-w)p\le\eta^{-1}$ in the first
case, and if $w,1-w\in\cH^\alpha(H_0)$, $b\in\cH^{\beta'}(H_0)$,
$\epsilon\le w\le1-\epsilon$, and $p_-\le p\le p_+$ in the second. By
the paragraph that defines the MAR family, these are conditions of
the types (H), (I), and (W) of \cref{lem:local}, for functions $K$
with $k_0=2$ or $k_0=1/2$, or with $K\equiv1$ for the density. Their
hypotheses hold, since $2<H$, $1/2<H_0$, $\eta<1/2<1-\eta$,
$\epsilon<1/2<1-\epsilon$, and $r_\pm/2=\eta^{\mp1}$ in the first case
and $r_\pm=p_\pm$ in the second.

Suppose first that $s<d/4$. By \cref{lem:local}, applied with a
small enough fixed value of the input $\epsilon$ of
\cref{prop:local-lower} (not the overlap bound of
$\mathcal T^{\rm sep}$), the MAR family with indices $(\alpha,\beta')$
and the interval $[r_-,r_+]$ lies in $\mathcal M$, and its image
$\cP'_n=K_j\cP^{\rm h}_n$ lies in $\mathcal T$, for all large $n$. Let
$\theta=2s/d$ and let $t_n=c_1\underline a_n(\theta,\tau_I)$ be the
scale at which the MAR mean is hard on the MAR family. For the MAR
mean, \cref{lem:hard-risk} gives the lower bracket on
$\mathcal M$. For the effects, \cref{lem:reductions}(c) shows that
$\mu_j$ is hard at scale $t_n$ on $(\cP'_n)$. By \eqref{eq:kernel-ate},
$T_{\rm ATE}$ differs from $\mu_j$ on $\cP'_n$ by a sign and an
additive constant, so that it is also hard at scale $t_n$. For the ATT
and the ATU, we apply \cref{lem:reductions}(b), with $h(O)=(Y,A)$,
to the relations for $\mu_0$ and $\mu_1$ on the right of
\eqref{eq:att-atu-identities}. Here $nt_n^2\to\infty$, because
\[
 \log(nt_n^2)=(1-2\theta)\log n-2\kappa\sqrt{\log n}
 +(\theta+2)\log\log n+2\log c_1,
\]
so that the target is hard at scale $t_n/(4L)$ on $(\cP'_n)$. In each
case, the lower bracket on $\mathcal T$ follows from
\cref{lem:hard-risk}.

If $s\ge d/4$, we use the laws $Q_t$ under which $X$ is uniform on
$[0,1]^d$ and independent of $Z\sim\pi_{u,v}$ with $(u,v)=(0,t/2)$,
for $|t|<1/8$. They have $p=1$ and the constant functions $w=1/2$ and
$b=1/2+t$, so that they satisfy the conditions above, $Q_t\in\mathcal M$,
$K_jQ_t\in\mathcal T$, and the MAR mean is $\psi(Q_t)=1/2+t$. The
density of $Q_t$ with respect to $Q_0$ is $1+t\mathsf s_v(Z)/2$, which
is affine in $t$ and at least $3/4$, and
\cref{lem:parametric-path} implies that
$\risk(\psi,\mathcal M)\ge cn^{-1/2}$. By \eqref{eq:kernel-constant},
each effect $T$ of $K_jQ_t$ equals $(2j-1)t$, and
\cref{lem:reductions}(c), applied to the target $(2j-1)T+1/2$ on
$\mathcal T$, gives $\risk(T,\mathcal T)\ge cn^{-1/2}$.

For the MAR mean on $\mathcal M_\eta$, the lower bracket also holds on
every class between a localized class $\cP_r(\pi_0)$ and
$\mathcal M_\eta$, by \cref{thm:generic}(b). The baseline
$\pi_0$ satisfies \eqref{eq:generic-nondegeneracy}: $\Sigma$ is
positive definite, and $0<\eta<1/4$ and $H>2$ give
$\max\{\delta,g_-\}=\eta<w_0=1/2<\eta^{-1}=g_+$, $|a_0|=2<H$, and
$|b_0|=1/2<H$. We have $\cP_r(\pi_0)\subseteq\mathcal M_\eta$ whenever
$0<r<1/2-\eta$, since $w\le1/2+r<1-\eta$ on $\cP_r(\pi_0)$.

\emph{The logarithmic asymptotics.} Let $\cP$ be $\mathcal M^{\rm sep}$
or $\mathcal T^{\rm sep}$, let $s<d/4$ and $\theta=2s/d$, and let
$\kappa_\zeta$ be the $\kappa$ of \eqref{eq:subcritical-scale} computed
from $\tau_{I_\zeta}$. Since
$\log a_n(\theta,\tau)=-\theta\log n-\kappa\sqrt{\log n}
+(\theta/2)\log\log n$, the lower bracket with $[p_-,p_+]$ and the
upper bracket with $I_\zeta$ imply that
\[
 -\kappa_p(\theta)\le\liminf_{n\to\infty}\frac{\log\risk+\theta\log n}{\sqrt{\log n}}
 \le\limsup_{n\to\infty}\frac{\log\risk+\theta\log n}{\sqrt{\log n}}
 \le-\kappa_\zeta .
\]
As $\zeta\to0$, we have $\rho_{I_\zeta}\to\rho_{[p_-,p_+]}$, hence
$\tau_{I_\zeta}\to\tau_p$ and $\kappa_\zeta\to\kappa_p(\theta)$, which
gives \eqref{eq:sep-log-risk}.
% ---------- C-covariances.tex
\chapter{Conditional covariances and variances}\label{sec:products-proof}

\emph{Bilinear and quadratic targets.} We apply
\cref{thm:generic} with $D=1$, $M_0=1$, $H=H_0$,
$[g_-,g_+]=[p_-,p_+]$, and $\delta=1/2$, with $(W,\lambda)=(0,1)$ for
$\E_P[m_Um_V]$ and $Q_2$, and $(W,\lambda)=(UV,-1)$ for
$\E_P[\Cov_P(U,V\mid X)]$ and $\bar\sigma^2$. Then $w=1$, $a=m_U$,
$b=m_V$, and $g=p$, so that the generic class is
$\cP_{\rm bil}(\alpha,\beta)$, or $\cP_{\rm quad}(s)$ when $U=V=Y$ and
$\alpha=\beta=s$. If $U,V$ are independent Rademacher variables, then
$(w_0,a_0,b_0)=(1,0,0)$, the residual covariance in
\eqref{eq:generic-nondegeneracy} is the identity, and the strict
inequalities there follow from \eqref{eq:design} and $H_0>1$. A Rademacher variable $Y$ satisfies the
diagonal alternative with residual variance one. This gives the
brackets of these four targets.

\emph{Upper brackets.} We write $m_j=\E_P[Y^j]$ for $j=1,2$. For a fixed
known measurable $f_0$ with $B_0=\|f_0\|_\infty<\infty$, we set
$\ell_0=\E_P[f_0(X)^2-2Yf_0(X)]$ and $E_{f_0}=\E_P[(f(X)-f_0(X))^2]$.
Since $\E_P[Yf_0(X)]=\E_P[f(X)f_0(X)]$, we have
\[
 E_*=Q_2-m_1^2,\qquad
 R_*^2=\frac{Q_2-m_1^2}{\max\{m_2-m_1^2,v_{\min}\}},\qquad
 E_{f_0}=Q_2+\ell_0,
\]
where the middle identity holds on $\cP_{\rm quad}(s;v_{\min})$. In
\cref{lem:reductions}, we use $h(O)=(Y,Y^2)$, with
$\mathcal R=[-1,1]\times[0,1]$, for $E_*$ and $R_*^2$, and
$h(O)=f_0(X)^2-2Yf_0(X)$, with $\mathcal R=[-B_0^2-2B_0,B_0^2+2B_0]$,
for $E_{f_0}$. Since the denominator is at least $v_{\min}$, the three
maps are Lipschitz on $[0,1]\times\mathcal R$. By the bracket of $Q_2$
and \cref{lem:reductions}(a), the three targets satisfy the upper
bracket $(2s/d,[p_-,p_+],2q_s)$ on their classes.

\emph{Lower brackets below the threshold.} Let $s<d/4$. We use the hard
laws of \cref{sec:app-hard-laws} in the diagonal case, with
Rademacher baseline, $(D,U,V,W,\lambda)=(1,Y,Y,0,1)$,
$\mathsf s_{\rm d}=Y$, indices $(s,s)$, and $[r_-,r_+]=[p_-,p_+]$. Then
$P(Y=y\mid X)=\{1+y(u+v)\}/2$, $f=u+v$, and the target is $Q_2$. The
conditions of $\cP_{\rm quad}(s)$ are of type (H) for $K=u+v$, with
$k_0=0$ and $\gamma_K=s$, and of type (W) for $K\equiv1$. By
\cref{lem:local}, for a suitable fixed $\epsilon$, the hard laws
$\cP^{\rm h}_n$ lie in $\cP_{\rm quad}(s)$ for all large $n$. Since
$Y^2=1$ and $|\E_PY|\le\|u+v\|_\infty\le Cn^{-s/d}$, they also lie in
$\cP_{\rm quad}(s;v_{\min})$ for all large $n$. On them $Q_2$ is hard at
scale $t_n=c_1\underline a_n(2s/d,\tau_p)$, and $nt_n^2\to\infty$. The
exact inverse relations
$Q_2=E_*+m_1^2=(m_2-m_1^2)R_*^2+m_1^2=E_{f_0}-\ell_0$
are $L$-Lipschitz on $I_1\times\mathcal R$, with $L=2,\sqrt5,1$, where
$I_1=[0,1]$ for $E_*$ and $R_*^2$ and $I_1=[0,(1+B_0)^2]$ for
$E_{f_0}$. By \cref{lem:reductions}(b), with $T_1=Q_2$ and
$e_n=0$, the three targets are hard on $(\cP^{\rm h}_n)$ at scale
$t_n/(4L)$, and their lower brackets follow from
\cref{lem:hard-risk}.

\emph{Root-$n$ lower bounds.} We fix $\varphi=\eta\sin(2\pi x_1)$, where
$\eta>0$ depends only on $s,H_0$ and satisfies
$\|\varphi\|_{C^s}\le H_0$ and $\|\varphi\|_\infty\le1/2$. Let $P_t$
have $p=1$ and $P_t(Y=y\mid X)=\{1+yt\varphi(X)\}/2$ for $|t|<1$.
These laws lie in $\cP_{\rm quad}(s;v_{\min})$, with $\E_{P_t}Y=0$ and
$\Var_{P_t}(Y)=1$. We set $A_\varphi=\int\varphi^2>0$ and
$B_\varphi=\int\varphi f_0$. Then $E_*(P_t)=R_*^2(P_t)=t^2A_\varphi$
and $E_{f_0}(P_t)=t^2A_\varphi-2tB_\varphi+\int f_0^2$. Since the
excess-loss derivatives at $t=1/4$ and $t=1/2$ differ by $A_\varphi/2$,
at one of these points their absolute value is at least
$A_\varphi/4$; the derivatives of $E_*$ and $R_*^2$ there are at least
$A_\varphi/2$. The densities relative to the law with uniform design
and Rademacher response are $1+tY\varphi(X)\ge1/2$, with bounded slope.
\Cref{lem:parametric-path} yields the root-$n$ lower bounds for
all three targets. Only the constants for $E_{f_0}$ may also depend on
$B_0$.

\emph{The trial.} The map
$(X,A,Y_{\rm obs})\mapsto(X,(2A-1)(2Y_{\rm obs}-1))$ sends the trial
class into $\cP_{\rm quad}(s)$, since its response lies in $[-1,1]$ and
its conditional mean is $f$. Conversely, we draw an independent
$A\sim\operatorname{Bernoulli}(1/2)$ and map $(X,Y)$ to
$(X,A,\{1+(2A-1)Y\}/2)$. Since the outcome lies in $[0,1]$ and its arm
means are $(1\pm f)/2$, this kernel sends $\cP_{\rm quad}(s)$ into the
trial class. Both kernels preserve $p$, $f$, and the two targets.
\Cref{lem:reductions}(c) transfers the upper bounds along the first
kernel and the lower and probability bounds along the second.
% ---------- D-further.tex
\chapter{Wald ratios and the overlap-weighted effect}\label{sec:additional-effects}

We recall the causal interpretations of the targets of
\cref{sec:further-applications} and prove \cref{cor:applications}.
The interpretations use identification assumptions; the classes restrict observed laws.

\section{Wald ratios}\label{sec:wald}

\begin{example}[Wald ratio]\label{ex:wald}
Let $A(z)$ and $Y(a)$ denote potential treatments and outcomes, and
suppose that $A=A(Z)$ and $Y=Y(A)$. If
$Z\perp(A(0),A(1),Y(0),Y(1))\mid X$ and $A(1)\ge A(0)$, then
$T_{\rm W}$ is the complier average treatment effect
$\E[Y(1)-Y(0)\mid A(1)>A(0)]$, the local average treatment effect
of~\cite{imbensangrist,angristimbensrubin}, in the form with covariates
given in~\cite{frolich,dml}.
Its numerator and denominator are average effects of the instrument,
that is, differences of two MAR means.
\end{example}

\begin{example}[Average conditional Wald ratio]\label{ex:conditional-wald}
Suppose that $P(Z=1\mid X)=1/2$ and $A\le Z$, so that the instrument is
assigned with equal probability and treatment is unavailable when
$Z=0$. Under the assumptions of \cref{ex:wald}, $b(X)$ is the
conditional complier effect. If also
$\E[Y(1)-Y(0)\mid X,A(1)]=\E[Y(1)-Y(0)\mid X]$, which says that the
covariates account for differences in mean effects between compliers
and never-takers, then $T_{\rm CW}$ is the ATE;
see~\cite[Theorem~1 and the following discussion]{wangtchetgen}. The two
Wald targets differ in general, since $T_{\rm W}$ averages $b(X)$ with
weights proportional to $q(X)$.
\end{example}

\begin{proof}[Proof of \cref{cor:applications} for the Wald ratios]
\emph{Wald ratio.} We write $\mathcal N=\E_P[\Delta_Y(X)]\in[-1,1]$
and $\mathcal D=\E_P[\Delta_A(X)]\in[c_W,1]$. With the instrument as the
treatment, the deterministic maps $(X,Z,A,Y)\mapsto(X,Z,Y)$ and
$(X,Z,A,Y)\mapsto(X,Z,A)$ send $\mathcal I_{\eta,c_W}(\alpha,\beta,H)$
into $\mathcal T_\eta(\alpha,\beta,\beta,H)$, and the ATEs of the images
are $\mathcal N$ and $\mathcal D$. By \cref{cor:applications}
for the ATE on $\mathcal T_\eta(\alpha,\beta,\beta,H)$ and
\cref{lem:reductions}(c), both satisfy the upper bracket
$((\alpha+\beta)/d,[\eta,\eta^{-1}],q_\alpha+q_\beta)$ on
$\mathcal I_{\eta,c_W}(\alpha,\beta,H)$, and applying \cref{lem:reductions}(a)
to the ratio, with $C_0=1$ and $\delta_{\mathcal D}=c_W$, together with
\cref{rem:combining-rates}, we obtain the upper bracket for
$T_{\rm W}$. For the lower bracket, the kernel
$(X,Z,Y)\mapsto(X,Z,Z,Y)$, which sets $A=Z$, maps
$\mathcal T_\eta(\alpha,\beta,\beta,H)$ into
$\mathcal I_{\eta,c_W}(\alpha,\beta,H)$: its output has $m_{A,z}=z$, whose
H\"older norm is $z\le1<H$, and $\Delta_A=1\ge c_W$, while the remaining
conditions are those of $\mathcal T_\eta(\alpha,\beta,\beta,H)$ with $Z$
as the treatment. Since $T_{\rm W}$ equals the ATE of the input on the
image, \cref{lem:reductions}(c) transfers its lower bracket.

\emph{Average conditional Wald ratio.} We take the response
$(Z,A,Y)$, the functions $(D,U,V,W,\lambda)=(2A,1,2(2Z-1)Y,0,1)$, and
$M_0=2$. Since
$P(Z=1\mid X)=1/2$ and $A\le Z$, we have
$\E_P[2A\mid X]=q$ and $\E_P[2(2Z-1)Y\mid X]=\Delta_Y$. Then $w=q$,
$a=1/q$, $b=\Delta_Y/q$, $g=qp$, and $\psi=\int(\Delta_Y/q)p=T_{\rm CW}$,
so that $\mathcal J_\eta(\alpha,\beta,H)\subseteq\cP(\alpha,\beta)$ with
$\delta=\eta$ and $[g_-,g_+]=[\eta,\eta^{-1}]$, and the upper bracket
follows from \cref{thm:generic}(a). For the lower
bracket, the kernel that draws $Z\sim\operatorname{Bernoulli}(1/2)$
independently and maps $(X,R,RB)$ to $(X,Z,ZR,ZRB)$ maps
$\mathcal M_\eta(\alpha,\beta,H)$ into $\mathcal J_\eta(\alpha,\beta,H)$:
the output has $A=ZR\le Z$, $q=P(R=1\mid X)$, and $\Delta_Y=qb$ with
$b=\E[B\mid R=1,X]$, so that its conditional Wald ratio is $b$ and its
target is the MAR mean of the input. \Cref{lem:reductions}(c)
therefore transfers its lower bracket.
\end{proof}

\section{The overlap-weighted effect}\label{sec:overlap-effect}

\begin{example}[Overlap-weighted effect]\label{ex:overlap-effect}
For fixed $\gamma$, minimizing $\E_P[(Y-\gamma A-h(X))^2]$ over $h$
gives $h=m-\gamma w$, and minimizing the resulting quadratic in
$\gamma$ gives \eqref{eq:overlap-effect}. This definition places no
restriction on $\E_P[Y\mid A,X]$, and the numerator and the denominator
of \eqref{eq:overlap-effect} are expected conditional covariances as in
\cref{sec:product-class}. Under consistency and conditional
exchangeability, \eqref{eq:overlap-effect} is the overlap-weighted
treatment effect~\cite{limorganzaslavsky}; see
also~\cite[Remark~9.3.3]{causalmlbook} for this interpretation of
regression on residuals when effects vary with $X$.
\end{example}

The weights $w(1-w)$ appear earlier in Crump et al.~\cite[Section~5.4,
Corollary~5.2]{crump2006}, as the weights that minimize the efficiency
bound for weighted sample average treatment effects under
homoscedasticity.

Robins et al.~\cite[Example~1c]{hoif2008} impose smoothness on
$m-\gamma_{\rm OW}w$ instead of on $m$. Since $|\gamma_{\rm OW}|\le1$
and $\cH^\alpha(H)\subseteq\cH^\beta(CH)$ when $\beta\le\alpha$, the two
conditions agree up to fixed changes of radius when $\beta\le\alpha$,
but they can differ when $\beta>\alpha$. Without smoothness assumptions
on the covariate density, and without the assumption that the
conditional treatment effect is constant, Robins et
al.~\cite[Remark~4.2]{hoif2008} refer to $n^{-2s/d}$, for $s<d/4$, as
their conjectured minimax rate for this effect. Here $s$ is the average
of the smoothness indices of the propensity and of their regression.
For $\beta\le\alpha$, \cref{cor:applications} confirms the exponent
$2s/d$, with $s=(\alpha+\beta)/2$, and shows that the minimax rate is
smaller than $n^{-2s/d}$ by a factor
$e^{-(\kappa_p(2s/d)+o(1))\sqrt{\log n}}$. Concurrently,
Park~\cite[Section~2]{park} reports minimax bounds under a known
positive lower bound on the denominator, with effective smoothness
$\min(\alpha+\beta,2\beta)$ in his notation, for the normalized
covariance $\E[\Cov(A,Y\mid X)]/\E[\Var(Y\mid X)]$ of two bounded
responses, which has the same form as $\gamma_{\rm OW}$ with the roles
of $A$ and $Y$ exchanged; the proof is in the supplement
of~\cite{park} (Theorem~S14), which was not available to us.

\begin{proof}[Proof of \cref{cor:applications} for the overlap-weighted effect]
Here $w=P(A=1\mid X)$ is the propensity, while the generic $D$ is one.
We write $\gamma_{\rm OW}=\mathcal N/\mathcal D$, where
$\mathcal N=\E_P[\Cov_P(A,Y\mid X)]\in[-1/4,1/4]$ and
$\mathcal D=\E_P[w(1-w)]\in[\epsilon(1-\epsilon),1/4]$, and we let
$\theta=\min\{(\alpha+\beta)/d,2\alpha/d\}$.

\emph{Upper bracket.} We apply \cref{thm:generic}(a) with
$D=1$, $M_0=1$, $H=H_0$, $\delta=1/2$, and $[g_-,g_+]=[p_-,p_+]$.
The choices $(U,V,W,\lambda)=(A,Y,AY,-1)$ and $(A,A,A,-1)$ give
$\mathcal N$ and $\mathcal D$, with upper brackets
$((\alpha+\beta)/d,[p_-,p_+],q_\alpha+q_\beta)$ and
$(2\alpha/d,[p_-,p_+],2q_\alpha)$, respectively. Applying
\cref{lem:reductions}(a) to the ratio with $C_0=1/4$ and
$\delta_{\mathcal D}=\epsilon(1-\epsilon)$, together with
\cref{rem:combining-rates}, we obtain the upper bracket with
exponent
$\theta$ and $\nu=q_\alpha+q_{\min(\alpha,\beta)}$.

\emph{A common response family.} Let $S_1=2A-1$ and $S_2=2Y-1$.
For $c\in\{0,1/4\}$ we define the baseline and scores by
\[
 \pi_0^{(c)}=\frac{1+cS_1S_2}{4},\qquad
 \mathsf s_i=\frac{2S_i}{1+cS_1S_2}\quad(i=1,2).
\]
Since the baseline factor cancels in score expectations, we have
$\E_{\pi_0^{(c)}}\mathsf s_i=0$ and
$\E_{\pi_0^{(c)}}[(S_j/2)\mathsf s_i]=\mathbf1_{i=j}$. The baseline means
of $A$ and $Y$ are $1/2$, and therefore $\mathsf s_1,\mathsf s_2$ are
scores for
$U=A$, $V=Y$, and $\mathsf s_1$ is a diagonal score for $U=V=A$.
For small functions $\ell_1,\ell_2$ of $x$, the law
$\pi_0^{(c)}(1+\ell_1\mathsf s_1+\ell_2\mathsf s_2)$ and its moments are
\[
 P(S_1=\sigma_1,S_2=\sigma_2\mid X)
 =\frac{1+c\sigma_1\sigma_2+2\ell_1\sigma_1+2\ell_2\sigma_2}{4},
 \qquad w=\frac12+\ell_1,\quad m=\frac12+\ell_2,
\]
\[
 \mathcal N=\frac c4-\int p\ell_1\ell_2,\qquad
 \mathcal D=\frac14-\int p\ell_1^2.
\]

\emph{Lower bracket when $\theta<1/2$.} If $\beta\le\alpha$, we use the
hard laws with indices $(\alpha,\beta)$, $c=0$, $(\ell_1,\ell_2)=(u,v)$,
and $(U,V,W,\lambda)=(A,Y,AY,-1)$. If $\beta>\alpha$, we use indices
$(\alpha,\alpha)$, $c=1/4$, $(\ell_1,\ell_2)=(u+v,0)$, and
$(U,V,W,\lambda)=(A,A,A,-1)$, with the diagonal score $\mathsf s_1$.
In both cases $D=1$, $[r_-,r_+]=[p_-,p_+]$, and the indices give the
exponent $\theta$. For membership in $\cP_{\rm OW}(\alpha,\beta)$, we
require $p_-\le p\le p_+$, $w\in\cH^\alpha(H_0)$,
$m\in\cH^\beta(H_0)$, and $\epsilon\le w\le1-\epsilon$. The density and
overlap conditions are of types (W) and (I). In the family with indices
$(\alpha,\beta)$, both smoothness conditions are of type (H), with
$k_0=1/2<H_0$, $\gamma_K=\alpha$ for $w$, and $\gamma_K=\beta$ for $m$.
In the family with indices $(\alpha,\alpha)$, the condition on $w$ is
of type (H), with $\gamma_K=\alpha$, while $m=1/2$ satisfies
$\|m\|_{C^\beta}=1/2<H_0$ directly. By \cref{lem:local}, the hard
laws lie in $\cP_{\rm OW}(\alpha,\beta)$ for all large $n$, and their
generic target, $\mathcal N$ in the first case and $\mathcal D$ in the
second, is hard on them at scale $t_n=c_1\underline a_n(\theta,\tau_p)$.

In the first case, eventually $\mathcal D\ge1/8$, and we have
\[
 \Bigl|\mathcal N-\frac{\gamma_{\rm OW}}4\Bigr|
 =|\gamma_{\rm OW}|\int pu^2
 \le8\|u\|_\infty^3\|v\|_\infty
 \le Cn^{-\theta-2\alpha/d}=o(t_n).
\]
The map $\gamma\mapsto\gamma/4$ is $1$-Lipschitz on $[-1,1]$, which
contains $\gamma_{\rm OW}$. In the second case, $\mathcal N=1/16$ and
eventually $\mathcal D\in[3/16,1/4]$, so that
$\gamma_{\rm OW}\in[1/4,1/3]$ and $\mathcal D=1/(16\gamma_{\rm OW})$;
this map is $1$-Lipschitz on $[1/4,1/3]$.
By \cref{lem:reductions}(b), with $k=0$, $T_1$ the generic target,
and $T_0=\gamma_{\rm OW}$, the target $\gamma_{\rm OW}$ is hard at
scale $t_n/4$ in both cases, and \cref{lem:hard-risk} gives the lower
bracket.

\emph{Root-$n$ lower bound.} Let $P_t$ have $p=1$ and response law
$\{1+tS_1S_2\}/4$ for $|t|<1/2$. Then $w=m=1/2$,
$\mathcal N=t/4$, $\mathcal D=1/4$, and $\gamma_{\rm OW}=t$.
These laws lie in $\cP_{\rm OW}(\alpha,\beta)$, and their densities
relative to the law with uniform design and uniform response are
$1+tS_1S_2\ge1/2$, with bounded slope.
\Cref{lem:parametric-path} implies the root-$n$ lower bound.
\end{proof}
